\documentclass[11pt]{amsart}
\usepackage[T1]{fontenc}
\usepackage{lmodern,amsmath,amssymb,amsthm,mathtools,microtype,booktabs,array,longtable}
\usepackage[margin=1.06in]{geometry}
\usepackage[hidelinks]{hyperref}
\emergencystretch=2em
\numberwithin{equation}{section}
\newtheorem{theorem}{Theorem}[section]
\newtheorem{lemma}[theorem]{Lemma}
\newtheorem{proposition}[theorem]{Proposition}
\newtheorem{corollary}[theorem]{Corollary}
\theoremstyle{definition}\newtheorem{definition}[theorem]{Definition}
\theoremstyle{remark}\newtheorem{remark}[theorem]{Remark}
\newcommand{\Q}{\mathbb Q}\newcommand{\R}{\mathbb R}\newcommand{\C}{\mathbb C}\newcommand{\N}{\mathbb N}\newcommand{\Z}{\mathbb Z}\newcommand{\E}{\mathbb E}
\newcommand{\ip}[2]{\langle #1,#2\rangle}\newcommand{\norm}[1]{\lVert #1\rVert}
\newcommand{\epsc}{\varepsilon_{\mathrm c}}
\DeclareMathOperator{\osc}{osc}\DeclareMathOperator{\spanop}{span}\DeclareMathOperator{\End}{End}\DeclareMathOperator{\diag}{diag}\DeclareMathOperator{\tr}{tr}\DeclareMathOperator{\rad}{rad}\DeclareMathOperator{\TV}{TV}
\title[Stochastic realization and orbit geometry]{Stochastic Realization, Finite Quotients,\\and Uniform Orbit Geometry}
\author{Qian Qi}
\date{27 September 2026}
\subjclass[2020]{Primary 93B15, 68Q45; Secondary 22C05, 60J10, 11J13}
\keywords{stochastic realization, finite quotient, Ramsey extraction, compact group, spectral gap, nonuniform width}
\begin{document}
\begin{abstract}
We study finite stochastic realizations of continuous compact-group interfaces, allowing clock-dependent transitions and a new machine at each horizon. Under eventual approximate returns, the minimum label width converges to the stationary minimum, and every smaller width has a uniform occupation bound. Stationary realizations admit a finite physical-group action without increasing width or error; initialization may remain randomized. This gives an equality-inclusive finite-quotient criterion even with infinitely many components. Quantitative width is governed by two geometric exponents: a small-ball bound uniform over all conditional-centroid directions, and the covering dimension of the target orbit. For full density-matrix outputs under spectral-gap alphabets in $SU(d)$, we prove width between $c(N/\log(N+2))^{d-1}$ and $C(N+1)^{d-1}$ throughout the fixed subcritical range with strict amplitude slack. The upper realization is exact and positive. At pure amplitude, every fixed positive subcritical error obeys the same orders, whereas a fixed algebraic free alphabet and a specified transcendental seed have exact width $2\cdot3^N-1$. We also retain phasewise classification, planar and profinite width laws, existential-real completeness for finite-group input, and provide positivity-preserving finite-precision budgets. Spectral gap and dense free generation are imported inputs, separately identified from the realization proofs.
\end{abstract}
\maketitle
\section{Introduction}
A finite stochastic register has a continuum of probability distributions. Its number of labels is therefore not determined by the cardinality of the numerical orbit it realizes. Allowing an external clock creates a further difficulty: the command rows may change at every position, and the entire realization may be redesigned at each horizon. The central question is whether these freedoms change the least width needed to preserve a fixed numerical experiment with a prescribed uniform error.

Our first theorem identifies the limiting clocked width with the stationary minimum, without a loss of labels or accuracy. The synchronization hypothesis is eventual approximation of the physical identity at every sufficiently large length. Under it, arbitrarily many cuts of width at most $k$ force a stationary realization on $k$ labels, even when the intervening registers are unrestricted. The proof selects approximately homogeneous composite transition tables by finite Ramsey theory. It applies to irreversible topological monoids as well as groups. It does not assume that a physical return acts trivially on the hidden register.

For compact groups a second reduction makes the stationary realization physically equivariant. Transition randomness can be removed, and hidden permutations above the physical identity can be averaged out, without increasing width or uniform error. The stationary minimum is therefore intrinsic to a finite continuous group action with randomized initialization. Combining the two reductions gives an equality-inclusive finite-quotient criterion for clocked boundedness, including positive critical radii on groups with infinitely many components. A moving-frame construction gives the exact limiting minimum in each return-length phase, again without additional labels.

\subsection{The interface and its two resources}
Unless a section explicitly specifies a different physical space, $H$ is a compact Hausdorff group and $\mathcal A$ is a finite nonempty alphabet whose letters $u_a\in H$ generate a dense subgroup. The execution convention is
\[
                  u_{a_1\cdots a_n}=u_{a_n}\cdots u_{a_1}.
\]
There are finite nonempty seed and query sets $\mathcal S,\mathcal J$ and continuous maps $F_{s,j}:H\to Y_j$. In the structural results $Y_j$ is a nonempty norm-compact convex subset of a real Banach space. The localization and algebraic sections specify finite-dimensional output spaces. For a categorical law, $Y_j=\Delta_{M_j}$ and the norm on differences is $\frac12\|\cdot\|_1$, namely total variation. A prescribed joint law is one query; incompatible measurements are not assigned an artificial joint law.

\begin{definition}[Stationary realization]\label{def:stationary55}
A realization on $k$ labels consists of initialization rows $\alpha_s\in\Delta_k$, row-stochastic matrices $T_a$ of order $k$, and decoder values $D_j(i)\in Y_j$, all independent of time and word length. Its error is
\[
       \sup_{s,j,w\in\mathcal A^*}
        \|\alpha_sT_{a_1}\cdots T_{a_n}D_j-F_{s,j}(u_w)\|_j.
\]
The empty word is included. Let $M_\varepsilon$ be the least $k$ at error at most $\varepsilon$, or $\infty$ when none exists. A permutation realization has permutation command matrices; its initialization need not be deterministic.
\end{definition}
The clocked resource $W_{N,\varepsilon}$ instead minimizes the maximum size of the registers $S_0,\ldots,S_N$ at a fixed horizon, with arbitrary row-stochastic transitions depending on time and horizon, and a terminal query decoder. Definition~\ref{def:machine54} states the model formally. Every advertised persistent label counts, including unused padding. Private information retained for future commands is part of the current label. The original width model does not charge the external clock, real table descriptions, their construction, exact arithmetic, or exact atomic sampling. Section~\ref{sec:complexity56} gives separate finite-description and fair-random-bit bounds; it does not reinterpret those resources as free computational memory. A numerical decoder vector is not an additional sampled-answer register.

For an open normal subgroup $U\le H$ put
\begin{equation}\label{eq:quotientradius55}
 r(U)=\max_{s,j,h\in H}\rad_{Y_j}(F_{s,j}(hU)),
 \qquad\rad_Y(A)=\min_{y\in Y}\sup_{x\in A}\|x-y\|.
\end{equation}
The coset images are compact by continuity. The radius depends continuously on the translate, and compactness gives the displayed extrema. The restriction of the center to the legal output set is essential. The component radius $R_0$ uses $U=H^\circ$, whether or not this subgroup is open; it is the infimum of the finite-quotient radii. At the critical error, the relevant issue is attainment of that infimum, not a limit from larger errors.

\subsection{The principal conclusions}
Theorem~\ref{thm:stationarization56} proves $W_{N,\varepsilon}\to M_\varepsilon$ under eventual approximate returns. Finite limits are eventually equal, and every $k<M_\varepsilon$ has a uniform narrow-cut occupation bound. Theorems~\ref{thm:equivariant56} and \ref{thm:completeboundary56} then give the intrinsic group-action minimum and its finite-quotient existence criterion. Theorem~\ref{thm:phaseminimum56} treats each length phase by a time-dependent change of physical frame. Table~\ref{tab:map56} separates these quantifiers and hypotheses.

The quantitative consequences have additional hypotheses. For a specified Diophantine planar alphabet the full-subcritical matching polynomial law is retained. A fixed $2$-adic interface has linear exact width and bounded width at every positive error. For a spherical interface in dimension $d$ with a spectral-gap alphabet, Theorem~\ref{thm:spherical56} gives the order $N^{(d-1)/2}$ up to a logarithmic factor in the lower bound. A rational alphabet generating $\mathrm{SO}(3)$ supplies a genuinely noncommuting example. Its spectral gap is imported from Bourgain--Gamburd, not inferred from a finite numerical test. Theorem~\ref{thm:ETR56} gives an $\exists\R$-complete finite input problem, while Theorem~\ref{thm:bits56} prices dyadic descriptions and fair random bits.


The quantitative theory also extends beyond spheres. Theorem~\ref{thm:uniformorbit57} separates a small-ball exponent uniform over all unit feature directions from the dimension of the target orbit. For conjugation on traceless Hermitian matrices a spectral split gives the uniform exponent $2(d-1)$, including directions with repeated eigenvalues. The pure-state orbit has the same dimension. Theorem~\ref{thm:matrixwidth57} therefore gives nearly matching width exponents $d-1$ for full legal density-matrix outputs in every dimension. Its exact upper construction stays inside the positive cone, rather than replacing the legal output set by an ambient Euclidean ball.

At an extreme target, the exact answer can be very different. Theorem~\ref{thm:extreme57} identifies the entire optimal exact width profile with the reachable orbit cardinalities. Theorem~\ref{thm:freeendpoint57} applies it to one fixed five-letter algebraic dense free alphabet in $SU(d)$: exact pure-state width is $2\cdot3^N-1$, while each fixed positive subcritical error has the polynomial bounds just described. The specified seed is transcendental, not a rational-input instance. These assertions concern a full matrix-valued terminal interface, not adaptive observations or the zero-error endpoint of a single scalar readout.

The compact nonnegative-matrix group structure used in purification is classical, and is separated from the external numerical error constraints throughout Section~\ref{sec:comparison54}. Neutral-letter collapse for advice automata is also a relevant predecessor to stationarization. The additional assertions proved here concern the same stochastic label bound and the same closed numerical error balls, arbitrary intervening widths, and approximate rather than exact returns. Neither classical theorem is relabeled as a new result.

\clearpage
\begin{table}[p]
\centering\small
\caption{Resources, hypotheses, and conclusions. All widths count persistent labels; all errors are worst-word numerical errors.}\label{tab:map56}
\begin{tabular}{p{.20\textwidth}p{.35\textwidth}p{.35\textwidth}}
\toprule
Object & Hypotheses & Conclusion\\
\midrule
$M_\varepsilon$ & Compact $H$; finite dense alphabet; compact convex outputs; no returns needed & Least finite continuous $H$-action with randomized initialization (Theorem~\ref{thm:equivariant56})\\[3pt]
$r(U)$ & $U$ open normal; arbitrary number of components & $M_\varepsilon<\infty$ iff some $r(U)\le\varepsilon$, equality included (Theorem~\ref{thm:finitequotient55})\\[3pt]
$W_{N,\varepsilon}$ & Eventual approximate returns; even an irreversible topological monoid & $W_{N,\varepsilon}\to M_\varepsilon$; finite limits eventually equal; all-profile occupation (Theorem~\ref{thm:stationarization56})\\[3pt]
$R_0$ & Compact group and eventual approximate returns & $R_0=\inf_Ur(U)$; critical boundedness iff the infimum is attained (Theorem~\ref{thm:completeboundary56})\\[3pt]
$M_{r,\varepsilon}$ & Block length $p$, phase subgroup $J$; block approximate returns & $W_{N,\varepsilon}\to M_{r,\varepsilon}$ for $N\equiv r$; no extra phase labels (Theorem~\ref{thm:phaseminimum56})\\[3pt]
$\varepsilon_r$ & Finitely many components; nonempty exact-return language & Component radius in the reachable phase; earlier localized occupation proof (Theorem~\ref{thm:phases55})\\[3pt]
Exact boundary & Cofinite exact returns; finite-dimensional outputs & Finite translate rank yields a finite observable quotient (Theorem~\ref{thm:zero55})\\[3pt]
Quantitative width & Separate arithmetic, profinite, or uniform-orbit spectral-gap hypotheses & Planar and dyadic laws; spherical and matrix-orbit exponents up to a logarithm (Theorems~\ref{thm:uniformorbit57}, \ref{thm:matrixwidth57})\\[3pt]
Extreme outputs & Affine group action; target orbit consists of legal extreme points & Exact profile equals orbit cardinalities; free example $2\cdot3^N-1$ (Theorems~\ref{thm:extreme57}, \ref{thm:freeendpoint57})\\[3pt]
Finite encodings & Explicit finite group and rational categorical tables & $\exists\R$-complete stationary decision; dyadic precision and random-bit budgets (Theorems~\ref{thm:ETR56}, \ref{thm:bits56})\\
\bottomrule
\end{tabular}
\end{table}
\clearpage

\section{Stationarization from recurrent narrow cuts}\label{sec:ramsey56}
The first reduction does not require reversible physical dynamics. In this section only, replace $H$ by a topological monoid $G$ with identity $1$, retain a finite command alphabet, and let $F_{s,j}:G\to Y_j$ be continuous. The output sets are norm-compact convex subsets of real Banach spaces. Define $M_\varepsilon$ and $W_{N,\varepsilon}$ by the same stochastic models. Compactness of the physical monoid is not required.

\begin{definition}[Eventual approximate returns]\label{def:returns56}
The alphabet has eventual approximate returns if, for every neighborhood $V$ of $1$, there is an integer $g(V)\ge0$ such that every integer $n\ge g(V)$ is the length of a word with product in $V$. Cofinite exact returns mean that a single $g$ works with product exactly $1$.
\end{definition}
The statement concerns physical products only. The hidden rows on a return word need not be close to an identity matrix. Approximation here is eventual at every length, not merely along a return subsequence.

\begin{theorem}[Width-preserving stationarization]\label{thm:stationarization56}
Assume eventual approximate returns. For every $\varepsilon\ge0$ and integer $k\ge1$, the following are equivalent:
\begin{enumerate}
\item a stationary realization with at most $k$ labels has error at most $\varepsilon$;
\item among error-$\varepsilon$ clocked realizations, the number of positive cuts having width at most $k$ is unbounded.
\end{enumerate}
Consequently, if $M_\varepsilon>k$, there is $L_{k,\varepsilon}<\infty$, independent of the horizon, such that every error-$\varepsilon$ realization satisfies
\begin{equation}\label{eq:occupation56}
 \#\{t\in\{1,\ldots,N\}:|S_t|\le k\}\le L_{k,\varepsilon}.
\end{equation}
All other registers may be arbitrarily large. Moreover,
\begin{equation}\label{eq:limitminimum56}
                  \lim_{N\to\infty}W_{N,\varepsilon}=M_\varepsilon
                  \quad\hbox{in }\N\cup\{\infty\}.
\end{equation}
If the right side is finite, equality holds for every sufficiently large $N$.
\end{theorem}
The theorem preserves the proposed width, including randomized initialization. It does not assert that the microscopic rows of the clocked machine become stationary. They can be entirely unrelated at different times.

\subsection{Finite witnesses and homogeneous transition tables}
For a matrix $A$ write $\|A\|_{\mathrm r}=\max_i\sum_j|A_{ij}|$. Put $a=|\mathcal A|$ and $D_* =\max(1,\max_j\operatorname{diam}Y_j)$. Let $\mathcal X_k$ be the compact parameter space of stationary $k$-label machines. For $m\ge0$ define
\[
 e_{k,m}=\min_{X\in\mathcal X_k}\max_{s,j,\,|w|\le m}
               \|X_{s,j}(w)-F_{s,j}(u_w)\|_j.
\]
The finite maximum is continuous, so the minimum is attained. If no stationary error-$\varepsilon$ machine exists, compactness of $\mathcal X_k$ and the finite intersection property give an $m$ with
\begin{equation}\label{eq:finitewitness56}
                       e_{k,m}=\varepsilon+\gamma,\qquad\gamma>0.
\end{equation}
Indeed the nested closed sets defined by the constraints $|w|\le m$ otherwise have a common point. This argument also applies when the full infinite-word error is not a continuous function of the parameters.

We shall use the finite edge-coloring theorem of Ramsey \cite{Ramsey}. Denote by $R_q(h)$ an integer such that every coloring of the edges of a complete graph on $R_q(h)$ vertices with at most $q$ colors has a monochromatic $h$-vertex subgraph. Only finiteness is essential. For definiteness one may take
\begin{equation}\label{eq:ramseybound56}
                   R_q(h)\le q^{qh+1}\qquad(q\ge2,\ h\ge2).
\end{equation}
To see this elementary bound, successively select a vertex and retain a largest color class among its remaining neighbors. With $q^{qh+1}$ initial vertices this produces at least $q(h-1)+1$ selected vertices, each of whose edges to later selected vertices has one color. At least $h$ of their associated colors agree; these vertices form the required subgraph. The deliberately excessive exponent avoids any endpoint rounding issue.

Partition $[0,1]$ into at most $\lceil k/\delta\rceil+1$ intervals of length at most $\delta/k$. Two tuples in the same entrywise cell of
\[
                  \mathsf{Stoch}(k)^{\{0\}\cup\mathcal A}
\]
have corresponding matrices within $\delta$ in $\|\cdot\|_{\mathrm r}$. The number of occupied cells is at most
\begin{equation}\label{eq:colors56}
             q=\max\bigl(2,(\lceil k/\delta\rceil+1)^{(a+1)k^2}\bigr).
\end{equation}
Each representative used below is chosen from an occupied cell and is therefore stochastic; an entrywise rounded matrix need not be stochastic and is not used as a representative.

\begin{proof}[Proof of Theorem~\ref{thm:stationarization56}]
The forward implication follows by using one stationary machine at every horizon. For the converse, suppose \eqref{eq:finitewitness56}. Choose $0<\eta<\gamma/4$. Continuity of multiplication and of the finitely many $F_{s,j}$ gives a neighborhood $V$ of $1$ with the following property. In any word of length at most $m$, inserting at most $m+3$ factors from $V$, including before and after the word, changes every output by less than $\eta$. There are only finitely many relevant command words and insertion patterns; continuity at the tuple of identity factors, followed by a finite intersection of neighborhoods, proves the assertion. This uses neither an invariant metric nor commutativity.

Let $g=g(V)$, put $b=g+1$, and choose
\[
       0<\delta\le \min\bigl(1,\gamma/(4D_*(m+1))\bigr).
\]
Take a clocked realization with so many narrow cuts that, after discarding cuts $t<g$ or $t>N-g$ and greedily retaining cuts separated by at least $b$, at least $R_q(m+2)$ remain. Such a choice is possible by hypothesis; at most $2g$ cuts are discarded and each retained cut accounts for at most $b$ of the others.

Identify each retained register with a subset of $\{1,\ldots,k\}$. Extend rows from dummy labels arbitrarily to stochastic rows without altering any row on the actual register. For retained cuts $t_i<t_j$ write $\ell=t_j-t_i\ge g+1$. Choose a length-$\ell$ return word with product in $V$, and let $A_{ij}$ be its actual composite kernel between these cuts. For each letter $c\in\mathcal A$, choose the word consisting of $c$ followed by a length-$(\ell-1)$ return word with product in $V$, and let $B_{ij,c}$ be its composite kernel. Thus the physical products of these words are respectively in $V$ and of the form $v u_c$ with $v\in V$. All the matrices are of order $k$, however large the intermediate registers may be.

Color the edge $\{i,j\}$, oriented from the earlier cut to the later cut, by the cell of $(A_{ij},(B_{ij,c})_c)$. Ramsey's theorem gives $m+2$ ordered cuts
\[
                         t_0<t_1<\cdots<t_{m+1}
\]
whose edges have one color. Fix a stochastic representative $(E,(B_c)_c)$ from that color. Choose fixed return words for the prefix of length $t_0$ and the suffix of length $N-t_{m+1}$, both with products in $V$. Let $\alpha_s$ be the actual state distribution after the prefix. Let $D_j(i)$ be the conditional expected original decoder after the suffix, given label $i$ at its start, with arbitrary legal values on dummy labels. Convexity gives $D_j(i)\in Y_j$.

Consider the stationary machine
\begin{equation}\label{eq:extractedmachine56}
                  \alpha_s,\qquad T_c=B_c,\qquad d_j=ED_j.
\end{equation}
Its decoder is legal. To test $w=c_1\cdots c_\ell$, $0\le\ell\le m$, use the first $\ell$ consecutive clique edges with kernels $B_{i-1,i,c_i}$ and then the single edge from $t_\ell$ to $t_{m+1}$ with kernel $A_{\ell,m+1}$. For $\ell=0$ there is only this last edge. Add the fixed prefix and suffix. This is a genuine word of length $N$; its physical product is obtained from $u_w$ by at most $\ell+3$ insertions from $V$. Its target output is consequently within $\eta$ of $F_{s,j}(u_w)$.

Each of its $\ell+1$ edge matrices differs from the corresponding matrix in \eqref{eq:extractedmachine56} by at most $\delta$. Stochastic multiplication contracts the $\ell^1$ distance between probability rows. A telescoping expansion, and subtraction of any fixed point of $Y_j$ from the decoder, therefore bound the difference between the actual output and $\alpha_s B_{c_1}\cdots B_{c_\ell}E D_j$ by
\[
                           D_*(\ell+1)\delta.
\]
Wordwise correctness of the original machine implies $e_{k,m}\le\varepsilon+\eta+D_*(m+1)\delta<\varepsilon+\gamma$, a contradiction. The representative $E$ is used only as part of the terminal decoder. No idempotence of $E$, consistency between different filler words, or hidden-state identity on a physical return was assumed.

The same argument shows that one may take
\begin{equation}\label{eq:ramseyloss56}
              L_{k,\varepsilon}=2g+(g+1)R_q(m+2).
\end{equation}
For exact cofinite returns, use exact fillers and omit $\eta$ and the neighborhood choice; $g$ is the given exact-return bound. In general $g$ depends on the finite witness and the continuity neighborhood. The remaining parameters depend only on the interface, the alphabet, $k$, and $\varepsilon$, never on a machine or horizon.

If $M_\varepsilon=\infty$, \eqref{eq:occupation56} for each $k$ gives $W_{N,\varepsilon}\to\infty$. If $M_\varepsilon=K<\infty$, the stationary machine gives $W_{N,\varepsilon}\le K$ for every $N$, and the bound for $k=K-1$ gives eventual equality when $K>1$. When $K=1$ equality holds for every $N$ because registers are nonempty.
\end{proof}

\begin{proposition}[Approximate returns without exact returns]\label{prop:approxreturns56}
Let the two circle commands have angles $\theta,\phi$, where $1,\theta/(2\pi),\phi/(2\pi)$ are rationally independent. They have no nonempty exact-return word, but satisfy Definition~\ref{def:returns56}.
\end{proposition}
\begin{proof}
A length-$n$ product has angle $n\phi+j(\theta-\phi)$, $0\le j\le n$. The positive orbit of the irrational rotation $\theta-\phi$ is dense. Compactness supplies, for every accuracy, a finite initial orbit segment that is a net of the circle. All longer initial segments remain nets; rotation by $n\phi$ preserves their accuracy. Thus a product is arbitrarily close to the identity at every sufficiently large length. An exact return would be a nontrivial integer relation among $1,\theta/(2\pi),\phi/(2\pi)$.
\end{proof}
A single irrational letter fails the eventual condition: there is only one word at each horizon. Its target can be placed in a one-label horizon-dependent decoder even when no finite stationary approximation at the same tolerance exists. Thus mere recurrence along a subsequence is not a substitute for Definition~\ref{def:returns56}.

\section{Stationary stochastic purification}\label{sec:purification55}
Write $\mathsf{Stoch}(k)$ for the compact semigroup of row-stochastic matrices of order $k$. The identity of a group inside this semigroup can be a singular idempotent; it must not be confused with the identity matrix.

\begin{lemma}[An idempotent over the physical identity]\label{lem:idempotent55}
Let $\mathcal T$ be a compact subsemigroup of $H\times\mathsf{Stoch}(k)$ whose projection onto $H$ is surjective. There is an element $e=(1,E)\in\mathcal T$ with $E^2=E$ and rank equal to the minimum matrix rank in $\mathcal T$. The monoid $e\mathcal T e$ is a compact group with identity $e$ and still projects onto $H$.
\end{lemma}
\begin{proof}
Only the integer ranks $1,\ldots,k$ occur, so their nonempty set has a minimum which is attained. Choose $(h,T)$ of that minimum rank $r$. The powers of $T$ are bounded, since they are stochastic. Thus its complex eigenvalues have modulus at most one and all Jordan blocks on the unit circle have size one: a larger peripheral Jordan block would produce polynomially growing powers. The closure of positive powers of an element of a compact group is a group: arbitrarily large positive powers approach its identity, and the preceding powers approach its inverse. Apply this observation to the compact monothetic group generated by $(h,\lambda_1,\ldots,\lambda_b)$ in $H\times(S^1)^b$, where the $\lambda_i$ are the peripheral eigenvalues of $T$. There is a net of integers tending to infinity for which $h^n\to1$ and all peripheral eigenvalues raised to the $n$th power tend to one. Along it $T^n$ tends to the spectral projection $E$ onto the peripheral subspace. This is a real stochastic idempotent. Its rank is at most $r$ and at least $r$ by minimality, so it is $r$. Nets avoid any metrizability assumption on $H$.

Every matrix $B$ in $e\mathcal T e$ satisfies $B=EBE$ and has rank $r$. Its action on the row space of $E$ is consequently invertible. Restriction embeds $e\mathcal T e$ as a compact submonoid of $H\times\mathrm{GL}(r,\R)$, with identity $(1,I_r)$; injectivity follows because the restriction determines $(xE)B=xB$ for every row $x$, by $B=EBE$, and hence determines every row of $B$. A compact submonoid of a topological group is a subgroup, by the same positive-power return argument. The inverse images of its inverses are still stochastic. Finally, sandwiching an element with physical coordinate $h$ between $e$'s leaves that coordinate unchanged. Surjectivity follows.
\end{proof}

\begin{lemma}[The recurrent simplex]\label{lem:simplex55}
If $E$ is a stochastic idempotent of rank $r$, then
\[
 \Delta_k E=\{p\in\Delta_k:pE=p\}
\]
is a simplex with exactly $r$ vertices. Every group of stochastic matrices with identity $E$ permutes these vertices.
\end{lemma}
\begin{proof}
Regard $E$ as a transition matrix. Its closed communicating classes each have a unique stationary probability row supported on that class. Every stationary row is a convex combination of these rows: mass on a transient class is zero, and stationarity within a closed irreducible class determines a one-dimensional space. For completeness, transient mass vanishes because each transient state has a positive probability of reaching a closed class in a bounded number of steps; a common bound over the finitely many states makes transient mass decrease by a fixed factor unless it is zero. The closed-class rows have disjoint supports, hence are linearly independent.

Each row of $E$ is stationary because $E^2=E$. Conversely, every stationary row equals its product with $E$. The span of the closed-class stationary rows is therefore precisely the row space of $E$, so there are $r$ classes and $r$ simplex vertices. If $B$ has a stochastic inverse relative to $E$, multiplication by $B$ and by this inverse are inverse affine maps of this simplex. Affine bijections permute its vertices.
\end{proof}

\begin{theorem}[Purification without a width or error loss]\label{thm:purification55}
For every compact-group interface in Definition~\ref{def:stationary55}, a stationary $k$-label realization of error at most $\varepsilon$ admits a permutation realization of error at most $\varepsilon$ on at most $k$ labels. There is no return-word hypothesis and no restriction on the number of connected components of $H$.
\end{theorem}
\begin{proof}
For the given machine form
\[
 \mathcal T=\overline{\{(u_w^{-1},T_w):w\in\mathcal A^*\}}
       \subset H\times\mathsf{Stoch}(k).
\]
The inverse on the physical coordinate is important: concatenation satisfies $u_{vw}^{-1}=u_v^{-1}u_w^{-1}$ and $T_{vw}=T_vT_w$. Thus $\mathcal T$ is a compact semigroup. Positive-word density makes its projection onto $H$ surjective. By continuity, its every element $(h,B)$ satisfies
\begin{equation}\label{eq:closederror55}
             \norm{\alpha_sBD_j-F_{s,j}(h^{-1})}_j\le\varepsilon
             \quad(s\in\mathcal S,\ j\in\mathcal J).
\end{equation}

Apply Lemma~\ref{lem:idempotent55} and denote the vertices of $\Delta_k E$ by $\nu_1,\ldots,\nu_r$. For each letter the element
\[
        e(u_a^{-1},T_a)e=(u_a^{-1},ET_aE)
\]
belongs to the group $e\mathcal T e$. By Lemma~\ref{lem:simplex55}, $B_a=ET_aE$ permutes the $\nu_i$. Use that permutation as the new command. If $V$ has rows $\nu_i$ and the induced row permutation is $\Pi_a$, the precise intertwining identity is $VB_a=\Pi_aV$. Write uniquely
\[
          \alpha_sE=\sum_{i=1}^r\beta_s(i)\nu_i,
          \qquad \widetilde D_j(i)=\nu_iD_j\in Y_j.
\]
The last inclusion uses convexity of the legal output set. For a nonempty word the output of the resulting $r$-label machine is $\alpha_s B_{a_1}\cdots B_{a_n}D_j$. For the empty word it is $\alpha_sED_j$. In both cases the corresponding joint element lies in $\mathcal T$ and has physical coordinate $u_w^{-1}$. Equation~\eqref{eq:closederror55} proves the error bound. Since $r\le k$, the width has not increased.
\end{proof}
The new machine need not reproduce the old machine's erroneous output word by word; it remains in the same closed error balls about the target. Nor do we assert $T_{w}=I$ on a physical return, invertibility of the original $T_a$, or deterministic initialization. These three incorrect shortcuts would invalidate the reduction.

\subsection{Finite quotients and the unrestricted stationary minimum}
For a tuple of permutation matrices $\sigma=(\Pi_a)_{a\in\mathcal A}$ of order $k$, set
\[
 \mathcal L_\sigma=
 \overline{\{(u_w^{-1},\Pi_w):w\in\mathcal A^*\}}
       \subset H\times\mathfrak S_k.
\]
The topology is the product of the compact topology on $H$ and the discrete topology on the finite group $\mathfrak S_k$. This is a compact subgroup. Its definition includes the possible discrepancy between physical relations and hidden-state relations.

\begin{theorem}[An exact stationary minimum]\label{thm:stationaryminimum55}
The value $M_\varepsilon$ is the least positive integer $k$ for which some tuple $\sigma$ of permutations of order $k$, some $\beta_s\in\Delta_k$, and some $d_j(i)\in Y_j$ satisfy
\begin{equation}\label{eq:permutationfeasibility55}
 \norm{\beta_s\Pi d_j-F_{s,j}(h^{-1})}_j\le\varepsilon
 \quad\text{for all }(h,\Pi)\in\mathcal L_\sigma\text{ and all }s,j.
\end{equation}
For fixed $k$ there are at most $(k!)^{|\mathcal A|}$ discrete choices (a crude bound before simultaneous conjugacy reduction), and the remaining feasible set is compact. The optimal stationary error at any fixed width is attained.
\end{theorem}
\begin{proof}
Necessity follows from Theorem~\ref{thm:purification55}, padding by unused labels if its output has fewer than $k$ states, and then taking the closure of word products. Sufficiency is the restriction of \eqref{eq:permutationfeasibility55} to words. The probabilities and legal decoders lie in a finite product of compact sets. Each displayed inequality defines a closed subset of that product. The worst error is continuous in these parameters, uniformly in the compact group variable, so its minimum exists. There are finitely many permutation tuples.
\end{proof}

\begin{theorem}[Finite observable quotient criterion]\label{thm:finitequotient55}
For arbitrary compact $H$ and the stated continuous interface,
\[
 M_\varepsilon<\infty
 \quad\Longleftrightarrow\quad
 r(U)\le\varepsilon\text{ for some open normal subgroup }U\le H.
\]
If such $U$ exists, $|\mathcal S|[H:U]$ deterministic-initialized permutation labels suffice. Conversely, a stationary $k$-label realization supplies such a $U$ with $[H:U]\le k!$.
\end{theorem}
\begin{proof}
Purify a stationary realization and use its joint group $\mathcal L_\sigma$. Put
\[
                  U=\{h\in H:(h,I)\in\mathcal L_\sigma\}.
\]
This is a closed normal subgroup of $H$: conjugate $(h,I)$ by an element above any $g\in H$, using surjectivity of the first projection. For a fixed permutation $\Pi$, a nonempty fiber in $\mathcal L_\sigma$ is a coset $hU$. These at most $k!$ fibers cover $H$, so $U$ has finite index. A closed subgroup of finite index in a compact group is open, since its complement is a finite union of closed cosets.

For each $s,j$ and each fiber, the single legal vector $\beta_s\Pi d_j$ is within $\varepsilon$ of every $F_{s,j}(h^{-1})$ in that fiber. Explicitly, $(hU)^{-1}=Uh^{-1}=h^{-1}U$ by normality, so inversion preserves the ordinary $U$-coset form. Hence $r(U)\le\varepsilon$. Conversely, store $(s,gU)$, update the coset by left multiplication by $u_a$, and decode a minimizing center of $F_{s,j}(gU)$. Compactness supplies the centers. This gives the asserted finite realization.
\end{proof}
For finitely many connected components, every open subgroup contains $K=H^\circ$, and $K$ itself is open normal. Therefore $M_\varepsilon<\infty$ holds exactly when
\[
 \varepsilon\ge\max_{s,j,c\in H/K}\rad_{Y_j}(F_{s,j}(c)),
\]
\emph{without} a return assumption. Executable returns enter the stronger clocked converse, not this stationary statement. The index bound $k!$ is an existence bound on a deterministic quotient, not a claim that the stationary minimum equals its index.

\begin{proposition}[Random initialization strictly saves states]\label{prop:C655}
Let $H=\Z/6\Z$, with an identity letter and a generator, one seed, and binary success probability
\[
 f(n)=\tfrac12+\tfrac14(-1)^n+\tfrac14\cos(2\pi n/3).
\]
At zero error the least deterministic component-quotient width is six, whereas the unrestricted stationary stochastic minimum is five.
\end{proposition}
\begin{proof}
The values at $n=0,\ldots,5$ are
\[
                  (1,\tfrac18,\tfrac58,\tfrac12,\tfrac58,\tfrac18).
\]
This sequence has least period six. A deterministic quotient of the cyclic component action preserving these outputs must therefore have six states. For a five-state realization take the disjoint union of a two-cycle and a three-cycle, both updating the index by $i\mapsto i+1$ modulo their respective lengths, initialize with mass $1/2$ at the zeroth vertex of each, and use decoder rows $(1,0)$ on the two-cycle and $(1,1/4,1/4)$ on the three-cycle. The identity command fixes all five labels. Direct averaging gives exactly $f(n)$, with legal binary probabilities.

If a stationary realization on at most four labels existed, purification would give one whose generator is a permutation of at most four letters. Its order is one of $1,2,3,4$, as follows by listing the possible cycle lengths. Every resulting numerical sequence has a period dividing that order, contradicting the least period six. Thus five is the unrestricted minimum.
\end{proof}

\section{Physical equivariance and the complete boundary}\label{sec:boundary56}
Return now to a compact Hausdorff group $H$ with a finite dense generating alphabet. The output sets may still be norm-compact convex subsets of real Banach spaces. Purification removes transition randomness; a further averaging removes hidden relations that have no physical counterpart.

\begin{theorem}[Physical-equivariant purification]\label{thm:equivariant56}
Every stationary $k$-label realization of error at most $\varepsilon$ admits one on $r\le k$ labels of the form
\begin{equation}\label{eq:equivariant56}
 \alpha_s\rho(h^{-1})d_j,\qquad h\in H,
 \qquad \rho:H\longrightarrow S_r\text{ a continuous homomorphism},
\end{equation}
with the same uniform error. Here permutations are represented by row-action matrices. In particular the command row is $\rho(u_a^{-1})$, all physical group relations hold on the new register, and initialization may remain random.
\end{theorem}
\begin{proof}
Apply Theorem~\ref{thm:purification55}, whose proof uses only finite convex combinations of decoder values. Denote the resulting initialization by $\beta_s$, its decoder by $D_j$, and its compact joint group by $\mathcal L\subset H\times S_k$, with the product topology and discrete permutation factor. Unused labels may be adjoined when necessary. Let $P$ be the second projection and
\[
                 N=\{\pi\in P:(1,\pi)\in\mathcal L\}.
\]
Then $N$ is normal in $P$. The averaging matrix $Q=|N|^{-1}\sum_{n\in N}n$ is stochastic, satisfies $Q^2=Q$, and commutes with $P$. For $(h,\pi)\in\mathcal L$, every $(h,n\pi)$ belongs to $\mathcal L$. Averaging their error inequalities and using convexity of the norm ball gives
\begin{equation}\label{eq:kernelaverage56}
                   \|\beta_sQ\pi D_j-F_{s,j}(h^{-1})\|_j\le\varepsilon.
\end{equation}

Let $O_1,\ldots,O_r$ be the $N$-orbits in the label set, and let $v_i$ be the uniform probability row on $O_i$. The rows of $Q$ are precisely these uniform orbit rows. Write $V$ for the matrix with rows $v_i$. Every $\pi\in P$ maps $N$-orbits bijectively onto $N$-orbits and therefore induces a permutation $\bar\pi$ with
\begin{equation}\label{eq:orbitintertwining56}
                             V\pi=\bar\pi V.
\end{equation}
The induced action of each $n\in N$ is trivial. Two elements of $\mathcal L$ with the same first coordinate differ by an element of $\{1\}\times N$. Hence $(h,\pi)\mapsto\bar\pi$ descends to a homomorphism $\rho:H\to S_r$. It is continuous: the surjection from the compact group $\mathcal L$ onto the Hausdorff group $H$ is a quotient map, and the original orbit action is continuous into a finite discrete group.

Set $\alpha_s(i)=\sum_{x\in O_i}\beta_s(x)$ and $d_j(i)=v_iD_j$. Then $\alpha_sV=\beta_sQ$ and $d_j=VD_j$. Equation~\eqref{eq:orbitintertwining56} gives
\[
                 \alpha_s\rho(h)d_j=\beta_sQ\pi D_j
                 \quad((h,\pi)\in\mathcal L).
\]
Replace $h$ by $h^{-1}$ in \eqref{eq:kernelaverage56}. This proves \eqref{eq:equivariant56}; all decoder values are legal by convexity, and $r\le k$.
\end{proof}

\begin{corollary}[Intrinsic stationary minimum]\label{cor:intrinsicminimum56}
The number $M_\varepsilon$ is the least cardinality of a finite continuous $H$-set for which some probability initialization for each seed and legal decoder for each query approximate the interface within $\varepsilon$. It is independent of the chosen finite dense alphabet. It is unchanged by replacing every target with
\[
                       h\longmapsto F_{s,j}(\varphi(h)q),
\]
where $\varphi$ is a continuous group automorphism and $q\in H$.
\end{corollary}
\begin{proof}
Theorem~\ref{thm:equivariant56} proves necessity, and restriction to command words proves sufficiency. For the last assertion, replace $\rho$ by $\rho\circ\varphi$ and $\alpha_s$ by $\alpha_s\rho(q^{-1})$ in \eqref{eq:equivariant56}. Applying the inverse change of variables proves equality of the minima.
\end{proof}
The finite quotient in Theorem~\ref{thm:finitequotient55} has a more precise description than its factorial bound. With $\mathcal L,P,N$ as above and $U=\{h:(h,I)\in\mathcal L\}$, the correspondence $(h,\pi)\mapsto(hU,\pi N)$ induces $H/U\simeq P/N$. Thus $[H:U]=|P/N|\le k!$. The orbit action in Theorem~\ref{thm:equivariant56} can have an additional kernel; no equality between this index and the minimal label count is asserted.

\begin{theorem}[Complete clocked finite-quotient classification]\label{thm:completeboundary56}
Assume eventual approximate returns on the compact group $H$. At every $\varepsilon\ge0$, including a positive critical radius, the following are equivalent:
\begin{enumerate}
\item $\sup_NW_{N,\varepsilon}<\infty$;
\item $M_\varepsilon<\infty$;
\item $r(U)\le\varepsilon$ for an open normal subgroup $U$ of $H$.
\end{enumerate}
For each $k<M_\varepsilon$ the all-profile occupation bound \eqref{eq:occupation56} holds, and $W_{N,\varepsilon}$ converges to $M_\varepsilon$. Put
\[
 R_0=\max_{s,j,h\in H}\rad_{Y_j}(F_{s,j}(hH^\circ)).
\]
Then $R_0=\inf_{U\lhd H\text{ open}}r(U)$. Above $R_0$ bounded realizations exist; below it every fixed width has bounded occupation. At $\varepsilon=R_0$ bounded clocked realization exists exactly when the infimum is attained by an open normal subgroup.
\end{theorem}
\begin{proof}
The first two conditions are equivalent by Theorem~\ref{thm:stationarization56}. The second and third are equivalent by Theorem~\ref{thm:finitequotient55}. That proof, including constrained center attainment, uses compactness and convexity of $Y_j$, not finite dimensionality.

For completeness, every open subgroup contains $H^\circ$: the image of the connected identity component in a finite discrete quotient is a singleton. Thus $r(U)\ge R_0$. The quotient $H/H^\circ$ is profinite, so its open normal subgroups form a neighborhood basis at the identity. The inverse images $U$ have intersection $H^\circ$. Given a neighborhood $O$ of $H^\circ$ in $H$, compactness supplies one such $U\subset O$: otherwise their intersections with the closed set $H\setminus O$ have the finite intersection property. Uniform continuity of the finite family of targets gives, for every $\eta>0$, a neighborhood $V$ of $1$ such that $\|F_{s,j}(hkv)-F_{s,j}(hk)\|_j<\eta$ for all $h\in H$, $k\in H^\circ$, and $v\in V$. Choose $U\subset H^\circ V$. Every $F_{s,j}(hU)$ is then within $\eta$ of $F_{s,j}(hH^\circ)$, so $r(U)\le R_0+\eta$. This proves the infimum formula and all remaining assertions.
\end{proof}
This result includes the formerly separate zero and positive boundary cases. The finite-rank proof of Theorem~\ref{thm:zero55} is retained below because it identifies an additional exact translate-space obstruction. No passage from an error strictly larger than $R_0$ to equality is used in Theorem~\ref{thm:completeboundary56}.

\subsection{Length phases without a width overhead}
Let $p\ge1$ be an integer and fix a command with product $a$. Let $J$ be the closure of all products of lengths divisible by $p$. The group argument of Lemma~\ref{lem:phasegroup55}, which uses only the block length, shows that $J$ is open normal, $H/J$ is cyclic of order dividing $p$, and all letters belong to $Ja$. Regard $\mathcal A^p$ as a finite alphabet on $J$. Assume that this block alphabet has eventual approximate returns. In particular this holds when the original nonempty exact-return language has greatest common divisor $p$.

For $0\le r<p$ let $M_{r,\varepsilon}$ be the intrinsic stationary minimum on $J$ for
\[
                         f_{r,s,j}(h)=F_{s,j}(ha^r).
\]
\begin{theorem}[Exact limiting width in each phase]\label{thm:phaseminimum56}
Under the preceding assumptions, for every $r$ and $\varepsilon\ge0$,
\begin{equation}\label{eq:phaseminimum56}
 W_{N,\varepsilon}\le M_{r,\varepsilon}\quad(N\equiv r\pmod p),
 \qquad
 \lim_{\substack{N\to\infty\\N\equiv r\ ({\rm mod}\ p)}}W_{N,\varepsilon}
          =M_{r,\varepsilon}.
\end{equation}
Finite limits are eventually attained. If $k<M_{r,\varepsilon}$, the number of positive cuts of width at most $k$ is bounded uniformly over that horizon phase, with no restriction on other widths. Boundedness in the phase is equivalent to the existence of an open normal subgroup $U\le J$ such that
\begin{equation}\label{eq:phasequotient56}
          \max_{s,j,h\in J}\rad_{Y_j}(F_{s,j}(hUa^r))\le\varepsilon.
\end{equation}
\end{theorem}
\begin{proof}
Suppose first that $M_{r,\varepsilon}=K<\infty$ and choose a physical-equivariant stationary realization of $f_r$ on $K$ labels, with representation $\rho:J\to S_K$. For a fixed $N\equiv r\pmod p$, let $g_t$ be the physical product after $t$ commands and set
\begin{equation}\label{eq:gauge56}
 z_t=a^{N-t}g_ta^{-r}\in J,
 \qquad c_{t,b}=a^{N-t}u_ba^{-(N-t+1)}\in J.
\end{equation}
Normality of $J$ and $u_ba^{-1}\in J$ justify the membership assertions. We have $z_0=a^{N-r}$, $z_t=c_{t,b}z_{t-1}$ when the $t$th letter is $b$, and $z_N=g_Na^{-r}$. Initialize with $\alpha_s\rho(z_0^{-1})$, use the clock-dependent permutation $\rho(c_{t,b}^{-1})$, and keep decoder $d_j$. The row product is $\alpha_s\rho(z_N^{-1})$, so its error from $F_{s,j}(g_N)=f_{r,s,j}(z_N)$ is at most $\varepsilon$. This uses the free external clock but no phase register or partial-block storage. It proves the upper bound with exactly $K$ labels.

Fix $k<M_{r,\varepsilon}$ and a cut residue $0\le\ell<p$. Put $b=(r-\ell)\bmod p$, chosen in $\{0,\ldots,p-1\}$. Restrict the input to the fixed prefix $a^\ell$, arbitrary $p$-letter blocks, and fixed suffix $a^b$. Integrate the prefix into initialization and the suffix into the decoder. The block target on $J$ is
\[
 G_{\ell,s,j}(h)=F_{s,j}(a^b h a^\ell)
     =f_{r,s,j}\bigl(\operatorname{Ad}_{a^b}(h)a^{\ell+b-r}\bigr).
\]
Here $a^{\ell+b-r}\in J$ and conjugation by $a^b$ is an automorphism of $J$. Corollary~\ref{cor:intrinsicminimum56} therefore identifies its stationary minimum with $M_{r,\varepsilon}$. Theorem~\ref{thm:stationarization56} on the block alphabet bounds its positive narrow cuts by a constant $L_\ell$. These cuts are precisely the original cuts in residue $\ell$ between the prefix and suffix, apart from the initial block cut. At most  $\ell+b+1\le2p$ further cuts need be charged for this residue, including small horizons shorter than the chosen prefix and suffix. Summing gives the uniform bound $\sum_{\ell=0}^{p-1}L_\ell+2p^2$. As before, a peak at most $k$ would force $N$ below this bound. This proves the limit and occupation assertions. Finally apply Theorem~\ref{thm:finitequotient55} to $f_r$ on $J$ to obtain \eqref{eq:phasequotient56}.
\end{proof}
If the order of $H/J$ is a proper divisor $q$ of $p$, the minima and quotient-radius criteria repeat with period $q$: multiplication by $a^q\in J$ is a right translation of the target. The earlier phase radii have the same repetition. Neither the physical quotient order nor the return-length greatest common divisor is substituted silently for the other.

\subsection{Compact outputs and finite hidden state}
Theorems~\ref{thm:stationarization56}, \ref{thm:equivariant56}, \ref{thm:completeboundary56}, and \ref{thm:phaseminimum56} hold for norm-compact convex Banach-valued outputs. A single-machine purification even needs only convexity: its new decoder values lie in the convex hull of finitely many old values. Compact output sets are used instead in finite-witness extraction, optimal-error attainment, and constrained-center existence. The representative-function and algebraic sections below retain their explicitly finite-dimensional hypotheses.

Finiteness on the hidden side has a different role. On countably many labels the full simplex is not compact in total variation; mass can escape to infinity, matrix ranks need not admit a useful finite minimum, and a recurrent simplex may have infinitely many vertices. On a general compact convex hidden state space affine automorphisms need not permute a finite vertex set: rotation of the disk already acts continuously on infinitely many extreme points. Thus these proofs give neither a finite-label purification theorem for arbitrary operator kernels nor an online observation theorem. Their terminal-output assertion is strictly the one stated in the models.

\section{The clocked component-radius threshold}\label{sec:intro54}
A finite stochastic register can carry infinitely many probability distributions. Thus an infinite numerical orbit is not, by itself, an obstruction to finite positive realization. A further difficulty arises when the realization may be redesigned at every horizon and its transitions may depend on a free external clock. The obstruction established here survives both freedoms. It is measured by the radius of an observable image inside its allowed output set.

For the clocked theorem in this section, let $H$ be a compact Hausdorff group with finitely many connected components. Put $K=H^\circ$ and $C=H/K$. Then $K$ is an open normal subgroup and $C$ is a finite group. Fix a finite alphabet $\mathcal A$ with elements $u_a\in H$ whose generated group is dense in $H$. Our convention is
\[
 u_{a_1\cdots a_n}=u_{a_n}\cdots u_{a_1}.
\]
The identity $e$ need not be an alphabet element. Instead assume the following executable return condition:
\begin{equation}\label{eq:returns54}
 \text{there is an integer }g\ge0\text{ such that, for every }n\ge g,
 \text{ some word of length }n\text{ has product }e.
\end{equation}
An identity letter gives $g=0$. The condition concerns physical products; the stochastic rows on a return word may perform arbitrary processing.

For nonempty finite seed and query sets $\mathcal S,\mathcal J$, let
\[
 F_{s,j}:H\longrightarrow Y_j
\]
be continuous, where $Y_j$ is a nonempty compact convex subset of a finite-dimensional real normed space $(V_j,\norm{\cdot}_j)$. A query is selected only after the command word. A prescribed joint law of finitely many outputs is treated as one query, not as a collection of marginal constraints. In particular, a categorical interface has
\[
 Y_j=\Delta_{M_j}=\{p\in\R^{M_j}:p_i\ge0,\ \textstyle\sum_i p_i=1\},
 \qquad \norm{x}_j=\tfrac12\sum_i|x_i|.
\]
Then $\norm{p-q}_j$ is total variation. Binary means correspond to $Y=[-1,1]$ with $\norm{x}=|x|/2$.

\begin{definition}[Clocked realization]\label{def:machine54}
At horizon $N$ choose finite nonempty registers $S_t$, $0\le t\le N$, initialization rows $E_s$, row-stochastic command matrices $T_{t,a}:S_{t-1}\to S_t$, and terminal decoder maps $D_j:S_N\to Y_j$. The realized output vector on $w=a_1\cdots a_N$ is
\[
 E_sT_{1,a_1}\cdots T_{N,a_N}D_j.
\]
Its error is the maximum, over all seeds, length-$N$ words, and queries, of its distance from $F_{s,j}(u_w)$. The command width is $\max_t|S_t|$, counting all declared labels, including unused padding labels. Write $W_{N,\varepsilon}$ for its least possible value at error at most $\varepsilon$.
\end{definition}
All retained private information, including randomness kept for future use, is charged to the current label. The rows may depend on $t$ and $N$, but not on an unrecorded input history or on the future query. The external clock and horizon, arbitrary real table descriptions, their construction and lookup, real arithmetic, and exact sampling are uncharged. This is a numerical realization invariant, not uniform workspace or program size. For a finite categorical interface, charging the terminal answer register changes the peak to $\max(W_{N,\varepsilon},\max_j M_j)$. Thus all boundedness and growth results are unchanged; the convention that also charges a sampled binary-answer register has a minimum of two labels. A decoder vector by itself is not an additional sampled register.

For $A\subset Y_j$ define its constrained radius by
\[
 \rad_{Y_j}(A)=\min_{y\in Y_j}\sup_{x\in A}\norm{x-y}_j.
\]
The minimum is attained when $A$ is compact. The critical error is
\begin{equation}\label{eq:radius54}
 \epsc=\max_{s,j,c\in C}\rad_{Y_j}\bigl(F_{s,j}(c)\bigr).
\end{equation}
Here a component is identified with its coset in $H$; its image is compact by continuity. The requirement that the center belong to $Y_j$ is essential for categorical outputs.

\begin{theorem}[Intrinsic radius classification]\label{thm:radius54}
Under \eqref{eq:returns54}, for every $\varepsilon\ge0$ the following are equivalent:
\begin{enumerate}
\item $\sup_N W_{N,\varepsilon}<\infty$;
\item one stationary finite-state realization with permutation command updates has error at most $\varepsilon$ at every horizon;
\item $\varepsilon\ge\epsc$.
\end{enumerate}
At the boundary, at most $|\mathcal S||C|$ command labels suffice. If $\varepsilon<\epsc$, then for every integer $k\ge1$ there is an integer $L_k<\infty$ such that every error-$\varepsilon$ realization satisfies
\begin{equation}\label{eq:occupation54}
 \#\{t\in\{1,\ldots,N\}:|S_t|\le k\}\le L_k.
\end{equation}
All intervening widths are unrestricted. In particular, $W_{N,\varepsilon}\to\infty$.
\end{theorem}
Neither a positive lower bound on state masses nor a reachable-rank, observability, or conditioning assumption occurs in the theorem. The upper realization is deterministic until its terminal decoder. It need not have the minimum width among all stochastic realizations. Theorem~\ref{thm:stationaryminimum55}, rather than a deterministic component partition, characterizes that stationary minimum.

For binary means, \eqref{eq:radius54} reduces to one quarter of the largest componentwise oscillation. For a genuine categorical law it is generally not half its diameter; an analytic three-outcome example below has critical total-variation error $2/3$, whereas every individual event has binary threshold $1/2$. The proof uses finitely many localized vector tests rather than a separate midpoint for each coordinate.

The finite-component assumption does not assert that $H$ is a Lie group. For example, take $H=\prod_{n\ge1}\R/\Z$, choose an element whose every finite coordinate tuple together with $1$ is rationally independent, and include this element and $e$ as commands. The positive orbit is dense and $H$ is connected but has no faithful finite-dimensional representation. The theorem applies to every finite continuous interface on this group. With infinitely many components the component construction need not be finite; Section~\ref{sec:profinite55} gives the finite-observable-quotient criterion and the strict-side clocked theorem in that setting.

There is also a quantitative counterpart. For planar rotations with a dual badly approximable $s$-tuple of angles and signal amplitude $0<\rho<1$, Theorem~\ref{thm:matched54} proves
\[
 W_{N,\varepsilon}=\Theta\bigl(N^{s/(2s+1)}\bigr)
 \quad\text{for every fixed }0\le\varepsilon<\rho/2.
\]
The upper realization is exact even when $\varepsilon>0$. The new lower estimate reaches the entire subcritical interval, by combining harmonic localization with the arithmetic orbit-packing mechanism. This quantitative theorem and the radius classification have distinct hypotheses: algebraic matrix entries alone do not imply the badly approximable angle condition.

\subsection{Executable tests}
Let $\mathcal P=\{u_w:w\in\mathcal A^*\}$. We use standard Haar averaging and the density of matrix coefficients of finite-dimensional unitary representations, in realified form \cite{Folland}. No quantitative mixing assumption for a preassigned random walk is imposed.

\begin{lemma}[Positive words and common lengths]\label{lem:positive54}
The semigroup $\mathcal P$ is dense in $H$, and $\mathcal P\cap K$ is dense in $K$. Every finite subset of $\mathcal P\cap K$ has executable representatives of a common positive length.
\end{lemma}
\begin{proof}
For an element $u$ of a compact group, positive powers return arbitrarily close to $e$ with exponent at least two. Indeed cover the compact closure of its powers by finitely many translates of a sufficiently small neighborhood. Among sufficiently many powers, three lie in one member of the cover; the first and last exponents differ by at least two, and their quotient is arbitrarily close to $e$. This argument also covers finite-order elements by taking a sufficiently large multiple of the order. If $u^r$ approaches $e$ with $r\ge2$, then $u^{r-1}=u^ru^{-1}$ approaches $u^{-1}$ and is still a positive power. Thus $u^{-1}$ is in the closure of positive powers. The closed semigroup $\overline{\mathcal P}$ therefore contains the group generated by the letters, and hence equals $H$. Since $H$ has finitely many components, $K$ is open, and a relatively open subset of $K$ is open in $H$; density gives the second assertion. For finitely many representative lengths $\ell_i$, choose $B\ge\max_i\ell_i+g$, $B\ge1$, and append return words of lengths $B-\ell_i$. The products are unchanged.
\end{proof}

\begin{lemma}[Finite executable gap]\label{lem:gap54}
Let $R:K\to O(E)$ be a continuous representation on a nonzero finite-dimensional Euclidean space, with no nonzero fixed vector. For every $k\ge1$ there are $B\ge1$, a probability law $\nu$ on finitely many executable length-$B$ words with products in $K$, and $d\in(0,1)$ such that
\begin{equation}\label{eq:gap54}
 \E_{w\sim\nu}\max_{1\le i\le k}\ip{v_i}{R(u_w)u}\le1-d
\end{equation}
for all unit vectors $u,v_1,\ldots,v_k\in E$.
\end{lemma}
\begin{proof}
On the compact product of unit spheres consider the minimum
\[
 \gamma=\min_{u,v_1,\ldots,v_k}\int_K
 \left(1-\max_i\ip{v_i}{R(h)u}\right)\,d\mu(h).
\]
Repeated centers are permitted. The integrand is nonnegative. If a minimizing tuple had integral zero, full support of Haar measure would give $R(h)u\in\{v_1,\ldots,v_k\}$ for every $h$. The orbit of $u$ is connected, hence a singleton, contradicting the fixed-vector hypothesis. Thus $\gamma>0$.

Partition the compact image $R(K)$ into finitely many measurable sets of small diameter in the Euclidean operator norm, and replace their Haar masses by atoms at nearby executable products in $K$. Cells of zero Haar mass can be discarded. The function $A\mapsto\max_i\ip{v_i}{Au}$ is uniformly $1$-Lipschitz in that norm. An approximation radius less than $\gamma/2$ therefore preserves a positive defect uniformly in the unit vectors. Lemma~\ref{lem:positive54} gives a common length. Decreasing the defect to a number in $(0,1)$ proves the assertion.
\end{proof}

\begin{lemma}[Conditional-centroid contraction]\label{lem:centroid54}
Let $Z\in E$ be a bounded random physical vector and $S$ the current register. Apply an independent word $w\sim\nu$ from Lemma~\ref{lem:gap54}, and let the actual block kernel generate a terminal state $S'$ having at most $k$ labels. Then
\begin{equation}\label{eq:centroid54}
 \E\norm{\E[R(u_w)Z\mid S']}\le(1-d)\E\norm{\E[Z\mid S]}.
\end{equation}
A fixed intervening return word cannot increase the same conditional norm, regardless of its intermediate or terminal widths.
\end{lemma}
\begin{proof}
Put $z_s=\E[Z\mid S=s]$ and write $T_w(s,i)$ for the actual composite row of the block. Since this row depends on the past only through $S$,
\[
 \Pr(S'=i)\E[R(u_w)Z\mid S'=i]
 =\sum_s\Pr(S=s)\E_w T_w(s,i)R(u_w)z_s.
\]
Choose $v_i$ in the direction of each nonzero terminal centroid, arbitrarily otherwise. Pair, sum in $i$, and use stochasticity to bound the sum by
\[
 \sum_s\Pr(S=s)\E_w\max_i\ip{v_i}{R(u_w)z_s}
 \le(1-d)\sum_s\Pr(S=s)\norm{z_s}.
\]
Zero-probability states contribute zero. No restriction was placed on registers inside the block. On a fixed return word the physical vector is unchanged at its endpoints; the new state is a stochastic image of the old one, so conditional Jensen gives nonincrease. Independence is used only at complete block boundaries.
\end{proof}

\section{Representative functions and localized radii}\label{sec:local54}
A real representative function on $K$ is a matrix coefficient of a finite-dimensional continuous orthogonal representation, or a finite sum of such coefficients. Their algebra contains the constants, is closed under multiplication by tensor products, and is uniformly dense in $C(K,\R)$ by the Peter--Weyl approximation theorem \cite{Folland}. A finite collection of representative functions can be written on one space
\[
 \mathcal V=\bigoplus_{i=1}^a\End(E_i),\qquad
 \Phi(h)=(R_i(h))_{i=1}^a,
\]
where the scalar product is the sum of Frobenius products. Left multiplication defines an orthogonal representation $R$ on $\mathcal V$ and satisfies $\Phi(gh)=R(g)\Phi(h)$. Adjoin a trivial summand when necessary. Let
\[
 P=\int_KR(h)\,d\mu(h),\quad E=(I-P)\mathcal V,\quad
 Z(h)=(I-P)\Phi(h),\quad Q_0=\norm{Z(e)}.
\]
Haar averaging is the orthogonal projection onto fixed vectors. Consequently $Z(gh)=R(g)Z(h)$, the representation on $E$ has no fixed vectors, and $P\Phi(h)=P\Phi(e)$ for $h\in K$.

Choose and fix a coefficient vector $a_p$ for each scalar representative function $p$, and a linear coefficient map $A_G$ for each vector-valued representative function $G$. Their norms need not be intrinsic; fixing any such choices suffices. For an arbitrary joint distribution of $h\in K$ and a finite state $S$, put $Q=\E\norm{\E[Z(h)\mid S]}$ and $\bar p=\int p\,d\mu$. The fixed/moving decomposition gives
\begin{align}
 \E[p(h)\mid S]&=\bar p+\ip{a_p}{\E[Z(h)\mid S]},\label{eq:calibration54}\\
 |\E p-\bar p|&\le\norm{a_p}Q,\qquad
 \norm{\E G-\bar G}\le\norm{A_G}Q.\label{eq:vectorcal54}
\end{align}
Here the norm of $A_G$ is the operator norm from the Euclidean feature space to the indicated output space. If $D:S\to Y$ and $M_Y=\max_{y\in Y}\norm{y}$, then
\begin{equation}\label{eq:decodercal54}
 \norm{\E[p(h)D(S)]-\bar p\E D(S)}\le M_Y\norm{a_p}Q.
\end{equation}
Indeed condition on $S$ in \eqref{eq:calibration54}, multiply by $D(S)$, and apply the triangle inequality. These are moment estimates, not total-variation convergence of a hidden physical distribution.

\begin{lemma}[Finite localization of a radius]\label{lem:radiuslocal54}
Let $f:K\to Y$ be continuous, with $Y$ compact and convex in a finite-dimensional normed space. Suppose $\rad_Y(f(K))>\varepsilon$. There are a vector-valued representative function $G$, a number $\eta>0$, and nonnegative representative functions $p_1,\ldots,p_a$ of Haar mean one such that
\[
 \norm{G-f}_\infty<\eta,\qquad
 r_*:=\min_{y\in Y}\max_i\norm{\bar{p_iG}-y}>\delta:=\varepsilon+\eta.
\]
\end{lemma}
\begin{proof}
Choose $\varepsilon<t<\rad_Y(f(K))$. For each $y\in Y$ some $h\in K$ has $\norm{f(h)-y}>t$. The corresponding open subsets of $Y$ cover $Y$, so a finite subcover yields $h_1,\ldots,h_a$ with
\[
 \min_{y\in Y}\max_i\norm{f(h_i)-y}>t.
\]
The strict inequality follows also at a minimizing $y$, since the maximum is everywhere greater than $t$ and $Y$ is compact.

Fix a sufficiently small positive $\eta$. Approximate the coordinates of $f$ uniformly by representative functions to obtain $\norm{G-f}_\infty<\eta$. For each $i$, choose a nonzero nonnegative continuous bump supported in a neighborhood where $f$ differs from $f(h_i)$ by less than $\eta$. Such a bump exists by normality of the compact Hausdorff space; it has positive Haar integral by full support. Approximate its square root uniformly by a representative function, square the approximation, and normalize its Haar integral. The resulting nonnegative representative function $p_i$ has mean one, and its weighted $f$-mean can be made arbitrarily close to that of the bump. Thus we may ensure
\[
 \norm{\bar{p_iG}-f(h_i)}<3\eta.
\]
The constrained radius changes by at most the maximum perturbation of the finitely many points. Taking $4\eta<t-\varepsilon$ gives $r_*>\varepsilon+\eta$.
\end{proof}

\begin{proposition}[Residual moment from a vector interface]\label{prop:residual54}
Fix the data of Lemma~\ref{lem:radiuslocal54}. Choose one feature space representing all $p_i$ and $p_iG$, and fix their coefficient choices before defining
\[
 L=\max_i\bigl(\norm{A_{p_iG}}+(M_Y+\delta)\norm{a_{p_i}}\bigr),
 \qquad \Delta=r_*-\delta>0.
\]
Suppose a terminal decoder $D(S)\in Y$ obeys
\[
 \norm{\E[D(S)\mid w]-G(h(w))}\le\delta
\]
for every complete deterministic word in a test law. Then $L>0$, $E\ne\{0\}$, $Q_0>0$, and
\begin{equation}\label{eq:residual54}
                         Q\ge\Delta/L.
\end{equation}
\end{proposition}
\begin{proof}
For $p=p_i\ge0$, wordwise correctness and the triangle inequality imply
\[
 \norm{\E[p(h)D(S)]-\E[p(h)G(h)]}\le\delta\E p(h).
\]
This is a conditional mean error, not a pointwise bound on the random decoder. Since $\bar p_i=1$, calibration gives explicitly
\[
                  \E p_i\le1+\norm{a_{p_i}}Q.
\]
Write $y=\E D(S)$, which belongs to $Y$ by convexity. Equations~\eqref{eq:vectorcal54}--\eqref{eq:decodercal54}, together with $\bar p=1$, give
\[
 \norm{y-\bar{p_iG}}\le\delta+
 \bigl(\norm{A_{p_iG}}+(M_Y+\delta)\norm{a_{p_i}}\bigr)Q.
\]
Maximizing in $i$ yields $r_*\le\delta+LQ$. Hence $LQ\ge\Delta>0$. In particular the moving feature and its initial norm cannot vanish: by equivariance, $Q_0=0$ would imply $Z(h)=0$ for all $h$ and therefore $Q=0$.
\end{proof}

\subsection{Proof of Theorem~\ref{thm:radius54}}
Choose $(s,j,c)$ with component radius greater than $\varepsilon$. By density there is an executable prefix $v$ with $u_v\in c$. Put $b=|v|$ and $f(h)=F_{s,j}(hu_v)$ on $K$. Since $Ku_v=c$, this map has the same constrained radius. Apply Lemma~\ref{lem:radiuslocal54} and Proposition~\ref{prop:residual54}.

The moving feature is nonzero independently of any particular test law. Otherwise all represented functions would be constant on $K$; since each $p_i$ has mean one, this would force $p_i=1$ and $G$ constant. Then all $\bar{p_iG}$ would equal a point at distance less than $\eta$ from $Y$, contradicting $r_*>\delta\ge\eta$.

For a fixed width $k$, apply Lemma~\ref{lem:gap54} to the resulting moving feature. Let $B,d$ be its block parameters and set
\[
 \chi=-\log(1-d)>0,\qquad A=\max(1,LQ_0/\Delta).
\]
Start a test word with $v$; between independent length-$B$ test blocks use fixed return words. The accumulated physical product is $hu_v$. Initially $h=e$, so the conditional feature norm after the prefix equals $Q_0$, regardless of the hidden-state distribution. If $J$ block endpoints have width at most $k$, Lemma~\ref{lem:centroid54} gives
\[
 Q_N\le Q_0(1-d)^J.
\]
The machine is correct on every deterministic word of length $N$. Thus any external law supported on such words inherits the conditional mean bound, with error $\delta=\varepsilon+\eta$ for $G$. Proposition~\ref{prop:residual54} implies
\begin{equation}\label{eq:J54}
                         J\chi\le\log A.
\end{equation}
The law depends on the advertised width profile, not on the machine's private state. When $A=1$, \eqref{eq:J54} excludes even one selected narrow block; no logarithmic loss has occurred.

For clarity we give the spacing count with the full return hypothesis. Discard narrow cuts before $b+B+g$ and after $N-g$. From the remaining cuts greedily choose the first and then the next at distance at least $B+g$. Put an independent block immediately before each selected cut. The initial gap, all intervening gaps, and the final suffix have lengths at least $g$, so \eqref{eq:returns54} fills them with executable returns. At most $b+B+2g-1$ cuts were discarded. Each chosen endpoint accounts for at most $B+g$ eligible cuts. Therefore one may take
\begin{equation}\label{eq:Lk54}
 L_k=\left\lceil b+B+2g-1+(B+g)\frac{\log A}{\chi}\right\rceil.
\end{equation}
The constants $B,d$ depend on the fixed moving representation and $k$; $\chi=-\log(1-d)$ depends on that defect. The localization, prefix length $b$, and $A$ depend only on the alphabet, interface, chosen component, and error. Consequently $L_k$ depends on these fixed data, the return certificate $g$, and $k$, but not on $N$ or the machine. When the eligible interval is empty the same estimate follows directly by counting. This proves occupation without restricting any register inside a block or gap. If the entire peak is at most $k$, then $N\le L_k$, which proves divergence directly. Monotonicity in the horizon is not needed when an identity letter is absent.

For the upper bound, use labels $(s,c)\in\mathcal S\times C$, including any labels not reached at a given horizon. Initialize at $(s,K)$ and update by left multiplication
\[
 (s,c)\longmapsto(s,[u_a]c).
\]
Choose a minimizing center of $F_{s,j}(c)$ in $Y_j$ as decoder at $(s,c)$ for query $j$. These are legal decoder values, and the error is at most \eqref{eq:radius54}. All command updates are permutations of one finite state set, independent of clock and horizon. This proves boundary attainment and the equivalence.\qed

\begin{proposition}[A finite return certificate]\label{prop:return54}
Condition~\eqref{eq:returns54} follows if there are finitely many executable identity-product words whose positive lengths have greatest common divisor one.
\end{proposition}
\begin{proof}
Let the lengths be $\ell_1,\ldots,\ell_a$. Their nonnegative sums modulo $\ell_1$ form the subgroup generated by their residues: a subsemigroup of a finite group is a subgroup. The gcd hypothesis makes this the whole residue group. Choose one nonnegative sum $n_r$ for each residue $r$ modulo $\ell_1$, and put $g=\max_r n_r$. Every $n\ge g$ equals its representative $n_r$ plus a nonnegative multiple of $\ell_1$. Concatenating the corresponding return words proves the claim.
\end{proof}
For example, the alphabet $\{R_\theta,R_{-\theta},R_{2\pi/3}\}$, with $\theta/2\pi$ irrational, contains no identity letter but has return lengths two and three. It satisfies \eqref{eq:returns54} with $g=2$. A single irrational rotation letter does not: there is only one word at each horizon, and a horizon-specific constant decoder can store its numerical answer with one command label. Thus synchronization cannot be silently omitted. Nor is a bounded irreversible semigroup enough: for the identity command $I=1$ and contraction command $a=\lambda\in(0,1)$, the mean $\rho\lambda^{\#a(w)}$ has an exact alive/absorbed two-state realization. The group generated with inverses is not bounded. The present theorem is not a compact-semigroup classification.

\section{Exact-return languages and length phases}\label{sec:phases55}
The existence of cofinite returns is not the only useful synchronization regime. In fact an exact-return language, when nonempty, has a rigid eventual form. Here and in this section a return has positive length.

\begin{lemma}[The return period]\label{lem:returnperiod55}
Suppose
\[
 \mathcal R=\{|w|>0:u_w=1\}\ne\varnothing,
 \qquad p=\gcd\mathcal R.
\]
There is an integer $g_p$ such that every multiple $dn$ with $n\ge g_p$ belongs to $\mathcal R$, and no length outside $p\N$ belongs to $\mathcal R$. In particular, every nonempty exact-return language is eventually one arithmetic progression.
\end{lemma}
\begin{proof}
Concatenation makes $\mathcal R$ additive. A finite subset of its elements already has gcd $p$: start with one length and successively decrease the gcd when necessary; a strictly decreasing chain of its positive divisors is finite. Divide these finitely many lengths by $p$. Their gcd is one, so the residue argument of Proposition~\ref{prop:return54} shows that their nonnegative sums contain every sufficiently large integer. Multiplying by $p$ proves the assertion. The exclusion of other lengths is the definition of the gcd.
\end{proof}

Fix a letter with product $a\in H$, and let $J$ be the closure of the products of words whose lengths are multiples of $p$, including the empty word.
\begin{lemma}[The phase subgroup]\label{lem:phasegroup55}
The group $J$ is an open normal subgroup of $H$. The quotient $H/J$ is cyclic of order dividing $p$, generated by $aJ$, and
\begin{equation}\label{eq:phasereachable55}
 \overline{\{u_w:|w|\equiv r\pmod p\}}=a^rJ
 \qquad(0\le r<p).
\end{equation}
If $H/H^\circ$ is finite, then $J^\circ=H^\circ$.
\end{lemma}
\begin{proof}
The compact closure defining $J$ is a group, by the positive-power argument, and is generated densely by the block alphabet $\mathcal A^p$. Also $a^p\in J$. If $x=u_w$ for a word $w$ of length $p$, then the executable word consisting of $p-1$ copies of the chosen letter, followed by $w$, followed by one more copy, has product $axa^{p-1}$ and length $2p$. Thus $axa^{p-1}\in J$, whence $axa^{-1}\in J$. It follows that $aJa^{-1}\subset J$. Applying this inclusion $p$ times and using $a^p\in J$ shows equality. Every letter $b$ has $ba^{p-1}\in J$, so $ba^{-1}\in J$. Thus all letters lie in $Ja$, normalize $J$, and their products of length $n$ lie in $Ja^n$. The finite union $\bigcup_{r=0}^{p-1}Ja^r$ is a closed group containing the letters, hence is $H$. This proves normality, finite index, and openness. Conversely, words consisting of $r$ copies of $a$ followed by arbitrary $p$-blocks have products dense in $Ja^r$, proving \eqref{eq:phasereachable55}. An open subgroup contains the identity component, and its own identity component is therefore the same one.
\end{proof}
The quotient order may be a proper divisor of $p$; the phase radii repeat whenever the physical cosets repeat. The phase index records executable length, not merely the physical component quotient.

Assume now that $H$ has finitely many components and put $K=H^\circ$. For $r\in\{0,\ldots,p-1\}$ define
\begin{equation}\label{eq:phaseradius55}
 \varepsilon_r=\max_{s,j,\ c\in H/K:\ c\subset a^rJ}
                    \rad_{Y_j}(F_{s,j}(c)).
\end{equation}
\begin{theorem}[The phasewise clocked threshold]\label{thm:phases55}
Under the sole synchronization assumption $\mathcal R\ne\varnothing$, for every $r$ and $\varepsilon\ge0$,
\[
 \sup_{N\equiv r\ (\mathrm{mod}\ p)}W_{N,\varepsilon}<\infty
             \quad\Longleftrightarrow\quad
             \varepsilon\ge\varepsilon_r.
\]
Equality is attained by a finite component-permutation realization on the indicated horizons. If $\varepsilon<\varepsilon_r$, then for each $k$ there is a constant $L_{k,r}$, independent of $N$, such that every error-$\varepsilon$ realization at $N\equiv r\pmod p$ satisfies
\[
                 \#\{1\le t\le N:|S_t|\le k\}\le L_{k,r}.
\]
The intervening widths are unrestricted. Consequently the full family is bounded exactly at and above $\max_r\varepsilon_r=\varepsilon_{\mathrm c}$; below it, divergence is asserted phase by phase, not necessarily along every horizon.
\end{theorem}
\begin{proof}
Every product of length $N\equiv r$ lies in $a^rJ$. Store its component and the seed and decode component centers as before. This proves the upper bound, including equality, with at most $|\mathcal S||H/K|$ labels. Centers on other components may be chosen arbitrarily when only this phase is under consideration.

For the converse fix a component $c\subset a^rJ$ and $s,j$ with radius greater than $\varepsilon$. We must count narrow cuts in every residue class, not only at block boundaries of one phase. For each $\ell\in\{0,\ldots,p-1\}$ choose $b_\ell\in\{0,\ldots,p-1\}$ congruent to $r-\ell$ modulo $p$. By \eqref{eq:phasereachable55} and the openness of components there is a fixed word $v_\ell$, of length $p_\ell\equiv\ell\pmod p$, with product in $a^{-b_\ell}c$. Indeed this component is contained in $a^\ell J$. Consider the block experiment on $J$ with alphabet $\mathcal A^p$ and target
\[
        G_\ell(h)=F_{s,j}(a^{b_\ell}h u_{v_\ell}),\qquad h\in J.
\]
It is a legal compact-convex interface. Its image on $K=J^\circ$ is $F_{s,j}(c)$, since $K$ is normal. Lemma~\ref{lem:returnperiod55} gives cofinite exact returns in block length for this alphabet.

Restrict the original machine to words that begin with $v_\ell$, end with $b_\ell$ copies of the chosen letter, and have $M=(N-p_\ell-b_\ell)/p$ arbitrary $p$-blocks in between. When $M\ge0$, absorb the fixed prefix into the initialization and the fixed suffix into the terminal decoder; stochastic averaging keeps that decoder in $Y_j$. This produces an error-$\varepsilon$ clocked block machine for $G_\ell$, with registers at original cuts $p_\ell+dm$, $0\le m\le M$. Theorem~\ref{thm:radius54} bounds its narrow cuts $1\le m\le M$ by a constant depending only on $k$ and the fixed data. The omitted cuts of residue $\ell$ number at most $p_\ell+b_\ell+1$: this also charges the block cut zero, while all counted cuts remain in $1,\ldots,N$. If $M<0$, the whole horizon is bounded by $p_\ell+b_\ell$. Sum these bounds over the $p$ residue classes. This counts all cuts $1,\ldots,N$, independently of all other register widths.
\end{proof}
The constants in this theorem depend on the fixed alphabet, return period and certificate, interface, error, width $k$, and chosen finite prefixes. They do not depend on the horizon or on the machine. The original constants $B,d,\chi,L_k$ in Section~\ref{sec:local54} use $d$ for a contraction defect. The return period here is denoted $p$ to keep the two parameters distinct.

\begin{proposition}[Different phases can have different thresholds]\label{prop:parity55}
There is an alphabet with return period two for which $W_{N,0}=1$ on every even horizon, but $W_{N,\varepsilon}\to\infty$ along odd horizons for every $0\le\varepsilon<\rho/2$.
\end{proposition}
\begin{proof}
Take $H=(\R/2\pi\Z)\times\Z/2\Z$ and letters $(\theta,1),(-\theta,1)$, with $\theta/2\pi$ irrational. The generated group is dense in $H$. An exact return must have zero signed count and hence even length; the two-letter return gives every positive even length. Here $J=(\R/2\pi\Z)\times\{0\}$. For one binary-mean query put $F(\phi,0)=0$ and $F(\phi,1)=\rho\cos\phi$, where $0<\rho\le1$. The even-phase target is identically zero and needs one label. The odd-phase radius in the norm $|\cdot|/2$ is $\rho/2$. Apply Theorem~\ref{thm:phases55}.
\end{proof}

\section{Infinite component groups and boundary attainment}\label{sec:profinite55}
This section retains independent finite-dimensional proofs of the strict-radius and exact translate-rank conclusions. Its output spaces are finite-dimensional. Let $H$ again be an arbitrary compact Hausdorff group, and set $K=H^\circ$. The quotient $H/K$ is profinite. We use two standard compact-group consequences of this fact and of Peter--Weyl theory \cite{Folland}: open normal subgroups form a neighborhood basis in a profinite group, and, for a continuous surjection $\pi:H\to G$ onto a compact Lie group, one has $\pi(K)=G^\circ$. To see the latter assertion, $\pi(K)$ is connected and normal, while $G/\pi(K)$ is a continuous group quotient of the profinite group $H/K$ and is therefore totally disconnected. This forces $G^\circ\subset\pi(K)$; the reverse inclusion is immediate. A continuous finite-dimensional representation of a profinite group has finite image, by the closed subgroup theorem and the no-small-subgroups property of a Lie group \cite{Hall}: an open normal subgroup can be placed inside a no-small-subgroups neighborhood, so its image is trivial.

Define the component radius
\begin{equation}\label{eq:R055}
                 R_0=\max_{s,j,g\in H}\rad_{Y_j}(F_{s,j}(gK)).
\end{equation}
Unlike $r(U)$, this quantity need not itself be realized by any finite quotient.

\begin{lemma}[Approximation by finite quotient radii]\label{lem:infradius55}
With $r(U)$ as in \eqref{eq:quotientradius55},
\[
                 \inf_{U\text{ open normal in }H}r(U)=R_0.
\]
\end{lemma}
\begin{proof}
Every open subgroup contains $K$, so $r(U)\ge R_0$. Fix $\eta>0$. Uniform continuity, for the finite family of interfaces, gives an identity neighborhood $V$ such that
\[
 \norm{F_{s,j}(xv)-F_{s,j}(x)}_j<\eta
       \quad(x\in H,\ v\in V,\ s,j).
\]
There is an open normal subgroup $U$ containing $K$ and contained in $KV$, by the profinite quotient basis. Write each $u\in U$ as $kv$ with $k\in K$, $v\in V$. Then $F_{s,j}(gu)$ lies within $\eta$ of $F_{s,j}(gk)$. A center for the latter component image consequently covers the former coset image with an additional error at most $\eta$. Thus $r(U)\le R_0+\eta$.
\end{proof}

\begin{lemma}[A legal finite-dimensional interface approximation]\label{lem:Lieapprox55}
For every $\eta>0$ there is a continuous homomorphism $\pi:H\to G\subset O(D)$ onto a compact Lie group and continuous maps $\widetilde F_{s,j}:G\to Y_j$ such that
\[
           \max_{s,j}\norm{F_{s,j}-\widetilde F_{s,j}\circ\pi}_\infty<\eta.
\]
Moreover, for every $g\in H$,
\[
 \left|\rad_{Y_j}(F_{s,j}(gK))-
 \rad_{Y_j}(\widetilde F_{s,j}(\pi(g)G^\circ))\right|<\eta.
\]
\end{lemma}
\begin{proof}
Approximate all coordinates of all $F_{s,j}$ by finitely many real representative functions. The direct sum of their orthogonal representations is $\pi$, and the approximants are continuous functions on its compact image $G$. Choose Euclidean inner products on the finitely many output spaces. Project each approximant to the closed convex set $Y_j$ by nearest-point projection. This projection is continuous and nonexpansive in the chosen Euclidean norm. Choose constants $a_j,b_j>0$ with $a_j\|v\|_j\le|v|_2\le b_j\|v\|_j$. An initial error less than $a_j\eta/b_j$ in $\|\cdot\|_j$ then gives a projected error less than $\eta$, by Euclidean nonexpansiveness. These constants are fixed before choosing the approximants. The closed subgroup $G\subset O(D)$ is a compact Lie group. Since $\pi(K)=G^\circ$, the two indicated images have Hausdorff distance less than $\eta$. Their constrained radii therefore differ by less than $\eta$.
\end{proof}

\begin{theorem}[Arbitrary compact groups on both strict sides]\label{thm:profinite55}
For any compact-group interface, $\varepsilon>R_0$ implies a finite stationary permutation realization. At $\varepsilon=R_0$, such a realization exists exactly when $r(U)\le R_0$ for some open normal $U$.

Suppose in addition that executable exact returns have all sufficiently large lengths. If $\varepsilon<R_0$, then for each $k$ there is a horizon-independent upper bound on the number of cuts of width at most $k$, with all intervening widths unrestricted. In particular $W_{N,\varepsilon}\to\infty$. No finite-component assumption is made in either assertion.
\end{theorem}
\begin{proof}
The stationary assertions follow from Lemma~\ref{lem:infradius55} and Theorem~\ref{thm:finitequotient55}. For the lower bound choose a component image of radius greater than $\varepsilon$ and then $\eta>0$ with $2\eta$ smaller than that strict difference. Apply Lemma~\ref{lem:Lieapprox55}. A machine correct for the original interface at error $\varepsilon$ is correct for the approximating interface at error at most $\varepsilon+\eta$. The chosen component radius of that interface is greater than $\varepsilon+\eta$. The projected commands still have cofinite exact returns and generate $G$ densely. Theorem~\ref{thm:radius54} therefore gives the desired occupation bound for the very same registers. This argument never assumes that $K$ is open or that executable products lying exactly in $K$ are dense.
\end{proof}
Theorem~\ref{thm:completeboundary56} now identifies this stationary boundary criterion with clocked boundedness, including $R_0>0$, and under the weaker approximate-return hypothesis. The zero-error boundary also admits the following independent translate-rank argument.

\begin{lemma}[Evaluation rank and translate dimension]\label{lem:evaluation56}
For a real function $f$ on a group, the supremum of the ranks of finite matrices $(f(y_lx_i))_{i,l}$ equals the dimension, possibly infinite, of the span of the functions $x\mapsto f(yx)$. These are left translates with the convention $(L_gf)(x)=f(g^{-1}x)$, after replacing $y$ by $g^{-1}$.
\end{lemma}
\begin{proof}
A finite-dimensional span bounds every evaluation rank. Conversely, if $f_1,\ldots,f_m$ are independent functions, their evaluation vectors $(f_1(x),\ldots,f_m(x))$ span $\R^m$: a proper span would have a nonzero annihilator, giving a linear dependence of the functions. Select $m$ evaluation vectors forming a basis. This gives an invertible evaluation matrix and proves the reverse inequality for every finite $m$.
\end{proof}

\begin{theorem}[The exact clocked boundary]\label{thm:zero55}
Assume cofinite executable exact returns. The following are equivalent:
\begin{enumerate}
\item $\sup_NW_{N,0}<\infty$;
\item all $F_{s,j}$ factor through one finite continuous quotient $H/U$, with $U$ open normal;
\item there is an exact stationary finite permutation realization.
\end{enumerate}
This includes groups with infinitely many connected components.
\end{theorem}
\begin{proof}
The implications from the second assertion to the third and from the third to the first are immediate. Suppose the first holds with bound $k$. Theorem~\ref{thm:profinite55} implies $R_0=0$, so every interface is constant on each $K$-coset.

Fix a real coordinate $f$ of one $F_{s,j}$. For any finite families of executable prefix products $x_i$ and suffix products $y_l$, use cofinite returns to pad all selected prefix words to one common length $A$ and all suffix words to one common length $B$. Choose an exact realization at horizon $A+B$ of width at most $k$. At its cut $A$ the matrix
\[
                            (f(y_lx_i))_{i,l}
\]
factors through the $k$ hidden labels: the left factors are the actual state distributions after the padded prefixes and the right factors are the actual conditional terminal expectations on the padded suffixes. Hence this matrix has rank at most $k$. Padding does not need to preserve hidden states; it only preserves the physical products used in the displayed entries. By density and continuity the same rank bound holds for arbitrary $x_i,y_l\in H$.

It follows that the span of the left translates of $f$ in $C(H,\R)$ has dimension at most $k$. Otherwise $k+1$ linearly independent translates would have an invertible evaluation matrix at some $k+1$ points, contrary to the preceding rank bound. Lemma~\ref{lem:evaluation56} formalizes this evaluation assertion. Equivalently it follows inductively: a nonzero function is nonzero somewhere, and a function outside the span of the preceding ones cannot have all its evaluations forced by theirs.

Take the sum $V$ of these translate spaces over all seed--query coordinates. It is finite-dimensional and invariant under left translation. Translation is continuous in the uniform norm and preserves that norm, so it gives a continuous representation of $H$ on $V$ with compact image; Haar averaging makes it orthogonal. Because $K$ is normal and every original coordinate is constant on $K$-cosets, $K$ acts trivially on every translate and thus on $V$. The representation factors through $H/K$. Its image is finite by the profinite finite-dimensional representation fact above. Its kernel $U$ is open normal. Since every original coordinate belongs to $V$, all interfaces are $U$-invariant and descend to $H/U$.
\end{proof}

\begin{proposition}[A nonattained zero-radius boundary]\label{prop:adic55}
On a compact group with infinitely many components, every positive error can have a bounded stationary realization while the zero-error clocked width is unbounded.
\end{proposition}
\begin{proof}
Let $H=\Z_2$ be the additive group of $2$-adic integers, with commands $0$ and $1$. Their positive products are dense, and the zero command supplies exact returns of every length. If $x=\sum_{n\ge0}x_n2^n$, $x_n\in\{0,1\}$, prescribe the binary success probability
\[
                         f(x)=\sum_{n\ge0}\frac{2x_n}{3^{n+1}}.
\]
Uniform convergence makes $f$ continuous. It is injective: at the first differing digit the contribution $2/3^{n+1}$ exceeds the sum $1/3^{n+1}$ of all possible later differences. Since $H$ is totally disconnected, $K=\{0\}$ and $R_0=0$. Truncation after $m$ digits factors through $\Z/2^m\Z$ and has uniform error at most $\sum_{n=m}^\infty2/3^{n+1}=3^{-m}$. Thus every positive tolerance admits a finite quotient machine. But the injective function $f$ cannot factor through any finite quotient. Theorem~\ref{thm:zero55} excludes bounded exact clocked width, and Theorem~\ref{thm:finitequotient55} excludes a finite exact stationary realization.
\end{proof}

\begin{theorem}[Linear exact width at a profinite boundary]\label{thm:adiclinear55}
For the interface of Proposition~\ref{prop:adic55},
\[
                         W_{N,0}=\Theta(N).
\]
More precisely, if $n$ is a power of two with $2n-1\le N$, then $W_{N,0}\ge n$, while $W_{N,0}\le N+1$. Every exact realization, with unrestricted intermediate widths, has at most $4k$ cuts in $\{1,\ldots,N\}$ of width at most $k$. Nevertheless, for each $\varepsilon>0$, $\sup_NW_{N,\varepsilon}<\infty$.
\end{theorem}
\begin{proof}
For nonnegative integers $z$ write $v_2(z+1)$ for the exponent of two in $z+1$. Incrementing the binary expansion, with $v_2(z+1)$ trailing ones, gives
\begin{equation}\label{eq:adicdifference55}
 b(z):=f(z+1)-f(z)=\frac53\,3^{-v_2(z+1)}-1.
\end{equation}
Let $n=2^m$, and form $B_n=(b(i+j))_{0\le i,j<n}$. Define a function on $\Z/n\Z$ by
\[
 \psi_m(x)=3^{-v_2(x)}\ (x\ne0),\qquad \psi_m(0)=3^{-m}.
\]
Here the valuation of a nonzero residue is independent of its representative. Since $1\le i+j+1\le2n-1$, its valuation equals that used in $\psi_m(i+j+1\bmod n)$, including the case $i+j+1=n$. Thus replacing row $i$ by the old row $(-i-1)\bmod n$ changes its entry at column $j$ to the kernel value at $j-i\bmod n$. This is the circulant with kernel $(5/3)\psi_m-1$.

On $\Z/n\Z$ one has the elementary identity
\[
            \psi_m=1-2\sum_{\ell=1}^m3^{-\ell}
                               \mathbf1_{2^\ell\mid x}.
\]
Use the eigenvector with entries $\exp(2\pi \mathrm i qj/n)$; its eigenvalue is the kernel sum against $\exp(2\pi \mathrm i qx/n)$. The sign is immaterial for subgroup indicators, whose transforms below are real. The Fourier eigenvalue at frequency zero is
\[
 \mu_0=\frac23 n6^{-m}>0.
\]
At frequency $q\ne0$ it is
\[
 \mu_q=-\frac{10n}{3}
           \sum_{\substack{1\le\ell\le m\\ n/2^\ell\mid q}}6^{-\ell}<0.
\]
Indeed a subgroup indicator has Fourier sum $n/2^\ell$ exactly when $n/2^\ell$ divides $q$, and zero otherwise. For $q\ne0$ the summand $\ell=m$ is always present. At $m=0$ the sole eigenvalue is $2/3$. Consequently $B_n$ is invertible for every dyadic $n$.

At an arbitrary command cut $t<N$, take a dyadic $n\le\min(t+1,N-t)$. For each $i<n$ choose a prefix of length $t$ with $i$ increment commands. For each $j<n$ choose two suffixes of length $N-t$ with respectively $j+1$ and $j$ increment commands. Fill all remaining positions with zero commands. The difference of their exact outputs is $b(i+j)$. Hence $B_n$ factors through the actual register $S_t$: one factor consists of the prefix state rows and the other of the differences of the two conditional suffix decoder columns. Invertibility gives $|S_t|\ge n$. The difference column need not be a legal decoder; it is used only for this rank bound on two legal experiments.

Taking $t=n-1$ proves $W_{N,0}\ge n$ whenever $2n-1\le N$. The largest dyadic $n\le(N+1)/2$ is greater than $(N+1)/4$, proving the linear lower order. Storing the number of increment commands uses at most $N+1$ labels and gives the exact upper bound.

For occupation, if $t<N$ and $|S_t|\le k$, the largest dyadic integer not exceeding $\min(t+1,N-t)$ is at most $k$, so $\min(t+1,N-t)<2k$. Such cuts lie within fewer than $2k$ positions of one of the two ends. Including the terminal cut gives at most $4k$ positive cuts. No width at another cut was restricted. The positive-error assertion is the finite-quotient truncation in Proposition~\ref{prop:adic55}.
\end{proof}
This supplies a sharp quantitative law on a genuinely infinite profinite group. Its linear exact-width obstruction and bounded positive-error realizations concern one fixed continuous interface; they are not a uniform approximation statement over all continuous outputs on that group.

\section{Binary, categorical, and observable interfaces}\label{sec:observable53}
\begin{corollary}[Component oscillation for binary means]\label{thm:classification53}
For a finite continuous binary-mean interface $f_{s,j}:H\to[-1,1]$, with the return hypothesis of Theorem~\ref{thm:radius54}, the exact bounded-width threshold in total variation is
\begin{equation}\label{eq:critical53}
 \epsc=\frac14\max_{s,j,c}\left(\max_{h\in c}f_{s,j}(h)-\min_{h\in c}f_{s,j}(h)\right).
\end{equation}
It is attained by a stationary permutation realization. Below it every fixed width has the occupation bound \eqref{eq:occupation54}.
\end{corollary}
\begin{proof}
In the norm $|\cdot|/2$, the radius of an interval is one quarter of its length. Its midpoint lies in $[-1,1]$. Apply Theorem~\ref{thm:radius54}.
\end{proof}
This recovers the scalar component theorem, including all of its nonuniform quantifiers. The vector theorem does not follow by applying this corollary separately to coordinates, because the coordinate centers need not belong to the output simplex.

\begin{proposition}[A three-outcome threshold]\label{prop:categorical54}
On the circle let $\omega=e^{2\pi i/3}$ and, for $j=0,1,2$, set
\[
 p_j(\phi)=\frac19\left|1+e^{i\phi}\omega^{-j}+e^{2i\phi}\omega^{-2j}\right|^2.
\]
For any dense circle command alphabet satisfying the return condition, the categorical law $p=(p_0,p_1,p_2)$ has critical total-variation error $2/3$. Each separately requested binary event $\{j\}$ has threshold $1/2$.
\end{proposition}
\begin{proof}
Orthogonality of the three characters gives $\sum_jp_j=1$; each entry is nonnegative. At $\phi=2\pi j/3$ the vector is the simplex vertex $e_j$. Hence every center $q\in\Delta_3$ has maximum distance at least
\[
 \max_j\TV(e_j,q)=1-\min_jq_j\ge2/3.
\]
The uniform center has distance at most $2/3$ from every point of $\Delta_3$, by convexity and the values at its vertices. Thus the constrained radius is exactly $2/3$. On the other hand each $p_j$ has range $[0,1]$, so its binary mean $2p_j-1$ has oscillation two. Apply Theorem~\ref{thm:radius54} and Corollary~\ref{thm:classification53}. In particular the categorical radius exceeds half the diameter, which here equals $1/2$.
\end{proof}

\begin{corollary}[Joint measure-once unitary laws]\label{cor:quantum54}
Let unitary commands have compact closure $H$ and satisfy the return condition. For a fixed unit seed $\psi_s$ and a finite measurement $E_{j,l}\ge0$, $\sum_lE_{j,l}=I$, prescribe the law
\[
 F_{s,j}(h)_l=\ip{h\psi_s}{E_{j,l}h\psi_s}.
\]
The critical error is its componentwise simplex radius in total variation. For closure $U(d)$ and one orthonormal-basis measurement, this radius is $1-1/d$.
\end{corollary}
\begin{proof}
The displayed map is a continuous simplex-valued polynomial in realified matrix entries. A closed unitary matrix group has finitely many components, so Theorem~\ref{thm:radius54} applies. For $U(d)$, every unit vector is reachable from the seed, and its squared coordinate magnitudes range over the whole simplex. The simplex vertices force radius at least $1-1/d$, attained by the uniform distribution.
\end{proof}
A query selects one complete measurement, not simultaneous answers to incompatible measurements. Several compatible finite outputs, with specified correlations, instead form one joint categorical law. Language recognition at a cutpoint is a different requirement.

For linear binary interfaces on a compact matrix group, write
\[
 f_{s,j}(h)=q_j^Thx_s\in[-1,1],\quad
 V_{\rm r}=\spanop\{hx_s:h\in H,s\in\mathcal S\},
\]
\[
 V_{\rm n}=\{x\in V_{\rm r}:q_j^Thx=0\text{ for all }h\in H,j\in\mathcal J\}.
\]
Both spaces are invariant. The observable space $V_{\rm o}=V_{\rm r}\cap V_{\rm n}^\perp$ is the orthogonal complement of $V_{\rm n}$ \emph{inside} $V_{\rm r}$, canonically realizing the observable quotient.

\begin{theorem}[Intrinsic zero-error criterion]\label{thm:observable53}
For these linear means, bounded exact clocked width is equivalent to finiteness of the generated group's image in $O(V_{\rm o})$, and to trivial action of $K$ on $V_{\rm o}$. An exact stationary permutation realization then exists.
\end{theorem}
\begin{proof}
Corollary~\ref{thm:classification53} identifies bounded exact width with constancy of all interface means on every component. Under that condition, for $k\in K$ and $g,h\in H$, normality gives
\[
 q_j^Th(k-I)gx_s=q_j^T(hkh^{-1})hgx_s-q_j^Thgx_s=0.
\]
Thus $(k-I)V_{\rm r}\subset V_{\rm n}$. On the invariant orthogonal complement $V_{\rm o}$ the difference also lies in $V_{\rm o}$, hence vanishes. The observable action factors through the finite component group. Conversely, if $K$ acts trivially on $V_{\rm o}$, the same difference is invisible and all means are componentwise constant. A finite image of the dense generated group has the same finite image closure from $H$, whose connected subgroup $K$ acts trivially. The zero-dimensional quotient is included.
\end{proof}

\begin{corollary}[The planar threshold]\label{cor:circle53}
Let an orthogonal planar alphabet contain $I$ and an irrational rotation, with any additional orthogonal letters. For signed coordinate seeds $\pm\rho e_1,\pm\rho e_2$, coordinate queries, and $0<\rho\le1$, the threshold is $\epsc=\rho/2$. At and above it one command label with a fair binary decoder suffices. Below it every fixed width has a finite occupation budget.
\end{corollary}
\begin{proof}
The identity component is $SO(2)$, and every specified mean has range $[-\rho,\rho]$ on each component of $H=SO(2)$ or $O(2)$. All midpoint decoders equal zero, so the seed/component register collapses to one label.
\end{proof}
For the noncommuting rational alphabet
\[
 I,\quad R=\begin{pmatrix}3/5&-4/5\\4/5&3/5\end{pmatrix},\quad
 F=\begin{pmatrix}1&0\\0&-1\end{pmatrix},
\]
one has $FRF=R^{-1}\ne R$. The rotation has infinite order: otherwise $(3+4i)/5$ and its inverse would be algebraic integers, whereas their rational sum $6/5$ is not an integer. This example has the sharp threshold above. We do not infer a badly approximable angle from rational matrix entries.

\begin{proposition}[Independent toral coordinates]\label{prop:torus53}
Suppose a component is identified with the full product torus $(\R/2\pi\Z)^m$, surjectively in independent coordinates, and a mean on it equals
\[
 a_0+\sum_{l=1}^m(a_l\cos\theta_l+b_l\sin\theta_l).
\]
Its contribution to \eqref{eq:critical53} is $\tfrac12\sum_l\sqrt{a_l^2+b_l^2}$.
\end{proposition}
\begin{proof}
Each summand has the symmetric range with amplitude $\sqrt{a_l^2+b_l^2}$, and all extrema are simultaneously attainable in the full product. Divide the resulting oscillation by four.
\end{proof}
On a proper subtorus, extrema must instead be taken subject to its relations. Ambient coordinates do not certify independence. Similarly, invisible directions do not create an obstruction: for $H=\{\diag(R_\theta,1)\}$, seed and query $e_3$ give constant mean one, whereas seed $\rho e_1$ and query $e_1$ have threshold $\rho/2$.

\begin{proposition}[Compact similarity]\label{prop:similarity53}
The classification applies to a finite alphabet in $GL(D,\R)$ whose generated group, including inverses, is uniformly bounded, with the executable return hypothesis retained. Conjugating to an invariant Euclidean metric changes neither output probabilities nor stochastic widths.
\end{proposition}
\begin{proof}
Uniform bounds on matrices and inverses make the closure a compact group. The matrix $Q=\int_Hh^Th\,d\mu(h)$ is positive definite and satisfies $A^TQA=Q$. Hence $Q^{1/2}AQ^{-1/2}$ is orthogonal. Transport the continuous output maps through this conjugacy; for a linear mean transport $x_s$ to $Q^{1/2}x_s$ and $q_j$ to $Q^{-1/2}q_j$. Exact identity products and their word lengths are preserved.
\end{proof}

\subsection{A curved categorical image with distinct component centers}
\begin{proposition}\label{prop:curved55}
Let $H=(\R/2\pi\Z)\times\Z/2\Z$, with an identity letter, an irrational rotation in the first factor, and a toggle of the second factor. Put
\[
 q_0=(\tfrac12,\tfrac13,\tfrac16),\qquad
 q_1=(\tfrac16,\tfrac13,\tfrac12),\qquad t=\tfrac1{12},
\]
and prescribe the three-outcome law
\[
 F(\theta,c)_l=q_{c,l}+t\cos(\theta-2\pi l/3),\qquad l=0,1,2.
\]
Each component image is a proper closed curve in the interior of the simplex, with unique constrained total-variation center $q_c$ and radius $t$. The union of both component images has radius $1/4$. Thus the sharp clocked threshold is $1/12$, and the stationary minimum at that threshold is exactly two labels.
\end{proposition}
\begin{proof}
The three cosines sum to zero and every coordinate is at least $1/12$, so $F$ is a legal law in the simplex interior. The first two independent trigonometric coordinates identify a nondegenerate ellipse in its affine plane. For a zero-sum vector $z\in\R^3$, one has $\frac12\norm{z}_1=\max_l|z_l|$. Each coordinate of the curve ranges over $[q_{c,l}-t,q_{c,l}+t]$. Hence for every center $y\in\Delta_3$,
\[
 \sup_\theta\TV(F(\theta,c),y)
       =t+\max_l|q_{c,l}-y_l|.
\]
This proves uniqueness of $q_c$ and radius $t$. For the union, the least possible value of $\max_{c,l}|q_{c,l}-y_l|$ is $1/6$, forced by the two endpoints in the first coordinate and attained by $y=(1/3,1/3,1/3)$. Its radius is therefore $t+1/6=1/4$. A one-label stationary machine has a constant decoder and cannot reach error $t$; two component labels attain it. The clocked threshold follows from Theorem~\ref{thm:radius54}.
\end{proof}

\begin{remark}
The total-variation diameter of a probability simplex is one, so half its diameter is $1/2$; the three-vertex example above has radius $2/3$. In the unitary corollary, $U(d)$ is connected. These facts explain respectively why coordinate midpoints do not solve the categorical problem and why the full-unitary component is the whole group.
\end{remark}

\section{Matching width laws throughout the subcritical interval}\label{sec:rates54}
We now distinguish the dimension of the command arithmetic from the dimension of the physical orbit. Let $s\ge1$ and let the alphabet contain the rotations $R_{2\pi\alpha_1},\ldots,R_{2\pi\alpha_s}$ and $I$. For the matching upper bound the alphabet is exactly these letters, optionally with the reflection $F=\diag(1,-1)$. Use signed coordinate seeds of radius $\rho$ and coordinate queries. Extra letters do not affect the lower bounds, but are not silently included in the upper construction.

For a real number $x$, write $\|x\|_{\R/\Z}$ for distance to the nearest integer. The dual Diophantine condition is
\begin{equation}\label{eq:DA54}
 \|n\cdot\alpha\|_{\R/\Z}\ge b\|n\|_1^{-\tau}
 \qquad(0\ne n\in\Z^s),\qquad b>0,\quad\tau\ge s.
\end{equation}
In particular $1,\alpha_1,\ldots,\alpha_s$ are rationally independent. When $\tau=s$ this is the dual badly approximable condition. All implicit constants below may depend on $\alpha,b,s,\rho,\varepsilon$, but not on $N$ or a stochastic realization.

\begin{theorem}[Full-subcritical matching law]\label{thm:matched54}
Suppose \eqref{eq:DA54} holds with $\tau=s$, and let $0<\rho<1$. For every fixed $0\le\varepsilon<\rho/2$,
\begin{equation}\label{eq:matched54}
                 W_{N,\varepsilon}=\Theta\bigl(N^{s/(2s+1)}\bigr).
\end{equation}
The upper bound is achieved by an exact horizon-dependent realization. At $\varepsilon\ge\rho/2$, one command label suffices. If $\tau\ge s$ is arbitrary, the lower bound $W_{N,\varepsilon}=\Omega(N^{s/(2\tau+1)})$ still holds throughout the same subcritical interval.
\end{theorem}
The exponent for one badly approximable rotation is $1/3$, not logarithmic. The statement includes zero error when the signal has strict amplitude slack $\rho<1$. The threshold theorem continues to include $\rho=1$, but the exact polynomial upper bound here is not asserted at that endpoint. The earlier arithmetic argument obtained this scale in a smaller noise range. The proof below supplies the missing full-noise lower estimate and states the exact upper realization in full.

\begin{proposition}[An explicit family in every command dimension]\label{prop:algebraicangles54}
Let $\xi$ be a real algebraic integer of degree $s+1$. Then
$\alpha=(\xi,\xi^2,\ldots,\xi^s)$ satisfies \eqref{eq:DA54} with $\tau=s$ and some effectively computable $b>0$.
\end{proposition}
\begin{proof}
For $0\ne n\in\Z^s$ put $H=\|n\|_1$, and choose $n_0\in\Z$ so that
$z=n_0+\sum_{l=1}^s n_l\xi^l$ has modulus $\|n\cdot\alpha\|_{\R/\Z}$.
The degree assumption gives $z\ne0$. Write $\xi_1=\xi,\ldots,\xi_{s+1}$ for the conjugates and put $M_j=\max_{1\le l\le s}|\xi_j|^l$. Then
$|n_0|\le H M_1+1/2$ and
\[
 \left|n_0+\sum_{l=1}^s n_l\xi_j^l\right|
 \le H(M_1+M_j+1/2)\qquad(j\ge2).
\]
The norm of the nonzero algebraic integer $z$ is a nonzero integer. Thus
\[
 |z|\ge H^{-s}\prod_{j=2}^{s+1}(M_1+M_j+1/2)^{-1}.
\]
Choose a positive rational $b$ below this computable positive constant and $1/2$. This proves the claim. It concerns algebraic \emph{angles}, not algebraic entries of their rotation matrices.
\end{proof}

\subsection{Harmonic localization}
For the harmonic calculation assume that the alphabet contains $I,R_\theta$ and that one seed-query pair has mean $\rho\cos\phi$ after accumulated angle $\phi$, where $0<\rho\le1$. Other seeds, queries, or commands impose additional constraints and do not invalidate the lower bound. Let $\lambda=e^{i\theta}$ have infinite order. Put $\delta=2\varepsilon<\rho$, and choose an integer
\begin{equation}\label{eq:r53}
 r\ge1,\qquad r>\frac{\delta}{\rho-\delta},\qquad m=r+1.
\end{equation}
For example $r=\lfloor\delta/(\rho-\delta)\rfloor+1$ works. Define
\begin{equation}\label{eq:circleconstants53}
 \Delta_r=2\left(\frac{\rho r}{r+1}-\delta\right)>0,\qquad
 A_r=\frac{2(1+\rho+\delta)(2r+1)m}{\Delta_r}>1.
\end{equation}

\begin{lemma}[Explicit harmonic localization]\label{lem:harmonic53}
The nonnegative trigonometric polynomials
\[
 p_\pm(\phi)=\frac{(1\pm\cos\phi)^r}{2^{-r}\binom{2r}{r}}
\]
have Haar mean one and weighted means
\[
 \int p_\pm(\phi)\rho\cos\phi\,\frac{d\phi}{2\pi}
                         =\pm\frac{\rho r}{r+1}.
\]
For the moving feature vector
\[
 Z_m(\phi)=(\cos\phi,\sin\phi,\ldots,\cos m\phi,\sin m\phi),
\]
any terminal realization with wordwise mean error at most $\delta$ obeys
\begin{equation}\label{eq:harmonicresidual53}
 \E\norm{\E[Z_m(\phi)\mid S]}\ge
 \frac{\Delta_r}{2(1+\rho+\delta)(2r+1)}.
\end{equation}
\end{lemma}
\begin{proof}
The identity $1+\cos\phi=2\cos^2(\phi/2)$, or the constant Laurent coefficient, gives
\[
 C_r:=\int(1+\cos\phi)^r\frac{d\phi}{2\pi}
            =2^{-r}\binom{2r}{r},\qquad
 C_{r+1}/C_r=\frac{2r+1}{r+1}.
\]
Hence the weighted cosine mean is $C_{r+1}/C_r-1=r/(r+1)$. Translation by $\pi$ gives the negative case.

We use the unweighted real coefficient vector in the basis $1,\cos\phi,\sin\phi,\ldots$, without orthonormal Fourier rescaling. The absolute sum of the real cosine coefficients of either $p_\pm$, including the constant coefficient, is
\[
 B_r=\frac{4^r}{\binom{2r}{r}}\le2r+1.
\]
For $p_+$ its coefficients are nonnegative and their sum is $p_+(0)$; for $p_-$ they differ only by alternating signs. The inequality follows because the central binomial coefficient is the largest of the $2r+1$ coefficients summing to $4^r$. Multiplication by $\rho\cos\phi$ has coefficient absolute sum at most $\rho B_r$, since $\cos u\cos v=(\cos(u+v)+\cos(u-v))/2$. Thus the Euclidean norms of the moving coefficient vectors of $p_\pm$ and $p_\pm\rho\cos\phi$ are at most $2r+1$ and $\rho(2r+1)$, respectively. For completeness, let $c=\E D$, $Q=\E\norm{\E[Z_m\mid S]}$. For either weight, the conditional coefficient identity gives $|c-\int p_\pm\rho\cos\phi|\le\delta+(1+\rho+\delta)(2r+1)Q$. Subtracting the two inequalities gives \eqref{eq:harmonicresidual53}. Correctness is averaged word by word before using this identity. The starting feature norm is $\sqrt m\le m$; $A_r$ deliberately uses the latter bound to obtain a rational closed form.
\end{proof}

The integer $r$ in \eqref{eq:r53} is the localization order; it is unrelated to the number $s$ of rotation letters. For fixed subcritical error it is finite and independent of $N$ and $k$.

\begin{theorem}[A full-noise arithmetic occupation bound]\label{thm:arithmetic54}
Assume \eqref{eq:DA54} and $0\le\varepsilon<\rho/2$. Fix the localization parameters $r,m,A_r$ in \eqref{eq:r53}--\eqref{eq:circleconstants53}. For an integer $k\ge1$, put
\[
 L=\left\lceil(2k)^{1/s}\right\rceil,\qquad B=s(L-1).
\]
Every error-$\varepsilon$ realization of peak at most $k$ satisfies
\begin{equation}\label{eq:finiteDA54}
 N<B\left(1+b^{-2}(mB)^{2\tau}\log A_r\right).
\end{equation}
With no restriction on other widths, the number of cuts of width at most $k$ is at most
\begin{equation}\label{eq:DAoccupation54}
 \left\lceil B-1+B b^{-2}(mB)^{2\tau}\log A_r\right\rceil.
\end{equation}
\end{theorem}
\begin{proof}
Use $s$ consecutive groups of $L-1$ command positions. In group $i$, choose $J_i$ uniformly from $\{0,\ldots,L-1\}$, apply $J_i$ copies of the $i$th rotation, and fill the other positions with identities. This defines an external law on length-$B$ words with $L^s\ge2k$ equally likely products. On the harmonic feature use
\[
 R_m(t)=\diag(R_t,R_{2t},\ldots,R_{mt}).
\]
For two distinct count vectors $J,J'$ and any frequency $1\le l\le m$, condition~\eqref{eq:DA54} and $\|l(J-J')\|_1\le mB$ give
\[
 |e^{2\pi i l(J-J')\cdot\alpha}-1|
 =2\bigl|\sin\bigl(\pi l(J-J')\cdot\alpha\bigr)\bigr|
 \ge4b(mB)^{-\tau}=:\sigma.
\]
Here $2\sin(\pi x)\ge4x$ for $0\le x\le1/2$. Also $b\le1/2$, $m\ge2$, and $B\ge1$, so $\sigma\le1$.

For a unit vector $u=(u_1,\ldots,u_m)$ in the frequency blocks, the squared distance between the two represented orbit points is
\[
 \sum_{l=1}^m\norm{u_l}^2
       |e^{2\pi i l(J-J')\cdot\alpha}-1|^2\ge\sigma^2.
\]
An open ball of radius $\sigma/2$ about any unit center contains at most one point. At least half the $L^s$ points therefore have distance at least $\sigma/2$ from all of any $k$ centers. Since the inner-product defect between unit vectors equals half their squared distance, the finite word law has defect at least
\[
                  d_k=\sigma^2/16=b^2(mB)^{-2\tau}.
\]
Apply conditional-centroid contraction to successive blocks and Lemma~\ref{lem:harmonic53} to the terminal query. The starting norm is $\sqrt m\le m$. Thus
\[
 \left\lfloor\frac NB\right\rfloor
 \le\frac{\log A_r}{-\log(1-d_k)}
 \le b^{-2}(mB)^{2\tau}\log A_r,
\]
which proves \eqref{eq:finiteDA54}. The greedy endpoint count, with identity gaps and no prefix, gives \eqref{eq:DAoccupation54}. Internal registers remain unrestricted. Finally $B<s(2k)^{1/s}$, so \eqref{eq:finiteDA54} implies the asserted $\Omega(N^{s/(2\tau+1)})$ lower bound.
\end{proof}

\begin{remark}
Testing a coordinate unit vector in \eqref{eq:DA54} gives $b\le1/2$. In the block construction, $L=\lceil(2k)^{1/s}\rceil\ge2$, hence $B=s(L-1)\ge1$. The integer $r$ is the localization order and $s$ is the number of rotations. For the algebraic-angle example, $1,\xi,\ldots,\xi^s$ are linearly independent over $\Q$ by the degree $s+1$ hypothesis; this ensures that the algebraic integer used in its norm bound is nonzero.
\end{remark}

\subsection{An exact polygon realization}
Write $\eta_\alpha(q)=\max_i\|q\alpha_i\|_{\R/\Z}$. The following elementary construction explains why a common denominator, rather than the number of numerical orbit points, controls the upper width.

\begin{proposition}[Exact enclosure upper bound]\label{prop:polygon54}
Let $0<\rho<1$, $a=(1+\rho)/2$, and $\gamma=\log(a/\rho)>0$. Choose an integer $q_0\ge4$ with $\cos(\pi/q_0)\ge a$. If $q\ge q_0$ and
\begin{equation}\label{eq:enclosure54}
                4\pi^2N\,\eta_\alpha(q)/q^2\le\gamma,
\end{equation}
then the indicated rotation alphabet, with or without $F$, has an exact horizon-$N$ realization with $q$ command labels.
\end{proposition}
\begin{proof}
Let $P_q$ be the regular polygon with vertices $v_j=(\cos(2\pi j/q),\sin(2\pi j/q))$, $0\le j<q$. For each $i$ choose a nearest integer $p_i$ to $q\alpha_i$ and put
\[
 t=\pi/q,\qquad h_i=2\pi|\alpha_i-p_i/q|\le t,\qquad
 \Lambda_i=\frac{\cos(t-h_i)}{\cos t}.
\]
Writing $n_j=(\cos((2j+1)t),\sin((2j+1)t))$, its facet description is
\[
               P_q=\{x:\ip{n_j}{x}\le\cos t\ (0\le j<q)\}.
\]
The facet inequalities of the regular polygon give $R_{2\pi\alpha_i}P_q\subset\Lambda_iP_q$: the nearest facet normal to a rotated vertex has angular difference $t-h_i$. Further,
\[
 \log\Lambda_i
 \le\log(1+\tan(t)h_i)\le\tan(t)h_i
 \le4\pi^2\eta_\alpha(q)/q^2,
\]
since $\cos h+\tan(t)\sin h\le1+\tan(t)h$ and $\tan t\le2t$ for $0\le t\le\pi/4$. Let $\Lambda=\max(1,\Lambda_1,\ldots,\Lambda_s)$. The standard reflection $F=\diag(1,-1)$ sends $v_j$ to $v_{-j\ ({\rm mod}\ q)}$ and preserves the chosen vertex set. The identity and reflection therefore preserve $P_q$, so every command $U$ satisfies $\Lambda^{-1}UP_q\subset P_q$.

For each command and each vertex choose a probability row with barycenter $\Lambda^{-1}Uv_j$ in $P_q$. These are actual common command rows, not independent prefix/suffix factorizations. At most three successors suffice by planar convexity. Initialize the seed $\rho u_s$, $u_s\in\{\pm e_1,\pm e_2\}$, with a row whose barycenter is $au_s$. This is possible since $a\le\cos(\pi/q)$, the polygon's inradius. After a length-$N$ word the mean vertex is
\[
                         a\Lambda^{-N}U_wu_s.
\]
By \eqref{eq:enclosure54}, $\rho\Lambda^N\le a$. Hence the terminal columns
\[
                       d_j(v)=\frac{\rho\Lambda^N}{a}\,e_j^Tv
\]
belong to $[-1,1]$ and return exactly $\rho e_j^TU_wu_s$. The rows can be stationary within a horizon; the terminal scaling and chosen denominator may depend on $N$. No one finite all-horizon exact machine is asserted.
\end{proof}

\begin{lemma}[Comparable simultaneous denominators]\label{lem:denominators54}
Assume \eqref{eq:DA54}. For every integer $Q\ge2$ there is $q$ with $1\le q\le Q^s$, $\eta_\alpha(q)\le1/Q$, and
\[
 q\ge c Q^{s/(\tau+1-s)},\qquad
 c=\left(\frac{b}{s^{\tau+1}}\right)^{s/(\tau+1-s)}.
\]
In particular, for $\tau=s$ one has $cQ^s\le q\le Q^s$.
\end{lemma}
\begin{proof}
Partition the torus into $Q^s$ half-open cubes and apply the pigeonhole principle to $0,\alpha,\ldots,Q^s\alpha$. Their difference gives $q\le Q^s$ and integers $p_i$ with $|q\alpha_i-p_i|\le1/Q$.

Put $M=\lfloor q^{1/s}\rfloor$. The $(M+1)^s>q$ vectors in $\{0,\ldots,M\}^s$ cannot all have distinct residues $n\cdot p$ modulo $q$, even when some nearest integers $p_i$ are negative. Subtract two to obtain a nonzero vector $n$ with $q\mid n\cdot p$ and $H:=\|n\|_1\le s q^{1/s}$. Therefore
\[
 bH^{-\tau}\le\|n\cdot\alpha\|_{\R/\Z}
 \le H/(qQ),\qquad
 q^{(\tau+1-s)/s}\ge bQ/s^{\tau+1}.
\]
The displayed bound follows. This argument supplies the needed quantitative transference directly; no unspecified denominator subsequence is being used.
\end{proof}

\begin{proof}[Proof of Theorem~\ref{thm:matched54}]
The lower estimate follows from Theorem~\ref{thm:arithmetic54}. For $\tau=s$, take the constants $a,\gamma,q_0,c$ above and an integer
\[
 Q\ge\max\left(2,(q_0/c)^{1/s},
                \left(\frac{4\pi^2N}{c^2\gamma}\right)^{1/(2s+1)}\right).
\]
Choosing the least such integer, Lemma~\ref{lem:denominators54} gives $q\le Q^s=O(N^{s/(2s+1)})$, $q\ge q_0$, and
\[
 4\pi^2N\eta_\alpha(q)/q^2
 \le4\pi^2N/(c^2Q^{2s+1})\le\gamma.
\]
Proposition~\ref{prop:polygon54} supplies the exact upper bound. At and above $\rho/2$ use the fair decoder from Corollary~\ref{cor:circle53}. This proves all claims.
\end{proof}
The command count $s$ is an arithmetic parameter, not a replacement for the dimension of the physical circle. The constants deteriorate as $\varepsilon\uparrow\rho/2$ through the localization degree, and as $\rho\uparrow1$ through the enclosure slack. Both dependences are explicit in the finite inequalities. The theorem makes no uniform claim at these limiting parameters.

\begin{corollary}[Unit amplitude at positive error]\label{cor:unitnoise55}
For the exact rotation alphabet of Theorem~\ref{thm:matched54}, with or without its specified reflection, suppose the angle vector is dual badly approximable. At $\rho=1$, for every fixed $0<\varepsilon<1/2$,
\[
                     W_{N,\varepsilon}=\Theta(N^{s/(2s+1)}).
\]
\end{corollary}
\begin{proof}
The harmonic localization and orbit-packing lower bound in Theorem~\ref{thm:arithmetic54} permits $\rho=1$ and every fixed $\varepsilon<1/2$. For the upper bound choose $0<\eta<2\varepsilon$ and put $\rho'=1-\eta\in(0,1)$, decreasing $\eta$ if necessary. Apply the exact polygon construction and comparable-denominator lemma to $\rho'$. On every seed, word, and query its mean differs from the unit-amplitude target by at most $\eta$, so the binary total-variation error is at most $\eta/2<\varepsilon$. Its width has the required order. This is a positive-error upper bound, not an exact realization at unit amplitude.
\end{proof}
The case $\rho=1$, $\varepsilon=0$ is not decided by this corollary. In particular, taking $\eta=0$ would destroy the strictly positive enclosure slack required by Proposition~\ref{prop:polygon54}.

\section{Spherical experiments with a spectral gap}\label{sec:nonabelian56}
The next result concerns a full vector orbit, not an abelian factor of a nonabelian system. Write $\sigma$ for normalized area on $S^{d-1}$, where $d\ge2$. Let a finite orthogonal alphabet contain the identity, and let $\mu$ be a probability law on its letters. Assume that its averaging operator
\[
                  Pf(x)=\E_{R\sim\mu}f(Rx)
\]
satisfies
\begin{equation}\label{eq:spectralgap56}
        \|P^n f\|_{L^2(\sigma)}\le\lambda^n\|f\|_{L^2(\sigma)}
        \quad\left(\int f\,d\sigma=0\right),\qquad 0<\lambda<1.
\end{equation}
The physical target is $F(g)=\rho g x_0$ for a fixed $x_0\in S^{d-1}$, with output set the closed Euclidean unit ball and error measured in Euclidean norm.

\begin{theorem}[A spherical width law up to a logarithm]\label{thm:spherical56}
Assume \eqref{eq:spectralgap56}, $0<\rho<1$, and $0\le\varepsilon<\rho$. There are constants $c,C>0$ such that
\begin{equation}\label{eq:sphericalrate56}
 c\left(\frac{N}{\log(N+2)}\right)^{(d-1)/2}
       \le W_{N,\varepsilon}\le C(N+1)^{(d-1)/2}.
\end{equation}
The upper realization is exact. Every error-$\varepsilon$ realization has at most
\[
                      C k^{2/(d-1)}\log(k+2)
\]
positive cuts of width at most $k$, with other widths unrestricted. The constants may depend on $d,\lambda,\rho,\varepsilon$, but not on the horizon or the realization. At unit amplitude the same bounds hold for every fixed $0<\varepsilon<1$; exact unit amplitude is not included in the upper assertion.
\end{theorem}

\begin{lemma}[Spherical quantization and executable contraction]\label{lem:sphericalgap56}
Under \eqref{eq:spectralgap56}, there are $c_d>0$ and integers $B_k\le C_{d,\lambda}\log(k+2)$ such that
\[
 \E_{w\sim\mu^{B_k}}\max_{i\le k}\langle v_i,u_w u\rangle
                  \le1-c_d k^{-2/(d-1)}
\]
for all unit vectors $u,v_1,\ldots,v_k$. The defect may be decreased so that it lies in $(0,1)$. When $d=3$ one may take defect $1/(2k)$ and
\begin{equation}\label{eq:sphere3constants56}
                  B_k=\left\lceil\frac{\log(64k^2)}{-\log\lambda}\right\rceil.
\end{equation}
\end{lemma}
\begin{proof}
Normalized spherical cap measure obeys
\[
       c_d h^{d-1}\le\sigma\{x:\|x-u\|_2<h\}\le C_d h^{d-1}
                     \qquad(0<h\le1),
\]
uniformly in $u$. This follows by writing area in polar angle as a constant times $\sin^{d-2}\theta\,d\theta$ and using $h=2\sin(\theta/2)$; for $d=2$ the same formula is the arc-length formula. For $f(x)=\max_i\langle v_i,x\rangle$, choose $h=c k^{-1/(d-1)}\le1$ so that the union of the $k$ radius-$h$ caps has area at most $1/2$. Outside them, $1-f(x)\ge h^2/2$. Thus
\begin{equation}\label{eq:haarsphere56}
                         \int f\,d\sigma\le1-a_d k^{-2/(d-1)}.
\end{equation}
The function $f$ is $1$-Lipschitz for chordal distance and bounded by one in absolute value.

Let $K_\delta$ average over the chordal cap of radius $\delta\le1$. It commutes with all rotations, fixes constants, and satisfies $\|K_\delta f-f\|_\infty\le\delta$. The cap density at a center has $L^2$ norm at most $C_d\delta^{-(d-1)/2}$. Hence \eqref{eq:spectralgap56} and Cauchy--Schwarz imply, uniformly in the center,
\[
 \left|P^B f(u)-\int f\,d\sigma\right|
 \le\delta+C'_d\delta^{-(d-1)/2}\lambda^B.
\]
Take $\delta$ to be a fixed sufficiently small multiple of $k^{-2/(d-1)}$ and then $B=O_{d,\lambda}(\log(k+2))$ so that the right side is at most half the defect in \eqref{eq:haarsphere56}. The law of the $B$ successive letters is a finite executable word law, proving the general assertion.

For $d=3$, a chordal cap of radius $h\le2$ has normalized area $h^2/4$. The union bound gives $\Pr(1-f<t)\le kt/2$, so integration for $0\le t\le2/k$ gives $\int(1-f)\,d\sigma\ge1/k$. A radius-$\delta$ cap density has $L^2$ norm $2/\delta$, and $\|f-\int f\|_2\le2$. Thus the mixing error is at most $\delta+4\lambda^B/\delta$. Choose $\delta=1/(4k)$ and \eqref{eq:sphere3constants56}; the error is at most $1/(2k)$.
\end{proof}

\begin{proof}[Proof of Theorem~\ref{thm:spherical56}]
Let $Z=u_wx_0$ and let $D(S)$ be the terminal decoder. For every deterministic word, correctness and $\|Z\|_2=1$ give
\[
           \langle\E[D(S)\mid w],Z\rangle\ge\rho-\varepsilon.
\]
Under any external word law this implies
\begin{equation}\label{eq:sphericalresidual56}
       \rho-\varepsilon\le\E\langle D(S),Z\rangle
             \le\E\|\E[Z\mid S]\|_2,
\end{equation}
since the decoder lies in the unit ball. Lemma~\ref{lem:centroid54}, applied with the executable law of Lemma~\ref{lem:sphericalgap56}, contracts the last conditional norm by $1-a_k$ at each selected narrow endpoint, where $a_k=c_d k^{-2/(d-1)}\in(0,1)$. Initially the norm is one. Fixed identity-letter gaps cannot increase it, even though their stochastic rows may process the register. After $J$ selected endpoints, \eqref{eq:sphericalresidual56} gives
\[
            J[-\log(1-a_k)]\le\log\frac1{\rho-\varepsilon}.
\]
Greedy spacing by $B_k$ as in the occupation proof discards at most $B_k-1$ initial cuts and charges at most $B_k$ eligible cuts to each selected endpoint. Therefore the occupation count is at most
\begin{equation}\label{eq:sphericaloccupation56}
 B_k\left(1+a_k^{-1}\log\frac1{\rho-\varepsilon}\right).
\end{equation}
If every register has width at most $k$, this bounds $N$. Inverting the bound gives the lower order in \eqref{eq:sphericalrate56}. Explicitly, when $k\le(N+1)^{d-1}$ its logarithm is $O_d(\log(N+2))$, and otherwise the asserted lower bound is automatic.

For the upper bound, choose a maximal $\delta$-separated set of vertices on $S^{d-1}$ containing $x_0$, with $0<\delta\le1/2$. The disjoint caps of radius $\delta/2$ show that it has at most $C_d\delta^{-(d-1)}$ vertices. Maximality makes it a $\delta$-net. Its convex hull $Q$ satisfies
\begin{equation}\label{eq:sphericalpolytope56}
                r\mathbb B^d\subset Q\subset\mathbb B^d,
                         \qquad r=1-\delta^2/2,
\end{equation}
because for every unit direction $u$ some vertex $v$ has $\langle u,v\rangle\ge1-\delta^2/2$. The inclusion follows from the supporting half-space characterization of a closed convex body. For each command $R$ and vertex $v_i$, the point $rRv_i$ belongs to $Q$. Choose barycentric coordinates of that point as the $i$th row of one common stochastic command matrix. The identity command is treated in the same way, with the same factor $r$.

Initialize at the vertex $x_0$ and decode vertex $v_i$ as $(\rho/r^N)v_i$. Induction on the row means gives output exactly $\rho u_wx_0$ on every length-$N$ word. The decoder is legal if $r^N\ge\rho$. Choose
\[
                 \delta^2=\min\left(\tfrac14,
                                      \frac{-\log\rho}{N+1}\right).
\]
Since $\log(1-\delta^2/2)\ge-\delta^2$, the required inequality holds. The number of vertices is $O_{d,\rho}((N+1)^{(d-1)/2})$. Both rows and decoder may depend on the horizon, as permitted by the clocked model.

For $\rho=1$ and $\varepsilon>0$, the same lower proof applies. An exact upper realization at amplitude $\rho'=1-\varepsilon/2$ differs from the unit-amplitude target by $\varepsilon/2$ in Euclidean norm, proving the stated upper bound. There is no assertion of exactness at unit amplitude.
\end{proof}

\begin{corollary}[A rational noncommuting alphabet on $\mathrm{SO}(3)$]\label{cor:SO356}
Let the alphabet consist of the identity, the matrices
\[
 R_x=\begin{pmatrix}1&0&0\\0&3/5&-4/5\\0&4/5&3/5\end{pmatrix},
 \qquad
 R_z=\begin{pmatrix}3/5&-4/5&0\\4/5&3/5&0\\0&0&1\end{pmatrix},
\]
and their inverses. For $F(g)=\rho g e_3$, with $0<\rho<1$ and $0\le\varepsilon<\rho$, the bounds are
\[
             c\frac{N}{\log(N+2)}\le W_{N,\varepsilon}\le C(N+1).
\]
At error $\varepsilon\ge\rho$ a single zero decoder suffices. Every subcritical realization has at most $Ck\log(k+2)$ positive cuts of width at most $k$.
\end{corollary}
\begin{proof}
The angle $\theta$ of both rotations has $2\cos\theta=6/5$. It is not a rational multiple of $2\pi$: otherwise $2\cos\theta$ would be a rational algebraic integer, hence an integer. The closures of the two cyclic groups are the full circle subgroups about their distinct axes. Their Lie algebras and bracket span $\mathfrak{so}(3)$, so the generated compact closure is $\mathrm{SO}(3)$. The matrices do not commute.

Under the double cover $\mathrm{SU}(2)\to\mathrm{SO}(3)$, their unit-quaternion lifts are $(2+i)/\sqrt5$ and $(2+k)/\sqrt5$. They have algebraic matrix entries and generate a dense subgroup of $\mathrm{SU}(2)$: its closed image is all of $\mathrm{SO}(3)$, so its Lie algebra has dimension three. The spectral-gap theorem of Bourgain--Gamburd \cite{BG2012} applies. Adding the identity with positive mass makes the symmetric averaging operator have norm strictly less than one on mean-zero $L^2$. The spherical representation is a subrepresentation of the regular representation, so the uniform law on the five displayed letters satisfies \eqref{eq:spectralgap56} on $S^2$. Apply Theorem~\ref{thm:spherical56}. The radius of the full sphere $\rho S^2$ inside the unit ball is $\rho$, attained at zero.
\end{proof}
The spectral gap in this corollary is an imported theorem, not a numerical assertion extracted from finite orbit checks. The logarithmic gap between the lower and upper bounds is explicit. No matching pure linear law or explicit numerical spectral-gap constant is claimed.

\section{A uniform orbit criterion for quantitative width}\label{sec:uniform57}
The spherical estimate has an extension in which the dimension of the ambient feature space is not the relevant exponent. Two geometric quantities enter: a uniform small-ball exponent for every possible conditional-centroid direction, and the covering dimension of the particular target orbit. They need not agree. Keeping this distinction is essential for matrix-valued experiments.

Let $G$ be a compact connected Lie group with normalized Haar measure $m_G$ and a bi-invariant Riemannian distance. Let $R:G\to O(E)$ be a continuous orthogonal representation on a finite-dimensional real Euclidean space. Fix $z_0\in E$ of norm one, put
\[
 \mathcal O=R(G)z_0,\qquad Q=\operatorname{conv}(\mathcal O),
 \qquad q=\dim\mathcal O,
\]
and assume that $\mathcal O$ spans $E$ and that $E$ has no nonzero fixed vector. Thus $q\ge1$. A finite alphabet in $G$ contains the identity and carries a probability law $\mu$ for which
\begin{equation}\label{eq:groupgap57}
 \|P^nf\|_2\le\lambda^n\|f\|_2\quad\left(\int_Gf\,dm_G=0\right),
 \qquad Pf(g)=\E_{a\sim\mu}f(u_ag),\quad 0<\lambda<1.
\end{equation}
Letters outside the support of $\mu$ are allowed. All command products and stochastic rows retain the conventions of Definition~\ref{def:machine54}.

The uniform geometric hypothesis is that, for some $p>0$, $A<\infty$ and $h_0>0$,
\begin{equation}\label{eq:smallball57}
 \sup_{u,v\in S(E)}m_G\{g:\|R(g)u-v\|<h\}\le Ah^p
                      \qquad(0<h\le h_0).
\end{equation}
It is imposed on all unit vectors $u$, not only on $\mathcal O$. A conditional mean of points in $\mathcal O$, after normalization, need not belong to $\mathcal O$.

\begin{lemma}[From orbit small balls to executable contraction]\label{lem:orbitgap57}
Under \eqref{eq:groupgap57}--\eqref{eq:smallball57}, constants $c,C>0$ and integers $B_k\le C\log(k+2)$ can be chosen so that
\begin{equation}\label{eq:orbitgap57}
 \E_{w\sim\mu^{B_k}}\max_{1\le i\le k}
       \langle v_i,R(u_w)u\rangle\le1-c k^{-2/p}
\end{equation}
for every $u,v_1,\ldots,v_k\in S(E)$. The constant $c$ can be decreased so that the defect lies in $(0,1)$ for every $k\ge1$.
\end{lemma}
\begin{proof}
Choose $h=c_0k^{-1/p}\le\min(h_0,1)$ with $A c_0^p\le1/2$. Outside the union of the $k$ open balls in \eqref{eq:smallball57},
\[
 1-\max_i\langle v_i,R(g)u\rangle\ge h^2/2.
\]
The union has measure at most $1/2$. Consequently the Haar mean of
$f(g)=\max_i\langle v_i,R(g)u\rangle$ is at most $1-a_k$, where
$a_k=c_0^2k^{-2/p}/4$. Repeated centers are allowed. All such functions are bounded by one in absolute value and have a common Lipschitz constant $L$: differentiating the finite-dimensional representation bounds the derivative uniformly on unit vectors.

Write $D=\dim G$. In exponential coordinates at the identity, Haar measure has a smooth positive density. Hence a metric ball of radius $\delta\le\delta_0$ has measure at least $c_G\delta^D$. Average $P^Bf$ over this ball. Left translations are isometries, so $P^Bf$ has Lipschitz constant at most $L$. Cauchy--Schwarz and \eqref{eq:groupgap57} give
\[
 |P^Bf(1)-\bar f|
 \le L\delta+2c_G^{-1/2}\delta^{-D/2}\lambda^B.
\]
Take $\delta$ to be a sufficiently small fixed multiple of $a_k$, also bounded by $\delta_0$. Then choose $B_k\ge1$ to make the second term at most $a_k/4$. The first term is at most $a_k/4$ as well. Since $\log(a_k^{-1})=O(\log(k+2))$, this requires $B_k\le C\log(k+2)$. Products under $\mu^{B_k}$ are genuine words of one common length. Their reverse chronological order has the same law because the letters are independent with the same law. Thus $P^{B_k}f(1)$ is the expectation in \eqref{eq:orbitgap57}. This proves the lemma with defect $a_k/2$.
\end{proof}

\begin{lemma}[Quadratic inner approximation of an orbit hull]\label{lem:orbithull57}
There are constants $C_1,C_2,\delta_0>0$ such that, for $0<\delta\le\delta_0$, points $z_1,\ldots,z_K\in\mathcal O$, including $z_0$, can be chosen with
\begin{equation}\label{eq:orbithull57}
 K\le C_1\delta^{-q},\qquad
 (1-C_2\delta^2)Q\subseteq\operatorname{conv}\{z_1,\ldots,z_K\}\subseteq Q,
 \qquad 0<1-C_2\delta^2<1.
\end{equation}
\end{lemma}
\begin{proof}
The orbit is a compact smooth embedded homogeneous manifold. Finitely many coordinate charts with bounded derivatives and inverse derivatives give an intrinsic $\delta$-net of size at most $C\delta^{-q}$; adjoining $z_0$ changes only $C$. Use the induced Riemannian metric on the orbit. If the functional $x\mapsto\langle b,x\rangle$ attains its maximum at $x_b\in\mathcal O$, its derivative along the orbit vanishes there. The second derivative along unit-speed orbit geodesics has absolute value at most $C\|b\|$, uniformly in their starting points and directions, by compactness of the unit tangent bundle and smoothness of the embedding. A net point within intrinsic distance $\delta$ of $x_b$ therefore satisfies
\[
 \langle b,z_i\rangle\ge h_Q(b)-C\|b\|\delta^2,
 \qquad h_Q(b)=\max_{x\in Q}\langle b,x\rangle.
\]
Haar averaging gives $\int R(g)z_0\,dm_G=0$. Moreover $0$ lies in the interior of $Q$ in $E$. Otherwise a nonzero supporting functional would be nonnegative on the orbit and have Haar mean zero; continuity and full Haar support would make it vanish on the orbit, contrary to its spanning $E$. Thus $r\mathbb B_E\subset Q$ for some $r>0$, and $h_Q(b)\ge r\|b\|$. The support inequality becomes
$\max_i\langle b,z_i\rangle\ge(1-C\delta^2/r)h_Q(b)$.
The support-function criterion for inclusion of compact convex sets proves \eqref{eq:orbithull57}. Decrease $\delta_0$ to obtain the last bounds.
\end{proof}

\begin{theorem}[Geometric width and occupation]\label{thm:uniformorbit57}
Assume \eqref{eq:groupgap57}--\eqref{eq:smallball57}. Take one seed and one query with target $F(g)=\rho R(g)z_0$, output set $Q$, and Euclidean error. For $0<\rho\le1$ and $0\le\varepsilon<\rho$, put $\zeta=\rho-\varepsilon\in(0,1]$. Every error-$\varepsilon$ realization satisfies
\begin{equation}\label{eq:orbitoccupation57}
 \#\{1\le t\le N:|S_t|\le k\}
 \le C\log(k+2)\left(1+k^{2/p}\log\frac1\zeta\right),
\end{equation}
where $C$ depends on the geometric and spectral-gap data but not on $\rho,\varepsilon,N,k$ or the realization. In particular, at fixed parameters,
\begin{equation}\label{eq:orbitrates57}
 W_{N,\varepsilon}\ge c\left(\frac{N}{\log(N+2)}\right)^{p/2}.
\end{equation}
For $0<\rho<1$ there is an exact common-row realization with
\begin{equation}\label{eq:orbitupper57}
 W_{N,0}\le C\left(1+\frac{N+1}{-\log\rho}\right)^{q/2}.
\end{equation}
For $\rho=1$ and fixed $0<\varepsilon<1$, the upper order is $O_\varepsilon((N+1)^{q/2})$. When $p=q$, the lower and upper exponents agree, with the displayed logarithmic loss in the lower bound. The constrained error threshold is exactly $\rho$, attained by the constant zero decoder.
\end{theorem}
\begin{proof}
Put $Z=R(u_w)z_0$ and let $D(S)\in Q\subseteq\mathbb B_E$ be the decoder. For each deterministic word, the error bound gives
$\langle\E[D(S)\mid w],Z\rangle\ge\rho-\varepsilon$.
Under any external word law, conditioning on the register yields
\[
 \zeta\le\E\langle D(S),Z\rangle
          \le\E\|\E[Z\mid S]\|.
\]
The conditional-centroid Lemma~\ref{lem:centroid54} applies to the law of Lemma~\ref{lem:orbitgap57}. The uniformity over every direction in \eqref{eq:smallball57} is precisely what justifies this application. It contracts the conditional norm by $1-a$ at a selected endpoint of width at most $k$, where $a=c k^{-2/p}\in(0,1)$. The initial norm is one. Fixed identity-letter fillers cannot increase it, regardless of their rows or intervening widths. After $J$ selected endpoints,
\[
 J[-\log(1-a)]\le\log(1/\zeta).
\]
Greedily select narrow cuts at separation at least $B_k$, omitting the first $B_k-1$ cuts. Each selected endpoint accounts for at most $B_k$ cuts. The count is at most $B_k(1+a^{-1}\log(1/\zeta))$, proving \eqref{eq:orbitoccupation57}, also when $\zeta=1$. In that case even one selected endpoint is impossible. Inverting the fixed-parameter count as in Theorem~\ref{thm:spherical56} gives \eqref{eq:orbitrates57}.

For the upper bound take the net in Lemma~\ref{lem:orbithull57}, and write $r=1-C_2\delta^2$. The orbit hull is invariant under every command. Hence $rR(u_a)z_i$ is in the convex hull of the net points. Choose its barycentric coordinates as row $i$ of a stochastic matrix $T_a$, the same at every epoch within this horizon. Initialize at $z_0$. The conditional row means then evolve as $r^tR(u_w)z_0$. Decode label $i$ by $(\rho/r^N)z_i$, which belongs to $Q$ if $r^N\ge\rho$, since $0\in Q$. Choosing
\[
 \delta^2=\min\left(\delta_0^2,\frac1{2C_2},
                          \frac{-\log\rho}{2C_2(N+1)}\right)
\]
gives this inequality by $\log(1-x)\ge-2x$ for $0\le x\le1/2$. The net size gives \eqref{eq:orbitupper57}. At unit amplitude and positive error use the exact construction for amplitude $1-\varepsilon/2$; its discrepancy from the desired target is $\varepsilon/2$.

Finally, Haar mean zero and $\|R(g)z_0\|=1$ imply
\[
 \int_G\|\rho R(g)z_0-y\|^2\,dm_G=\rho^2+\|y\|^2.
\]
Every constrained center therefore has worst error at least $\rho$, while $y=0$ attains it. Below it \eqref{eq:orbitrates57} rules out bounded width. None of the constructions requires transition matrices to satisfy the physical group relations, and extra finitely many commands do not change the upper construction.
\end{proof}
The geometric ingredients in Lemma~\ref{lem:orbithull57} are the usual local geometry of compact homogeneous spaces; related covering estimates are developed in \cite{Szarek57}. The additional realization statement is the common stochastic-row construction and its combination with a contraction uniform over conditional-centroid directions. A small-ball estimate restricted to the target orbit alone does not imply Theorem~\ref{thm:uniformorbit57}.

\section{Positive matrix realizations in every dimension}\label{sec:matrix57}
Let $d\ge2$ and let $\mathcal H_d^0$ be the real Euclidean space of traceless Hermitian $d\times d$ matrices, with scalar product $\langle A,B\rangle=\operatorname{tr}(AB)$ and Frobenius norm $\|\cdot\|_{\rm F}$. Let
\[
 \mathcal D_d=\{D=D^*\succeq0:\operatorname{tr}D=1\},\quad
 c_d=I/d,\quad r_d=\sqrt{1-1/d}.
\]
This section concerns numerical outputs in $\mathcal D_d$. The requirement that every decoder value be positive with trace one is part of the model; it is not imposed merely on the final average.

\begin{lemma}[Small balls on the complex Grassmannian]\label{lem:grassmann57}
Let $P$ be a fixed rank-$m$ orthogonal projection in $\C^d$, with $1\le m<d$. For Haar $U\in\mathrm{SU}(d)$ there is $C_{d,m}<\infty$ such that
\[
 \Pr\{\|UPU^*-Q\|_{\rm F}<h\}\le C_{d,m}h^{2m(d-m)}\quad(0<h\le1)
\]
for every rank-$m$ projection $Q$. For rank one, writing
$\operatorname{dist}_{\rm pr}(P,Q)=\sqrt{1-\operatorname{tr}(PQ)}$, one has exactly
\begin{equation}\label{eq:projectivecap57}
 \Pr\{\operatorname{dist}_{\rm pr}(UPU^*,Q)<h\}=h^{2(d-1)}
                         \qquad(0\le h\le1).
\end{equation}
\end{lemma}
\begin{proof}
Conjugation is transitive, so fix $Q=\operatorname{diag}(I_m,0)$. A projection sufficiently close to $Q$ has range the graph of a unique matrix $Z\in\C^{(d-m)\times m}$, with
\[
 P(Z)=\begin{pmatrix}I\\Z\end{pmatrix}(I+Z^*Z)^{-1}(I\quad Z^*).
\]
This notation means the product of the three displayed matrices. The map and its inverse are smooth near zero, and its derivative at zero is $Z\mapsto\left(\begin{smallmatrix}0&Z^*\\Z&0\end{smallmatrix}\right)$, of Frobenius norm $\sqrt2\|Z\|_{\rm F}$. The invariant probability measure is normalized Riemannian volume on this compact homogeneous manifold; in this chart its density is smooth and bounded on a sufficiently small closed neighborhood. The inverse chart puts a Frobenius $h$-ball inside a Euclidean $Ch$-ball in real dimension $2m(d-m)$. This proves the bound for small $h$; increasing the constant covers the remaining interval. Transitivity makes the constant independent of $Q$.

For rank one choose $Q=e_1e_1^*$. A uniform unit vector can be obtained by normalizing independent standard complex Gaussian coordinates. If $X$ is the squared modulus of the first coordinate before normalization and $Y$ is the sum of the other squared moduli, then $X$ is exponential of mean one and $Y$ is an independent sum of $d-1$ such variables. For $0<z<1$,
\[
 \Pr\left\{\frac{X}{X+Y}>1-z\right\}
 =\E\exp\left(-\frac{1-z}{z}Y\right)=z^{d-1}.
\]
Put $z=h^2$. The endpoint cases follow by continuity. This proves \eqref{eq:projectivecap57}, and also shows that $\operatorname{dist}_{\rm pr}$ is the Frobenius chordal distance divided by $\sqrt2$.
\end{proof}

\begin{lemma}[A uniform conjugacy-orbit exponent]\label{lem:conjugacysmall57}
There is $C_d<\infty$ such that
\begin{equation}\label{eq:conjugacycap57}
 \sup_{\substack{A,B\in\mathcal H_d^0\\\|A\|_{\rm F}=\|B\|_{\rm F}=1}}
 m_G\{U:\|UAU^*-B\|_{\rm F}<h\}
                 \le C_d h^{2(d-1)}\qquad(0<h\le1).
\end{equation}
The exponent cannot be increased uniformly over $A$ and $B$.
\end{lemma}
\begin{proof}
Write the eigenvalues of $A$ as $a_1\ge\cdots\ge a_d$. Since their sum is zero and their squared sum is one, $a_1-a_d\ge d^{-1/2}$. Some adjacent gap therefore satisfies
\begin{equation}\label{eq:separatedcut57}
 a_m-a_{m+1}\ge\gamma_d:=\frac1{\sqrt d(d-1)}.
\end{equation}
Let $P_A$ be the spectral projection onto the first $m$ eigenvectors. It is well defined despite possible multiplicities within either block.

Two conjugates $A',A''$ of $A$, with eigenbases $x_i,y_j$ in this order, obey the exact identity
\begin{equation}\label{eq:spectralidentity57}
 \|A'-A''\|_{\rm F}^2
       =\sum_{i,j}(a_i-a_j)^2|\langle x_i,y_j\rangle|^2.
\end{equation}
Expanding both traces proves it. The cross-block terms and \eqref{eq:separatedcut57} imply
\begin{equation}\label{eq:projectorbound57}
 \|P_{A'}-P_{A''}\|_{\rm F}
                    \le\gamma_d^{-1}\|A'-A''\|_{\rm F},
\end{equation}
since the squared norm on the left is exactly the sum of the two cross-block overlap sums.

If the set in \eqref{eq:conjugacycap57} is nonempty, choose one member $U_0$. Every other member has $\|UAU^*-U_0AU_0^*\|_{\rm F}<2h$. Thus its rank-$m$ spectral projection lies within $2h/\gamma_d$ of $U_0P_AU_0^*$. The Haar law of $UP_AU^*$ is the invariant Grassmannian law. Lemma~\ref{lem:grassmann57} bounds the measure by
$C_{d,m}(2h/\gamma_d)^{2m(d-m)}$ for $h\le\gamma_d/2$. Since $m(d-m)\ge d-1$, and there are finitely many $m$, this gives \eqref{eq:conjugacycap57} with a constant independent of all eigenvalue multiplicities. Larger $h$ are covered by increasing $C_d$.

For sharpness take $A=B=(P-c_d)/r_d$ with $P$ rank one. By \eqref{eq:projectivecap57}, for sufficiently small $h$ the probability is
\[
                    (r_d^2h^2/2)^{d-1}.
\]
It cannot be bounded by a constant times $h^{p'}$ as $h\downarrow0$ for any $p'>2(d-1)$.
\end{proof}
The spectral split, rather than an assumption that all centroids remain rank one, supplies the uniformity needed by the stochastic contraction. The conjugation representation has ambient dimension $d^2-1$, but its worst small-ball exponent is $2(d-1)$.

\begin{theorem}[Unitary matrix-orbit width]\label{thm:matrixwidth57}
Let a finite alphabet in $\mathrm{SU}(d)$ contain the identity and support a probability law satisfying \eqref{eq:groupgap57}. Fix a rank-one projection $P_0$. For $0<a\le1$, set
\begin{equation}\label{eq:matrixexperiment57}
 F_a(U)=c_d+a(UP_0U^*-c_d),\qquad Y=\mathcal D_d,
\end{equation}
with Frobenius error. Its exact bounded-width threshold is $a r_d$, attained by the one-label decoder $c_d$. For $0\le\varepsilon<a r_d$, put $\zeta=a-\varepsilon/r_d$. Every realization satisfies
\begin{equation}\label{eq:matrixoccupation57}
 \#\{1\le t\le N:|S_t|\le k\}
 \le C\log(k+2)\left(1+k^{1/(d-1)}\log\frac1\zeta\right).
\end{equation}
For fixed $0<a<1$ and $0\le\varepsilon<a r_d$,
\begin{equation}\label{eq:matrixrate57}
 c\left(\frac{N}{\log(N+2)}\right)^{d-1}
            \le W_{N,\varepsilon}\le C(N+1)^{d-1},
\end{equation}
and the upper realization can be exact. The same orders hold for $a=1$ and every fixed $0<\varepsilon<r_d$.

More explicitly, for $a<1$ and $\varepsilon=0$ set $a'=a$; for $\varepsilon>0$ set $a'=a-\varepsilon/(2r_d)$. Whenever $a'<1$, an error at most $\varepsilon/2$ realization has
\begin{equation}\label{eq:matrixupper57}
 K\le C_d\left(1+\frac{N+1}{-\log a'}\right)^{d-1}
\end{equation}
labels. In the zero-error case this realization is exact. In \eqref{eq:matrixoccupation57} the constant depends only on $d$ and the spectral-gap data, not on the amplitude or tolerance. Extra finitely many unitary commands are allowed in both bounds.
\end{theorem}
\begin{proof}
Center at $c_d$ and divide all outputs by $r_d$. The resulting legal set is
\[
 Q=\operatorname{conv}\{(P-c_d)/r_d:P\text{ rank one}\}\subseteq\mathbb B_{\mathcal H_d^0}.
\]
The last inclusion follows from $\operatorname{tr}D^2\le1$ for $D\in\mathcal D_d$. The orbit spans $\mathcal H_d^0$: if a traceless Hermitian $A$ has $\operatorname{tr}(AP)=0$ for every rank-one $P$, then $v^*Av=0$ for every $v$, hence $A=0$. Conjugation-fixed matrices are scalar, so there are no fixed vectors in this traceless space. The orbit is complex projective space, of real dimension $2(d-1)$. Lemma~\ref{lem:conjugacysmall57} and Theorem~\ref{thm:uniformorbit57}, with $p=q=2(d-1)$, prove the lower and occupation bounds, as well as the radius statement after multiplying the norm by $r_d$.

We give a direct inner approximation to make positivity of the upper rows transparent. Choose a maximal $\delta$-separated family of rank-one projections $P_i$ for $\operatorname{dist}_{\rm pr}$, containing $P_0$. Its $\delta/2$ balls are disjoint, so \eqref{eq:projectivecap57} gives
\[
                         K\le(2/\delta)^{2(d-1)}.
\]
For any traceless Hermitian $A$, put $M=\lambda_{\max}(A)$. When $A\ne0$, $M>0$ and $\lambda_{\min}(A)\ge-(d-1)M$. Approximate a top eigenline by $P_i$ at distance at most $\delta$. Spectral decomposition then gives
\begin{equation}\label{eq:matrixsupport57}
 \operatorname{tr}(AP_i)\ge M-(M-\lambda_{\min}(A))\delta^2
                         \ge(1-d\delta^2)M.
\end{equation}
The trivial $A=0$ case is automatic. Comparing support functions in $\mathcal H_d^0$, with $K_0=\mathcal D_d-c_d$, yields
\[
 (1-d\delta^2)K_0\subseteq\operatorname{conv}\{P_i-c_d\}\subseteq K_0.
\]
Put $r=1-d\delta^2>0$. For each actual command $U_b$ choose one row of barycentric coordinates satisfying
\begin{equation}\label{eq:matrixrow57}
 \sum_jT_b(i,j)(P_j-c_d)=rU_b(P_i-c_d)U_b^*.
\end{equation}
Rows are nonnegative and sum to one by construction. Initialize at $P_0$ and decode label $i$ as
\[
                  D_i=c_d+(a'/r^N)(P_i-c_d).
\]
This is in $\mathcal D_d$ when $r^N\ge a'$: it is a convex combination of $c_d$ and $P_i$. Choosing
\[
 \delta^2=\min\left(\frac1{4d},\frac{-\log a'}{2d(N+1)}\right)
\]
ensures this inequality. Induction in \eqref{eq:matrixrow57} gives mean output exactly $F_{a'}(u_w)$. The target discrepancy is $|a-a'|r_d=\varepsilon/2$ when $\varepsilon>0$. The cardinal bound proves \eqref{eq:matrixupper57}. The case $a=1,\varepsilon=0$ is deliberately excluded from this dilation construction and is treated exactly in the next section.
\end{proof}
The family contains the qubit sphere as $d=2$, but for $d\ge3$ the pure-state orbit is a proper curved subset of the unit sphere in $\mathcal H_d^0$. Thus \eqref{eq:matrixrate57} does not follow by substituting the ambient dimension into the spherical theorem. Its logarithmic gap is explicit; no sharp constant or disappearance of that gap is claimed.

\begin{corollary}[Algebraic generating alphabets]\label{cor:algebraicunitary57}
Theorem~\ref{thm:matrixwidth57} applies to the identity together with any finite symmetric algebraic-entry set generating a dense subgroup of $\mathrm{SU}(d)$, equipped with a lazy symmetric law. Such a finite alphabet exists for every $d\ge2$.
\end{corollary}
\begin{proof}
The spectral gap is the imported theorem of Bourgain--Gamburd \cite{BG2012}. Laziness makes its one-sided gap an absolute $L^2$ contraction on mean-zero functions. Existence can be seen by placing two algebraic irrational-angle $\mathrm{SU}(2)$ rotations in every adjacent pair of coordinate directions. For example, use the quaternion lifts in Corollary~\ref{cor:SO356} within each pair. Each pair generates its full $\mathrm{SU}(2)$ subgroup. Their Lie algebras contain the adjacent off-diagonal matrix units and their brackets, which generate $\mathfrak{su}(d)$. The connected compact closure is consequently $\mathrm{SU}(d)$. All entries remain algebraic. The spectral-gap constant is not computed by the finite regression checks.
\end{proof}

\subsection{A finite whole-law encoding}
Let $q_0=d^2-1$ and choose a Frobenius-orthonormal basis $H_1,\ldots,H_{q_0}$ of $\mathcal H_d^0$. Since $\|H_l\|_{\rm op}\le1$, the $2q_0$ matrices
\[
 E_{l,\pm}=\frac{I\pm H_l/2}{2q_0}
\]
are positive and sum to $I$. Put $\Phi(D)_{l,\pm}=\operatorname{tr}(E_{l,\pm}D)$ and $Y_\Phi=\Phi(\mathcal D_d)$. Every entry is at least $1/(4q_0)$. The map is injective, because
\begin{equation}\label{eq:tomography57}
 D-c_d=\sum_l2q_0\bigl(\Phi(D)_{l,+}-\Phi(D)_{l,-}\bigr)H_l.
\end{equation}
On the vector space parallel to the affine span of $Y_\Phi$, use the norm of the reconstructed matrix on the right. Then $\Phi$ is an affine isometry, and replacing each legal matrix decoder by its image preserves every clocked or stationary width exactly. This is one whole categorical law with a compact convex legality constraint, not separate binary event tolerances. For ordinary total variation the explicit norm comparison is
\[
 \frac{\|D-D'\|_{\rm F}}{4q_0}
 \le\TV(\Phi(D),\Phi(D'))
 \le\frac{\|D-D'\|_{\rm F}}{4\sqrt{q_0}}.
\]
It follows directly by comparing the $\ell^1$ and $\ell^2$ norms of the basis coefficients. Enlarging the legal output set from $Y_\Phi$ to the entire simplex is a different realization problem. In particular the extreme-output argument below cannot be transferred to that enlarged set: its target vectors are interior points of the simplex.

\section{Extreme outputs and the exact endpoint}\label{sec:extreme57}
The strict amplitude slack in the polytope construction is not a technical
convenience that can simply be removed. At an extreme output, convex
averaging has no room to hide a different physical state. This observation
gives an exact finite-horizon formula, independently of spectral gap or
synchronization.

\begin{theorem}[The extreme-orbit formula]\label{thm:extreme57}
Let a group act by affine bijections on a compact convex set $Y$. Let
$x_0\in Y$ have an orbit consisting of extreme points of $Y$, and let a
finite nonempty command alphabet act through this group. There is one seed
and one terminal query, with target $u_w x_0$. Put
\[
 \mathcal O_t=\{u_wx_0:|w|=t\},\qquad n_t=|\mathcal O_t|.
\]
Every exact clocked stochastic realization at horizon $N$ satisfies
\begin{equation}\label{eq:extremeprofile57}
                         |S_t|\ge n_t\quad(0\le t\le N).
\end{equation}
One deterministic realization attains all these inequalities simultaneously.
Consequently
\begin{equation}\label{eq:extremewidth57}
                         W_{N,0}=n_N.
\end{equation}
Neither an identity letter nor an exact return is required.
\end{theorem}
\begin{proof}
Fix $t$ and one suffix $v$ of length $N-t$. Such a suffix exists because the
alphabet is nonempty; at $t=N$ use the empty word. Let $D_v(i)\in Y$ be the
conditional expected terminal decoder after this fixed suffix, when the
cut-$t$ label is $i$. Convexity guarantees that it lies in $Y$. For every
prefix $w$ of length $t$, exact correctness says
\[
             u_vu_wx_0=\sum_i p_w(i)D_v(i),
\]
where $p_w$ is the actual distribution at that cut. Extremality forces
$D_v(i)=u_vu_wx_0$ whenever $p_w(i)>0$. Since the action of $u_v$ is
injective, prefixes with different physical outputs have disjoint
nonempty supports in $S_t$. This proves \eqref{eq:extremeprofile57}, including
the terminal cut.

Conversely, at cut $t$ store the current point of $\mathcal O_t$. A command
sends this point to its image in $\mathcal O_{t+1}$; the terminal decoder is
the stored point. The initialization is deterministic. Finally, any one
fixed command gives an injection $\mathcal O_t\to\mathcal O_{t+1}$, so $n_t$
is nondecreasing. The peak of this realization is $n_N$, proving
\eqref{eq:extremewidth57}.
\end{proof}
The same conclusion holds for an orbit on a Euclidean unit sphere with
legal output set its convex hull: every orbit point is extreme by strict
convexity of the ambient Euclidean ball. For density matrices, a rank-one
projection is extreme. Indeed, if $P=\sum_i p_iD_i$ with $D_i\ge0$ and
$\tr D_i=1$, then $\tr((I-P)D_i)=0$ for every positive weight. Positivity
forces $D_i$ to be supported on the range of $P$, and hence $D_i=P$.

\subsection{A fixed nonabelian family with an exponential exact endpoint}
The next example uses two established inputs, explicitly separated from
the realization argument. Breuillard and Gelander prove that a dense
subgroup of a connected semisimple real Lie group contains a dense free
subgroup on two generators \cite[Theorem~1.1]{BreuillardGelander57}; their
later correction supplies repairs and details for other parts of the
original proof \cite{BreuillardGelanderCorrection57}. The spectral-gap input
for dense algebraic unitary generators is \cite{BG2012}. No algorithm for
finding free generators, or for certifying a spectral-gap constant, is
asserted here.

\begin{lemma}[A line with trivial algebraic projective stabilizer]
\label{lem:transcendentalseed57}
Fix a real number $\tau$ transcendental over $\Q$ and put
\[
 v_\tau=\frac{(1,\tau,\ldots,\tau^{d-1})^{\mathsf T}}
                  {(\sum_{j=0}^{d-1}\tau^{2j})^{1/2}},\qquad
 P_\tau=v_\tau v_\tau^*.
\]
An invertible matrix with algebraic entries that preserves the line
$\C v_\tau$ is scalar.
\end{lemma}
\begin{proof}
If $Mv_\tau=\lambda v_\tau$, the eigenvalue $\lambda$ is algebraic because
it is a root of the characteristic polynomial of the algebraic matrix $M$.
Each row of $(M-\lambda I)(1,\tau,\ldots,\tau^{d-1})^{\mathsf T}=0$ is a
polynomial identity in $\tau$ with algebraic coefficients. Transcendence
over $\Q$ implies transcendence over the field of algebraic numbers: an
algebraic relation over a finite number field would make $\tau$ algebraic
over $\Q$ by transitivity of algebraic extensions. Every row polynomial
therefore has all coefficients zero. Thus $M=\lambda I$.
\end{proof}
One may take a fixed computable transcendental such as
$\tau=\sum_{n=1}^\infty10^{-n!}$. For completeness, its truncations have
rational denominator $10^{m!}$ and strictly positive tail at most
$2\,10^{-(m+1)!}$. For a nonzero integer polynomial $A$ of degree $D$, a
rational $p/q$ not among its finitely many roots satisfies
$|A(p/q)|\ge q^{-D}$. If $A(\tau)=0$, a bounded derivative near $\tau$
would instead bound $|A(p/q)|$ by a constant times $q^{-(m+1)}$ along these
truncations, a contradiction for large $m$. This also records precisely
why the chosen seed is not part of the rational-input specialization.

\begin{theorem}[Exponential exact width and polynomial noisy width]
\label{thm:freeendpoint57}
For every $d\ge2$ there is a fixed five-letter algebraic unitary alphabet
\[
                       \{I,U,U^{-1},V,V^{-1}\}\subset SU(d)
\]
and a fixed rank-one seed $P_\tau$ as above such that the full matrix
interface $F(g)=gP_\tau g^*$, with legal decoder set $\mathcal D_d$, has
\begin{equation}\label{eq:freeexact57}
                      W_{N,0}=2\,3^N-1\qquad(N\ge0).
\end{equation}
For each fixed $0<\varepsilon<r_d=\sqrt{1-1/d}$, its Frobenius-error width
instead satisfies
\begin{equation}\label{eq:freenoise57}
 c_{d,\varepsilon}\left(\frac{N}{\log(N+2)}\right)^{d-1}
 \le W_{N,\varepsilon}\le C_{d,\varepsilon}(N+1)^{d-1}.
\end{equation}
At $\varepsilon\ge r_d$ one label suffices. For the same alphabet and seed,
every fixed amplitude $0<a<1$ admits the exact polynomial upper bound in
\eqref{eq:matrixupper57} and the corresponding lower bound in
\eqref{eq:matrixrate57}.
\end{theorem}
\begin{proof}
The subgroup of $SU(d)$ consisting of matrices with algebraic entries is
dense. One direct verification is to approximate a unitary matrix by a
nonsingular matrix with entries in $\Q+i\Q$, apply Gram--Schmidt, and
multiply its final column by the inverse determinant of the resulting
unitary matrix. All entries produced are algebraic; continuity near the
original unitary matrix gives density. Apply
\cite[Theorem~1.1]{BreuillardGelander57} to this subgroup. It contains
algebraic $U,V$ freely generating a dense free group $\Gamma$ of rank two.
The scalar matrices in $SU(d)$ form a finite group. Since a free group is
torsion-free, $\Gamma$ contains no nonidentity scalar matrix.
Lemma~\ref{lem:transcendentalseed57} now shows that the orbit map
$\Gamma\to\mathcal D_d$, $g\mapsto gP_\tau g^*$, is injective.

With the identity available for padding, the products at length $N$ are
exactly the elements of reduced free-word length at most $N$. There are
four choices for the first letter of a nonempty reduced word and three
for each subsequent letter. Thus the ball has cardinality
\[
               1+4\sum_{j=1}^N3^{j-1}=2\,3^N-1.
\]
Apply Theorem~\ref{thm:extreme57} to obtain \eqref{eq:freeexact57}. In fact the
entire pointwise minimum profile is $|S_t|=2\,3^t-1$.

Assign a symmetric positive law to the four nonidentity commands and
positive mass to the identity. The algebraic dense-generator theorem
\cite{BG2012}, with laziness, gives \eqref{eq:groupgap57}. Consequently
Theorem~\ref{thm:matrixwidth57} applies to this same fixed alphabet and
seed. It yields \eqref{eq:freenoise57}, the threshold $r_d$, and the asserted
strict-amplitude conclusion.
\end{proof}
This is a full legal-matrix-output statement. It does not resolve the
zero-error unit-amplitude problem for a single scalar planar readout:
that readout need not expose every orbit point as an extreme decoder
value. Nor does it identify a sampled measurement outcome with a density
matrix stored in a classical label. The distinction between a prescribed
convex output set and a larger simplex is essential here.

\section{Effective algebraic-circle bounds}\label{sec:arithmetic53}
For algebraic unit-circle eigenvalues of infinite order, an explicit separation estimate replaces the Diophantine hypothesis of the preceding section. The parameters $r,m,A_r$ below are those of \eqref{eq:r53}--\eqref{eq:circleconstants53}; in this section $r$ denotes the localization order, not the number of command rotations. The alphabet contains $I,R_\theta$, and one interface mean is $\rho\cos\phi$.

\begin{lemma}[Separated powers give a block defect]\label{lem:separated53}
Suppose $|\lambda^n-1|\ge\sigma$ for every integer $1\le n\le2mk$, where $0<\sigma\le1$. Write $R_m(t)=\diag(R_t,R_{2t},\ldots,R_{mt})$ for the representation acting on $Z_m$. On this representation, the uniform law on the $2k$ words with products $R_\theta^j$, $0\le j<2k$, each padded to length $2k$, has defect at least $\sigma^2/16$ in \eqref{eq:gap54}.
\end{lemma}
\begin{proof}
Write a unit vector as $u=(u_1,\ldots,u_m)$ in its two-dimensional frequency blocks. For distinct indices $i,j\in\{0,\ldots,2k-1\}$,
\[
 \norm{R_m(i\theta)u-R_m(j\theta)u}^2
 =\sum_{l=1}^m\norm{u_l}^2|\lambda^{l(i-j)}-1|^2\ge\sigma^2.
\]
A ball of radius strictly less than $\sigma/2$ about any of $k$ centers contains at most one of these $2k$ points. Thus at least $k$ points have distance at least $\sigma/2$ from all centers (the same conclusion follows by taking the open balls of radius $\sigma/2$). For unit centers the defect is one half of the squared distance to the closest center. Averaging gives at least $(1/2)(\sigma^2/8)=\sigma^2/16$.
\end{proof}

\begin{theorem}[Full-noise algebraic-circle bound]\label{thm:algebraic53}
Assume $\lambda$ is algebraic and not a root of unity. Let its primitive irreducible integer polynomial be
\[
 P(z)=a z^d+a_{d-1}z^{d-1}+\cdots+a_0,\qquad a>0.
\]
Set $H_P=\max(a,|a_{d-1}|,\ldots,|a_0|)$ and
\[
                     B_P=a(1+H_P/a)^{d-1}>1.
\]
With $r,m,A_r$ as in \eqref{eq:r53}--\eqref{eq:circleconstants53}, every error-$\varepsilon$ realization of peak at most $k$ satisfies
\begin{equation}\label{eq:algwidth53}
 N<2k\left(1+16\,2^{2(d-1)}B_P^{4mk}\log A_r\right).
\end{equation}
Its number of command cuts $t\in\{1,\ldots,N\}$ with width at most $k$, even if other widths are unrestricted, is at most
\begin{equation}\label{eq:algoccupation53}
 2k-1+32k\,2^{2(d-1)}B_P^{4mk}\log A_r.
\end{equation}
Consequently, for fixed algebraic data, $\rho$, and $\varepsilon<\rho/2$,
\begin{equation}\label{eq:algrate53}
 W_{N,\varepsilon}\ge
 \frac{\log N-O(\log\log N)}{4m\log B_P}.
\end{equation}
All constants in the finite inequalities are explicit. If $\lambda=(a_1+ib_1)/c$ with integers $a_1,b_1$, $c>0$, $a_1^2+b_1^2=c^2$, and $\lambda$ not a root of unity, then \eqref{eq:algwidth53}--\eqref{eq:algrate53} hold with the factor $2^{2(d-1)}$ deleted and $B_P$ replaced by $c$.
\end{theorem}
\begin{proof}
List the conjugates with $\lambda_1=\lambda$. They are distinct by irreducibility in characteristic zero. None is a root of unity: otherwise irreducibility would make $P$ a cyclotomic polynomial up to its primitive sign. The resultant
\[
 \operatorname{Res}(P,z^n-1)=a^n\prod_{j=1}^d(\lambda_j^n-1)
\]
is therefore a nonzero integer. Each conjugate has modulus at most $R_P=1+H_P/a$. For completeness, if $|z|>R_P$, then
\[
 \sum_{j=0}^{d-1}|a_j||z|^j
 <\frac{H_P|z|^d}{|z|-1}<a|z|^d,
\]
so $P(z)\ne0$. Using $|z^n-1|\le2\max(1,|z|)^n$ for each other conjugate, the resultant bound gives
\begin{equation}\label{eq:resultant53}
 |\lambda^n-1|\ge
 2^{-(d-1)}\left(a\prod_{j=2}^d\max(1,|\lambda_j|)\right)^{-n}
 \ge 2^{-(d-1)}B_P^{-n}.
\end{equation}
For $1\le n\le2mk$, take $\sigma=2^{-(d-1)}B_P^{-2mk}$ in Lemma~\ref{lem:separated53}. The block defect is at least
\[
                  d_k=\frac{2^{-2(d-1)}B_P^{-4mk}}{16}.
\]
Repeated conditional contraction and Lemma~\ref{lem:harmonic53} imply
\[
 \left\lfloor\frac{N}{2k}\right\rfloor
 \le \frac{\log A_r}{-\log(1-d_k)}
 \le 16\,2^{2(d-1)}B_P^{4mk}\log A_r.
\]
This proves \eqref{eq:algwidth53}. The greedy occupation argument in the proof of Theorem~\ref{thm:radius54}, with $b=0$ and block length $2k$, proves \eqref{eq:algoccupation53}. Taking logarithms gives \eqref{eq:algrate53}: fix $C>1/(4m\log B_P)$. If $k\ge C\log N$ the claim is immediate; otherwise $\log k\le\log C+\log\log N$, which can be absorbed in the logarithm of \eqref{eq:algwidth53}. The implicit constant may depend on $P,\rho,\varepsilon$, not on $N$.

In the rational-circle case, $(a_1+ib_1)^n-c^n$ is a nonzero Gaussian integer, of modulus at least one. Thus $|\lambda^n-1|\ge c^{-n}$ directly. Repeat the proof with $\sigma=c^{-2mk}$. Here $c>1$ automatically. A nonminimal supplied denominator is still valid but weakens the bound.
\end{proof}
The threshold is sharp; these width rates are not asserted sharp. The factor $m$ records the increased localization degree as the error approaches the boundary. At smaller errors, the retained coordinate-calibration estimates can have better constants. One should use the larger of applicable lower bounds, rather than claiming numerical improvement at every parameter value.

\begin{proposition}[An explicit nonrational algebraic example]\label{prop:salem53}
Let $t=(1-\sqrt{13})/2$ and
\[
 \lambda=\frac{t+i\sqrt{4-t^2}}2.
\]
Then $|\lambda|=1$, $\lambda$ is not a root of unity, and one can take $d=4$ and $B_P=8$ in Theorem~\ref{thm:algebraic53}.
\end{proposition}
\begin{proof}
The number $t$ lies strictly between $-2$ and $2$ and solves $t^2-t-3=0$. Substitution of $t=z+z^{-1}$ gives
\[
                       P(z)=z^4-z^3-z^2-z+1.
\]
Its reduction modulo two is $z^4+z^3+z^2+z+1$. It has no linear factor over $\mathbb F_2$ and is not divisible by the only monic irreducible quadratic $z^2+z+1$; hence it is irreducible. Gauss's lemma proves irreducibility over $\mathbb Q$. The conjugate value $(1+\sqrt{13})/2>2$ gives a positive real conjugate of $\lambda$ greater than one. A root of unity cannot have such a conjugate. Finally $a=H_P=1$ and $d=4$, so $B_P=2^3=8$.
\end{proof}
For a reproducible rational calibration, take $\rho=1/10$ and $\varepsilon=1/25$. Then $\delta=2/25$, $r=5$, $m=6$, $\Delta_r=1/150$, and $A_r=23364$. Thus the rational-circle alphabet above obeys
\[
                 N<2k\bigl(1+16\cdot5^{24k}\log23364\bigr).
\]
This is a concrete bound inside the interval between the former small-error obstruction and the sharp threshold $1/20$. The exact coefficient and resultant checks accompanying the article verify the displayed arithmetic, not the universal theorem by finite experimentation.


\section{Finite encodings and boundary quotients}\label{sec:effective54}
The general continuous theorem is not an algorithm from arbitrary real data. We now give an explicit finite input class for which its conclusions are effective. A rational orthogonal input consists of a finite alphabet in $O(D,\Q)$ containing $I$, finite seed and query sets, and rational polynomials $F_{s,j,l}$ in the matrix entries specifying categorical probabilities. The promise $F_{s,j}(H)\subset\Delta_{M_j}$ can itself be checked once the compact closure is computed. Real algebraic output numbers are encoded by polynomial equations and isolating data.

We use two established algorithmic inputs: computation of a finite defining description of the Zariski closure of a finitely generated rational matrix group \cite{DJK,NPSW}, and real quantifier elimination with effective semialgebraic connected components \cite{BPR}. The former is a terminating closure algorithm, not merely an enumeration of valid polynomial identities with an unknown stopping point. Compact orthogonal groups are real algebraic, so the real points of the computed closure in $O(D)$ give the Euclidean closure $H$ \cite{BJKP}. The complex algebraic components must not be substituted for the real connected components used here.

\begin{theorem}[Effective rational classification]\label{thm:effective54}
For a rational orthogonal input, the following are computable:
\begin{enumerate}
\item a semialgebraic description of every connected component of $H$, the finite group $C$, and the component of each command;
\item the critical total-variation error $\epsc$ as an exact real algebraic number, and an attaining stationary component-permutation machine with algebraic decoders;
\item for each algebraic $0\le\varepsilon<\epsc$ and integer $k\ge1$, an integer occupation constant $L_k$ valid for every error-$\varepsilon$ clocked realization;
\item for every fixed $N$ and algebraic $\varepsilon\ge0$, the least width $W_{N,\varepsilon}$ and an algebraic realization attaining it.
\end{enumerate}
These are terminating computability statements. No polynomial running-time bound, practical implementation of general closure/quantifier elimination, or finite-random-bit cost bound is asserted.
\end{theorem}
\begin{proof}
Compute $H$ using the cited closure algorithm, then its real connected components and algebraic sample points. The component containing $I$ is $K$. Products of sample points and semialgebraic membership tests determine the component multiplication table and the command images. In particular $K$ is now an effectively given compact semialgebraic group.

For a component $c$ and a query, the condition that $q$ be a radius-$t$ center is the real first-order formula
\begin{equation}\label{eq:QEcenter54}
 q\in\Delta_{M_j},\qquad
 \forall h\in c:\quad
 \sum_{l=1}^{M_j}|F_{s,j,l}(h)-q_l|\le2t.
\end{equation}
Absolute values have a finite semialgebraic description. Quantifier elimination gives the set of admissible $(q,t)$. Its least $t$ exists by compactness and is real algebraic; algebraic sampling supplies an attaining $q$. Maximize over the finite interface and components. This proves (1)--(2), including equality at the critical error.

We detail effectivity of the occupation proof, since computing the threshold alone would not supply it. Choose a component with radius greater than the given $\varepsilon$, and enumerate positive words until one lies in that component. Its rational product gives the prefix used in the proof of Theorem~\ref{thm:radius54}. The resulting map $f(h)=F_{s,j}(hu_v)$ is polynomial, so it can be used directly in place of $G$, with no approximation error.

For any finite tensor degree, all required polynomial functions are coefficients of the explicitly given tensor representation. Its fixed subspace is
\[
 \{z:\ \forall h\in K,
                 (R(h)-I)z=0\}.
\]
This linear subspace is effectively semialgebraic. Real algebraic sampling, successively outside the span of the vectors already chosen, gives an algebraic basis in finitely many steps. Its Gram matrix gives the orthogonal projection $P$ exactly. Thus every polynomial Haar moment is computable algebraically from $P\Phi(e)$; numerical integration of Haar measure is unnecessary.

Enumerate finite lists of squares of rational polynomials on $K$, discarding zero Haar integrals and normalizing the others. For a list $p_i$, compute $b_i=\int p_if\,d\mu$ and its constrained simplex radius by quantifier elimination. Rational polynomials are uniformly dense among real polynomial restrictions on a compact matrix group; the localization proof therefore guarantees that some enumerated list has radius strictly greater than $\varepsilon$. Strict algebraic comparison detects that list. Fix a positive rational $\Delta$ below the radius difference and a positive rational upper bound $L$ for all coefficient constants in Proposition~\ref{prop:residual54}. These give the certified residual $Q\ge\Delta/L$.

Next enumerate finite rational probability laws on executable words of a common length whose products lie in $K$, together with rational $d\in(0,1)$. The assertion \eqref{eq:gap54}, quantified over the finitely many unit vectors in the algebraic moving space, is semialgebraic: each maximum of finitely many inner products can be expressed by a finite disjunction. It is therefore decidable. Lemma~\ref{lem:gap54} and sufficiently small rational perturbations of its atom weights ensure that this search terminates. Padding by the identity letter makes every finite choice executable at one common length.

The resulting $B,d,L,\Delta,Q_0$ determine \eqref{eq:Lk54}. To avoid exact comparisons involving logarithms, take a rational $A'\ge\max(1,LQ_0/\Delta)$ and use
\[
 \frac{\log A}{-\log(1-d)}\le\frac{A'}d.
\]
An integer above $b+B-1+BA'/d$ is a computable occupation bound. This proves (3).

For (4), fix a candidate width $k$ and use $k$ labels at every cut, padding smaller registers with unused labels. Stochastic rows, initialization, and categorical decoder probabilities are finitely many variables in compact simplexes. For each of the finitely many length-$N$ words, the output vector is polynomial in these variables, and the target vector is rational. The total-variation inequalities are semialgebraic. Quantifier elimination therefore decides feasibility and supplies an algebraic point when feasible. A deterministic register storing the seed and entire prefix is a finite exact realization, so increasing $k$ must terminate. This proves the least-width claim for each finite horizon, without claiming a finite-horizon cutoff for all stationary or nonuniform realization questions.
\end{proof}
The distinction between an effective theorem and an implemented solver matters. The accompanying finite checks verify displayed exact examples and proof constants; they do not execute the general group-closure and quantifier-elimination construction. The finite-bit compiler for rational finite groups remains in the complete predecessor supplement with its separate hypotheses and costs. It is not used to reprice arbitrary algebraic decoders as cost-free finite-bit procedures.

\subsection{The least component-quotient realization}
The upper bound $|\mathcal S||C|$ can often be reduced. Here is an exact finite characterization of that reduction, with its scope stated explicitly. Put $X=\mathcal S\times C$. A \emph{component-quotient realization} has a deterministic state of the form $\pi(s,c)$, where $\pi:X\to Q$ is a surjection; its command update is induced by $(s,c)\mapsto(s,[u_a]c)$. It is stationary and has no additional state recording a representative word within a component.

Call a partition $\mathcal B$ of $X$ equivariant if
\[
 (s,c)\sim(s',c')\quad\Longrightarrow\quad
 (s,dc)\sim(s',dc')\qquad(d\in C).
\]
For a block $B\in\mathcal B$ and query $j$, set
\[
 A_{B,j}=\bigcup_{(s,c)\in B}F_{s,j}(c).
\]
\begin{theorem}[Finite quotient minimum]\label{thm:quotient54}
The least number of command labels in an error-$\varepsilon$ component-quotient realization is
\begin{equation}\label{eq:quotient54}
 \min\left\{|\mathcal B|:\mathcal B\text{ equivariant and }
         \rad_{Y_j}(A_{B,j})\le\varepsilon\text{ for every }B,j\right\},
\end{equation}
with value $+\infty$ if there is no admissible partition. For rational orthogonal categorical inputs this minimum and an attaining quotient are computable. If $\epsc=0$ and $\varepsilon=0$, the coarsest admissible partition is the future-output equivalence
\[
 (s,c)\sim(s',c')\quad\Longleftrightarrow\quad
 F_{s,j}(dc)=F_{s',j}(dc')\quad\text{for every }d\in C,j,
\]
where componentwise constant maps are identified with their values.
\end{theorem}
\begin{proof}
An equivariant partition induces well-defined command permutations on its blocks. Choosing a minimizing center for each $A_{B,j}$ gives exactly the stated error. Conversely, the fibers of $\pi$ in any component-quotient realization form an equivariant partition. Its decoder must approximate every executable product in each corresponding component. Such products are dense in that component, so continuity extends the error bound to $A_{B,j}$. Its decoder is therefore a permitted center. This proves \eqref{eq:quotient54}. The command generators suffice to check equivariance, since their images generate the finite group $C$.

There are finitely many partitions of $X$. For polynomial categorical data their radius tests are of the form \eqref{eq:QEcenter54}, with a finite disjunction over the relevant seed/component pairs, so the minimum is computable. In the exact componentwise constant case, equality of all future outputs is an equivariant equivalence relation. Every admissible partition refines it, and its own blocks have singleton output images for each query. It is therefore the unique coarsest admissible partition.
\end{proof}
The minimization in \eqref{eq:quotient54} is a deterministic quotient minimum, distinct from the unrestricted stationary minimum in Theorem~\ref{thm:stationaryminimum55} and from the unrestricted nonuniform minimum at the boundary. Randomized encodings can identify numerical behaviors without identifying component pairs. What is classified here is exactly the finite deterministic quotient construction that attains the radius threshold. Theorem~\ref{thm:effective54}(4) separately computes the unrestricted minimum at each fixed horizon.

\subsection{Stationary minimization by permutation closures}
\begin{theorem}[Effective unrestricted stationary minimum]\label{thm:effectivemin55}
For rational orthogonal commands generating $H$, rational polynomial categorical outputs, and algebraic $\varepsilon\ge0$, one can compute $M_\varepsilon$ and an attaining algebraic machine when it is finite. An identity letter is not required. At each proposed width $k$, at most $(k!)^{|\mathcal A|}$ augmented orthogonal group closures and finite real quantifier-elimination problems suffice.
\end{theorem}
\begin{proof}
First compute $H$, its real connected components, and their constrained radii as in Theorem~\ref{thm:effective54}(1)--(2); those steps do not use an identity letter. The group $H\subset O(D)$ has finitely many components. Theorem~\ref{thm:finitequotient55} shows that $M_\varepsilon=\infty$ below their maximum radius. Otherwise $M_\varepsilon\le |\mathcal S||H/H^\circ|$.

For $k$ between one and this upper bound enumerate all command-permutation tuples $\sigma$. The matrices
\[
                         \diag(u_a^{-1},\Pi_a)\in O(D+k,\Q)
\]
generate a compact closure which is exactly the joint group $\mathcal L_\sigma$ of Theorem~\ref{thm:stationaryminimum55}. Compute it by the same algebraic closure procedure. In \eqref{eq:permutationfeasibility55}, the unknown seed rows and decoder rows are constrained to simplexes. The mean output is bilinear in these rows, the group variable is restricted by computed defining equations whose real zero set inside the augmented orthogonal group is exactly $\mathcal L_\sigma$, not merely by some equations vanishing on it, and $h^{-1}=h^T$ in the physical block. Consequently all target coordinates are polynomial in the universally quantified matrix entries. Total variation is semialgebraic. Real quantifier elimination decides the resulting existential--universal formula and supplies algebraic rows whenever it is feasible. The first feasible $k$ is the minimum by purification. The finite upper bound proves termination of the entire search.
\end{proof}
This is a structural reduction of stationary minimization, not merely an enumeration of arbitrary stochastic transition parameters. It does not equate stationary and finite-horizon minima. Theorem~\ref{thm:stationarization56} identifies their limit under its additional return hypothesis. The two-cycle/three-cycle example in Proposition~\ref{prop:C655} already prevents replacing the permutation-feasibility test by deterministic component partitions.

\subsection{Algorithmic cost and arithmetic representation}
There are three distinct layers of effectivity. First, the group-closure procedure produces explicit defining equations; real-component decomposition and algebraic sampling then identify the component table. Second, radius computation and the stationary permutation tests are finite quantifier-elimination tasks on those equations. At fixed $k$, the stationary seed and output variables number $k(|\mathcal S|+\sum_jM_j)$ before normalization equalities are eliminated; the group variables can be taken among the $(D+k)^2$ augmented matrix entries. Thus the displayed finite search is independent of a horizon parameter. Its factorial number of discrete cases and the cost of the real-algebraic subproblems should not be mistaken for an efficient algorithm. No bound for the full closure-to-machine process in terms of only the original bit length is extracted here.

Third, occupation certificates in Theorem~\ref{thm:effective54}(3) add searches over localizer degrees and executable word laws. Their termination follows from strict analytic margins; the present proof supplies no a priori useful bound on the successful degree or word length. Fixed-horizon minimization, by contrast, has finitely many word constraints at each $N$ but their number is exponential in $N$ before repetitions are identified. These stages should be priced separately rather than hidden behind the word ``effective.''

All algebraic Haar moments in the localizer search are represented in a real algebraic number field with isolating data, using the fixed-subspace projector and its algebraic Gram inverse. Signs and strict comparisons are exact real-algebraic comparisons; a quadrature tolerance is not a certificate. Every rational command word has a rational matrix product and hence a rational target vector for rational polynomial readouts. The orthogonal closure fact used above is precisely that the Euclidean closure of a subgroup of $O(D)$ is the real zero set, inside $O(D)$, of its vanishing polynomial ideal, as used in \cite{BJKP}; replacing real connected components by complex algebraic ones would give the wrong quotient.

\section{Finite descriptions and decision complexity}\label{sec:complexity56}
The qualitative reduction is insensitive to the encoding of real numbers. Complexity statements are not. We first specify a finite input model and then give an explicit finite-precision counterpart to label width.

\begin{theorem}[Finite-group stationary minimization is $\exists\R$-complete]\label{thm:ETR56}
The input consists of a finite group $G$ given by its full multiplication table, finite seed and query sets, rational categorical vectors $F_{s,j}(g)$ for every $g\in G$, a nonnegative rational tolerance $\varepsilon$, and an integer $k\ge1$ in binary. Deciding whether $M_\varepsilon\le k$ is complete, under polynomial-time many-one reductions, for the existential theory of the reals. Hardness already holds for the trivial group, one query, and zero error. In particular the rational orthogonal polynomial-interface stationary problem is $\exists\R$-hard even with constant readouts and trivial dynamics.
\end{theorem}
\begin{proof}
Let $n=|G|$. If $k\ge n|\mathcal S|$, a deterministic seed--group register realizes every target exactly, so the answer is yes. Otherwise $k$ is bounded by a polynomial in the explicit input length, even though its encoding is binary.

For membership, Theorem~\ref{thm:equivariant56} permits permutation matrices $P_g$ forming a representation of $G$. Introduce real variables for their entries, constrain each entry $x$ by $x(x-1)=0$, impose every row and column sum equal to one, and impose
\[
                         P_1=I,
                 \qquad P_gP_h=P_{gh}\quad(g,h\in G).
\]
Introduce simplex variables $\alpha_s$ and categorical decoder rows $d_j(i)$. For every $s,j,g$ and output coordinate $c$, use a nonnegative auxiliary variable $z_{s,j,g,c}$ with
\[
 z_{s,j,g,c}\ge\pm\bigl((\alpha_sP_{g^{-1}}d_j)_c-F_{s,j}(g)_c\bigr),
 \qquad\sum_c z_{s,j,g,c}\le2\varepsilon.
\]
This is an existential real formula with rational coefficients, degree at most three, and polynomial size. Padding a representation by fixed unused labels handles an optimum strictly smaller than $k$. Theorem~\ref{thm:equivariant56} proves equivalence with the original unrestricted stationary problem. There is no group-closure computation in this explicit finite-group input model.

For hardness, let $A$ be a rational nonnegative matrix and let $k$ be a proposed nonnegative rank. Delete zero rows; the all-zero case is decided directly. Normalize each remaining row by its positive row sum to obtain a rational row-stochastic matrix $F$. This transformation is polynomial-time and preserves nonnegative rank. Interpret the rows as seeds and the columns as the categories of one query on the trivial group. A zero-error realization gives a stochastic factorization $F=\alpha d$ of inner dimension at most $k$, so the nonnegative rank is at most $k$.

Conversely, from any nonnegative factorization $F=UV$ delete a zero row of $V$ together with its column of $U$, and let $s_i=\sum_cV_{ic}>0$. Put $d_{ic}=V_{ic}/s_i$ and $\alpha_{bi}=U_{bi}s_i$. Then $d$ is row-stochastic and $\sum_i\alpha_{bi}=\sum_cF_{bc}=1$, so $\alpha$ is row-stochastic. Identity command rows yield an exact stationary realization. Thus $M_0$ equals the real nonnegative rank of $F$. Shitov's universality theorem, specifically its complexity consequence \cite[Theorem~2]{Shitov}, makes this decision problem $\exists\R$-complete. The last assertion follows by using the one-dimensional identity matrix as the only physical command and constant rational polynomial readouts.
\end{proof}
The upper bound is for an explicitly tabulated finite group. The hardness transfers to rational matrix input, but a polynomial-size existential description of an arbitrary compact matrix closure has not been supplied. Theorem~\ref{thm:effectivemin55} remains the separate termination result for that broader input class. The complexity obstruction above already occurs in randomized initialization and decoding; it is not caused solely by closure algorithms.

\subsection{A finite-table resource counterpart}
For a categorical interface, a dyadic table of precision $b$ has each probability in $2^{-b}\Z$. Besides peak label width $k$, charge its binary table length, all retained working registers, and the number of fair random bits used. A horizon-$N$ clocked implementation can store initialization, $N$ command tables, and terminal decoders. Its probability-table length is
\[
 O\left(b\left(|\mathcal S|k+N|\mathcal A|k^2+
                                k\sum_jM_j\right)\right).
\]
A uniform implementation is one finite algorithm that produces or accesses these tables from the encoded input, horizon, and accuracy. This requirement is additional to the original nonuniform model. An arbitrary real table, without an effective encoding, is not such an algorithm.

\begin{lemma}[Dyadic rounding of a probability row]\label{lem:dyadic56}
A probability row $p$ with $m$ entries has a dyadic row $\widehat p$ of precision $b$ such that
\[
                   \TV(p,\widehat p)\le(m-1)2^{-b}.
\]
One exact draw from $\widehat p$ uses at most $b$ independent fair bits.
\end{lemma}
\begin{proof}
Round the first $m-1$ entries down to multiples of $2^{-b}$ and put the remaining mass in the last entry. This is a probability row. Its total variation error equals the mass transferred to the last entry, which is at most $(m-1)2^{-b}$. For sampling, draw a uniform integer in $\{0,\ldots,2^b-1\}$ using $b$ bits and compare it with the integer cumulative row sums. Degenerate rows can use fewer bits.
\end{proof}

\begin{theorem}[Precision and random-bit budgets]\label{thm:bits56}
A categorical horizon-$N$ realization of peak width $k$ and error at most $\varepsilon$ admits a dyadic realization on the same registers whose error is at most
\begin{equation}\label{eq:clockround56}
       \varepsilon+\bigl((N+1)(k-1)+M_*-1\bigr)2^{-b},
                         \qquad M_* =\max_j M_j.
\end{equation}
Sampling initialization, all transitions, and one terminal answer requires at most $(N+2)b$ fair bits.

For a stationary compact-group realization of width $k$, the corresponding error bound can instead be made independent of $N$:
\begin{equation}\label{eq:stationaryround56}
                    \varepsilon+(k+M_*-2)2^{-b}.
\end{equation}
The stationary dyadic realization has at most $k$ labels, permutation command tables, table length
\[
 O\left(|\mathcal A|k\log(k+1)+
                   b\left(|\mathcal S|k+k\sum_jM_j\right)\right),
\]
and uses at most $2b$ fair bits in total for one experiment of any length. Given exact algebraic encodings of the selected initialization and decoder rows, the rounding and sampling are effective finite procedures.
\end{theorem}
\begin{proof}
Apply Lemma~\ref{lem:dyadic56} to every row. Couple the initial draws maximally and then, as long as the two labels agree, maximally couple their next transition draws. The probability of any disagreement is at most the sum of the initialization and $N$ row errors, hence at most $(N+1)(k-1)2^{-b}$. Conditioned on matching terminal labels, the terminal categorical rows have total variation error at most $(M_*-1)2^{-b}$. The coupling inequality gives \eqref{eq:clockround56}, uniformly over words and queries. The sample count is the number of draws times $b$.

For the stationary claim, first use Theorem~\ref{thm:purification55}. Its transition matrices are permutations and are already dyadic at every precision. Round only initialization and terminal categorical rows. The initialization discrepancy cannot increase under a permutation, so the same coupling gives \eqref{eq:stationaryround56} with no factor $N$. There are no random transition draws, leaving only initialization and one terminal answer. A permutation is stored as $k$ integer images, giving the displayed description bound. Exact comparison of algebraic numbers with dyadic rationals computes the necessary floors and handles equality, after which integer cumulative sampling is immediate.
\end{proof}
For a prescribed extra error $\eta>0$, choosing $b$ with the added term in \eqref{eq:clockround56} or \eqref{eq:stationaryround56} at most $\eta$ gives an explicit precision budget. The persistent label uses $\lceil\log_2k\rceil$ bits; an implementation of the sampler also retains a $b$-bit random integer and bounded indexing registers while sampling. These temporary registers are not disguised as free persistent randomness. The description and random-bit bounds do not assert a polynomial time bound for constructing an algebraic optimum. They also do not turn an exact irrational realization into an exact finite-bit one: the extra tolerance is essential in that case.

\section{Positive finite-precision matrix decoders}\label{sec:matrixprecision57}
The label-width bounds leave real table descriptions uncharged. There is
nevertheless a finite-description counterpart for each supplied encoded
realization. Entrywise rounding of a density matrix is unsafe: it need not
preserve positivity. A Gram-factor construction avoids this problem.

\begin{lemma}[Rational positive rounding]\label{lem:positiveprecision57}
Let $D\in\mathcal D_d$ and let $L=D^{1/2}$. Approximate the real and imaginary
parts of every entry of $L$ by dyadic multiples of $2^{-b}$, each within
$2^{-b}$, to obtain $Q$. Set $\delta=\sqrt2d\,2^{-b}$. If $\delta\le1/4$, then
\begin{equation}\label{eq:gramround57}
 \widehat D=\frac{QQ^*}{\tr(QQ^*)}\in\mathcal D_d\cap
            M_d(\Q+i\Q),\qquad
 \norm{\widehat D-D}_{\mathrm F}\le8\sqrt2d\,2^{-b}.
\end{equation}
The construction is effective when $D$ is given by exact algebraic entries.
Its rational output has $O(d^2(b+\log(d+1)))$ bits for fixed encoding
conventions.
\end{lemma}
\begin{proof}
Write $Q=L+E$. Then $\norm E_{\mathrm F}\le\delta$ and
$\norm L_{\mathrm F}^2=\tr D=1$. Hence $Q\ne0$,
$\tr(QQ^*)\ge(1-\delta)^2$, and
\[
 \norm{QQ^*-D}_{\mathrm F}\le(2+\delta)\delta,
 \qquad |\tr(QQ^*)-1|\le(2+\delta)\delta.
\]
The first estimate follows from the two cross terms and $EE^*$, and the
second from $\norm Q_{\mathrm F}^2-\norm L_{\mathrm F}^2$. Since
$\norm D_{\mathrm F}\le1$, division by $\tr(QQ^*)$ gives an error at most
$2(2+\delta)\delta/(1-\delta)^2\le8\delta$. Positivity and trace one are
exact, not numerical tolerances. After multiplying $Q$ by $2^b$, the
normalization is a ratio of integer Gram entries to a positive integer
sum of squares. Their bit sizes are $O(b+\log(d+1))$. The positive square
root of an algebraic positive matrix has algebraic entries and can be
specified and approximated by real-algebraic arithmetic. This proves the
effectivity assertion without an efficiency claim for obtaining $D$.
\end{proof}

\begin{theorem}[A supplied realization with finite tables and fair bits]
\label{thm:matrixbits57}
Suppose an error-$\varepsilon$ horizon-$N$ realization for a matrix
interface uses at most $k$ labels, and its initialization, command rows,
and decoders are supplied with exact algebraic encodings. Let
$\sqrt2d\,2^{-b}\le1/4$. There is a rational finite-table implementation,
with the same persistent label widths, whose Frobenius error is at most
\begin{equation}\label{eq:matrixbits57}
 \varepsilon+\sqrt2\bigl((N+1)(k-1)+8d\bigr)2^{-b}.
\end{equation}
Each hidden draw uses $b$ independent fair bits, so at most $(N+1)b$ such
bits suffice. For a supplied stationary permutation realization the
additional error is at most
\begin{equation}\label{eq:stationarymatrixbits57}
                  \sqrt2\bigl(k-1+8d\bigr)2^{-b},
\end{equation}
and only $b$ fair bits are needed for initialization, independently of the
horizon. The numerical matrix decoder itself is not a quantum preparation
or a sampled terminal answer register.
\end{theorem}
\begin{proof}
Apply Lemma~\ref{lem:dyadic56} to each initialization and command row after
padding to $k$ labels. Each row changes by at most $(k-1)2^{-b}$ in total
variation. Induction, or a maximal coupling at every actual draw, bounds
the terminal distribution discrepancy by $(N+1)(k-1)2^{-b}$; this upper
bound may harmlessly exceed one. Replace each decoder by the legal
rational matrix in Lemma~\ref{lem:positiveprecision57}. The Frobenius
diameter of $\mathcal D_d$ is $\sqrt2$, so changing the distribution incurs
at most $\sqrt2$ times its total-variation discrepancy. The independent
decoder error is at most $8\sqrt2d\,2^{-b}$. This proves
\eqref{eq:matrixbits57}. A dyadic probability row of denominator $2^b$ is
sampled by one uniform $b$-bit integer and its cumulative integer table.

For a permutation realization the command rows already have exact
integer entries and need not be rounded or sampled. Only the
initialization distribution and decoder are changed. The same estimate
therefore gives \eqref{eq:stationarymatrixbits57} and the horizon-independent
fair-bit bound.
\end{proof}
The row tables and their cumulative sums require
$O((N|\mathcal A|k^2+|\mathcal S|k)b)$ bits in the general case, in addition
to $O(kd^2(b+\log(d+1)))$ bits for one matrix decoder. Temporary random
integers and table indices require their own $O(b+\log(k+1))$ workspace;
this workspace is not hidden in the persistent-label invariant. Multiple
queries multiply the decoder storage by their number. To implement a
sampled measurement one must additionally specify and charge its sampler.
The theorem compiles a supplied encoded realization; it neither computes
a width minimizer nor assigns finite encodings to arbitrary unspecified
real tables. Applying it after the exact inner-polytope construction uses
an explicit extra error budget, not an exact finite-precision claim at an
extreme zero-error endpoint.

\section{Relation to numerical automata and positive realization}\label{sec:comparison54}
The input/output object and the order of quantifiers are decisive in comparing realization theorems. Classical probabilistic automata assign numerical acceptance probabilities to words, while stochastic-language results also study threshold recognition \cite{Rabin,Paz}. A numerical equivalence of two fixed machines is not the assertion that for every horizon there exists a new bounded-width approximate machine. Conversely, language equivalence need not preserve the numerical probabilities.

For weighted automata over a field, the finite-Hankel-rank characterization gives finite linear representations and identifies their minimum dimension \cite{BR}. It does not enforce nonnegative rows, stochastic normalization, or simplex-valued decoders. Our common-row contraction concerns these constraints simultaneously. The invariant observable quotient in Theorem~\ref{thm:observable53} uses the same reachable/invisible linear reduction; its additional conclusion is a finite-group criterion for exact positive realization with the full horizon-dependent quantifier.

The classical positive-realization problem asks for a positive system representing specified input/output data, and distinguishes positive dimension from unrestricted linear dimension \cite{BF04}. Invariant cones and polyhedral sections are part of that theory, not new notions introduced here. Proposition~\ref{prop:polygon54} is an explicit common invariant-enclosure construction. The new matched statement pairs it with a full-subcritical lower estimate against all clocked stochastic machines. Neither one cutwise nonnegative factorization nor a finite-rank Hankel matrix supplies that lower estimate.

Brodsky and Pippenger identify bounded-error measure-once quantum languages as group languages \cite[Theorem 3.3]{BP}. Their probabilistic simulation result preserves a cutpoint language with a potentially changed cutpoint \cite[Theorem 3.6]{BP}; it is not uniform additive preservation of every acceptance probability. Our categorical unitary corollary fixes an entire terminal measurement law in total variation and allows horizon-dependent stochastic approximants. The distinction is already visible in the $d$-outcome threshold $1-1/d$.

The closure theorem and strict-threshold decision method of Blondel, Jeandel, Koiran, and Portier \cite{BJKP}, and the effective Zariski-closure algorithms \cite{DJK,NPSW}, are algorithmic premises of Section~\ref{sec:effective54}. Computing the closure is not a new claim of this article. The additional outputs are a constrained minimax noise threshold, an attaining categorical decoder, and a terminating construction of an all-machine occupation constant. Recent Lie-algebra-based closure computations \cite{deGraaf} give an alternative representation of the algebraic closure, but their complex connected components still require real-component analysis for the present application.

Peter--Weyl approximation, Haar averaging, and the representative-function algebra are classical compact-group and almost-periodic-function machinery \cite{Folland}. Uniform approximation by a finite sum of matrix coefficients gives a finite \emph{linear} representation of an output function. It does not give positive stochastic rows or the repeated information loss quantified by \eqref{eq:occupation54}. Similarly, Steinberg's minimal compact topological automaton and enveloping-monoid results \cite[Theorems 2.2 and 2.6]{Steinberg} concern recognition of a language. His compact topological-monoid criterion \cite[Proposition 2.8]{Steinberg} uses clopen recognition, not a continuous simplex-valued response with a prescribed error norm. The finite-quotient radius formula is specific to groups. The stationarization theorem itself now holds for irreversible topological monoids with eventual approximate returns (Theorem~\ref{thm:stationarization56}).

Distribution-based bisimulation metrics provide robust comparisons for specified probabilistic automata \cite{FSZ}. They do not by themselves impose a minimum over all horizon-specific stochastic realizations of an external group experiment. The finite future-output equivalence in Theorem~\ref{thm:quotient54} is deliberately a component-quotient statement, rather than a claim that approximate equivalence always has a unique minimal positive stochastic realization.

\begin{center}\small
\begin{tabular}{p{.27\textwidth}p{.31\textwidth}p{.32\textwidth}}
\toprule
Comparison object & Error or equivalence & Machine quantifier\\
\midrule
Compact stochastic groups \cite{Flor} & Idempotent and permutation structure; no external error specification & One matrix group or semigroup kernel\\
Advice regular languages \cite{FP} & Neutral-letter recognition & One-scan advice automata; Boolean output\\
Weighted series \cite{BR} & Exact scalar coefficients; linear rank & One linear representation\\
Group-language simulation \cite{BP} & Cutpoint language; not additive probability error & One automaton\\
Positive realization \cite{BF04} & Exact specified numerical response & One positive system\\
Topological recognition \cite{Steinberg} & Clopen language acceptance & One recognizing action\\
Bisimulation metrics \cite{FSZ} & Distributional comparison of given systems & Two specified automata\\
Theorem~\ref{thm:stationarization56} & Same closed numerical error and label bound & Arbitrarily many narrow cuts; one extracted stationary machine\\
Theorem~\ref{thm:radius54} & Worst-word norm error; TV for a whole categorical law & A new clocked stochastic machine at every horizon\\
Theorem~\ref{thm:matched54} & Binary TV, every fixed $\varepsilon<\rho/2$ & Same nonuniform lower model; exact upper construction\\
\bottomrule
\end{tabular}
\end{center}
The contribution combines stationary purification and its exact finite-quotient criterion with a sharp constrained clocked radius, phasewise arbitrary-width occupation, and, for the specified Diophantine family, a matching full-subcritical growth law. The elementary radius center, the arithmetic pigeonhole argument, the closure algorithm, and the harmonic-analysis ingredients are not asserted individually new. The accompanying theorem-level audit records the closest comparisons and their changed hypotheses; it is not a claim of exhaustive priority clearance.


\subsection{The classical matrix-semigroup boundary}
Flor's description of nonnegative matrix groups \cite[Theorems 2 and 3]{Flor} identifies their idempotent geometry and the finiteness of bounded groups; his account explicitly credits Brown and Schwarz for earlier compact and stochastic cases. We do not claim this matrix structure as new. Theorem~\ref{thm:purification55} uses a minimal-rank idempotent in a joint semigroup, not just in the matrix projection: its physical coordinate must be the group identity. Sandwiching each command then preserves all uniform error inequalities by closure, even when the original stochastic rows violate every nontrivial physical group relation. The resulting initialization may be random, a point essential to exact state minimization.

The stationary theorem has no return hypothesis. The clocked theorem has a different quantifier and needs synchronization; the empty-return example and the phase-dependent example show why the distinction cannot be removed by terminology. The exact-return-length reduction is an additive-semigroup argument, whereas the profinite zero-error converse uses finite Hankel rank and continuous finite-dimensional representations. The new stationarization theorem removes the reversibility requirement for the equality of limiting width and stationary width, but does not supply a finite-quotient radius formula for an arbitrary irreversible monoid.

The finite-quotient existence theorem does not say that $[H:U]$ is the optimal number of stochastic labels. The actual minimum is the permutation-feasibility invariant in Theorem~\ref{thm:stationaryminimum55}; the finite cyclic example separates them. For infinitely many components, the infimum of finite-quotient radii can fail to be attained, even at zero. Finally, the exactly matched quantitative law retains its specified Diophantine alphabet. Theorem~\ref{thm:spherical56} adds a dimension-dependent rate, matched up to a logarithm, for spherical representations with a spectral gap; Corollary~\ref{cor:SO356} gives a concrete rational noncommuting instance. The exact unit-amplitude endpoint is not covered by these upper constructions. These distinctions delimit the proved statements without weakening any of the retained conclusions.

\subsection{What the reductions add}
The algebraic ingredients of Lemmas~\ref{lem:idempotent55} and \ref{lem:simplex55} are the classical minimum-rank kernel mechanism, the group corner of a compact matrix semigroup, and the recurrent simplex of a stochastic idempotent. Flor's Theorem~2 describes nonnegative idempotents; his Theorem~3 and preceding Lemma~3 give the permutation structure of bounded nonnegative groups \cite{Flor}. The paper explicitly attributes earlier cases to Brown and Schwarz. The present proof adds the physical coordinate to the closure, chooses the idempotent over its identity, and tracks the closed error balls through sandwiching. Theorem~\ref{thm:equivariant56} further averages the finite kernel above the physical identity and descends to an actual continuous $H$-action. These are specified realization reductions using classical structure, not a new classification of matrix-semigroup kernels.

For deterministic Boolean recognition, neutral-letter removal of advice is an established phenomenon. Fijalkow--Paperman \cite[Theorem~10]{FP} prove the Crane Beach property for advice regular languages by padding distinguishing words and suffixes, and their account relates one-scan programs to earlier Ramsey arguments. In that setting a bounded number of deterministic states bounds the relevant distinguishability classes. A stochastic $k$-label register has a continuum of state distributions, so this discrete class count does not apply. Theorem~\ref{thm:stationarization56} instead colors approximate tuples of stochastic composite kernels, extracts one common tuple, and absorbs a final filler matrix into the decoder. It preserves the same numerical tolerance at equality and the same $k$, and bounds all narrow-cut occurrences with unrestricted intermediate widths. Eventual approximate returns further replace a literal neutral letter. The finite Ramsey theorem itself is classical \cite{Ramsey}.

The spherical lower estimate uses a classical spectral gap, cap measure, and smoothing; its upper estimate uses a polytope enclosing a contracted sphere. The contribution is their realization-theoretic combination with the all-profile conditional-centroid bound. The rational $\mathrm{SO}(3)$ example invokes \cite{BG2012} as a deep external input. Similarly, the lower-complexity input in Theorem~\ref{thm:ETR56} is Shitov's nonnegative-factorization universality theorem \cite{Shitov}; the stochastic normalization and finite-group existential upper reduction are given explicitly. Neither external theorem is claimed as part of our proof technology's novelty.

\subsection{Companion results}
The companion supplement gives additional results on compatible reachable sections, one-surplus geometry, stable sections, finite-horizon certificates, word distortion, two-state duality, finite-bit compilation, and global residual hierarchies. The scalar two-extreme argument and earlier quantitative constants are also retained there. The present main proofs are self-contained apart from the classical results explicitly cited; none of these companion derivations is used to conceal a missing step in a principal theorem.

\subsection{Orbit geometry, positive matrix outputs, and recognition}
The homogeneous-space covering estimates and the smooth orbit geometry used in Section~\ref{sec:uniform57} are classical; see, for example, Szarek~\cite{Szarek57}. The spectral-gap input is Bourgain--Gamburd~\cite{BG2012}. Dense algebraic free generators in Section~\ref{sec:extreme57} come from Breuillard--Gelander~\cite[Theorem~1.1]{BreuillardGelander57}; the later repair and supplementary proof details are recorded in~\cite{BreuillardGelanderCorrection57}. None of these results is claimed as a new group-theoretic theorem here. The realization argument requires, in addition, a contraction uniform over every conditional-centroid direction, one common positive row for each command and label, and a terminal decoder in the original legal matrix set.

Quantum finite automata already exhibit strong state-complexity separations in language recognition; Ambainis--Freivalds~\cite{AF57} is a relevant early comparison. Theorem~\ref{thm:freeendpoint57} instead asks a classical stochastic machine to reproduce one entire legal density-matrix vector at every input word, permits horizon-specific redesign and a free clock, and compares exact and fixed-error width for the same interface. Its exact lower bound uses extreme points of the decoder set; it does not apply unchanged to a single acceptance probability or to an enlarged unconstrained output simplex. A finite tomography map preserves width only when its legal image and reconstructed norm are retained. This comparison identifies different quantifiers and resources, not an exhaustive priority claim about all quantum or probabilistic automata formalisms.

The logarithm in the nonabelian lower bound remains explicit. Uniform orbit small balls and a spectral gap imply the displayed exponent with that loss, not a universally sharp width law for every compact representation. General online outputs, adaptive feedback, and unrestricted compact semigroup radius formulas require additional arguments. The same-width stationarization result for topological monoids remains available under its actual approximate-return hypothesis; it does not provide that missing radius formula.

\begin{thebibliography}{99}
\bibitem{BPR} S. Basu, R. Pollack, and M.-F. Roy, \emph{Algorithms in Real Algebraic Geometry}, second ed., Algorithms and Computation in Mathematics 10, Springer, Berlin, 2006.
\bibitem{BF04} L. Benvenuti and L. Farina, A tutorial on the positive realization problem, \emph{IEEE Transactions on Automatic Control} \textbf{49} (2004), no. 5, 651--664.
\bibitem{BR} J. Berstel and C. Reutenauer, \emph{Noncommutative Rational Series with Applications}, Encyclopedia of Mathematics and its Applications 137, Cambridge University Press, Cambridge, 2011.
\bibitem{BJKP} V. D. Blondel, E. Jeandel, P. Koiran, and N. Portier, Decidable and undecidable problems about quantum automata, \emph{SIAM Journal on Computing} \textbf{34} (2005), no. 6, 1464--1473.
\bibitem{BP} A. Brodsky and N. Pippenger, Characterizations of 1-way quantum finite automata, \emph{SIAM Journal on Computing} \textbf{31} (2002), no. 5, 1456--1478.
\bibitem{DJK} H. Derksen, E. Jeandel, and P. Koiran, Quantum automata and algebraic groups, \emph{Journal of Symbolic Computation} \textbf{39} (2005), nos. 3--4, 357--371.
\bibitem{FSZ} Y. Feng, L. Song, and L. Zhang, Distribution-based bisimulation and bisimulation metric in probabilistic automata, arXiv:1512.05027 (2015).
\bibitem{Flor} P. Flor, On groups of non-negative matrices, \emph{Compositio Mathematica} \textbf{21} (1969), no. 4, 376--382.
\bibitem{Folland} G. B. Folland, \emph{A Course in Abstract Harmonic Analysis}, second ed., Chapman \& Hall/CRC, Boca Raton, 2016.
\bibitem{deGraaf} W. A. de Graaf, Computing the Zariski closure of a finitely generated matrix group, arXiv:2505.02577v2 (2026).
\bibitem{Hall} B. C. Hall, \emph{Lie Groups, Lie Algebras, and Representations: An Elementary Introduction}, second ed., Graduate Texts in Mathematics 222, Springer, Cham, 2015.
\bibitem{NPSW} K. Nosan, A. Pouly, S. Schmitz, M. Shirmohammadi, and J. Worrell, On the computation of the Zariski closure of finitely generated groups of matrices, arXiv:2106.01853v6 (2025).
\bibitem{Paz} A. Paz, \emph{Introduction to Probabilistic Automata}, Academic Press, New York, 1971.
\bibitem{Rabin} M. O. Rabin, Probabilistic automata, \emph{Information and Control} \textbf{6} (1963), no. 3, 230--245.
\bibitem{Steinberg} B. Steinberg, Topological dynamics and recognition of languages, arXiv:1306.1468 (2013).
\bibitem{BG2012} J. Bourgain and A. Gamburd, A spectral gap theorem in $SU(d)$, \emph{Journal of the European Mathematical Society} \textbf{14} (2012), no. 5, 1455--1511.
\bibitem{FP} N. Fijalkow and C. Paperman, Monadic second-order logic with arbitrary monadic predicates, arXiv:1709.03117 (2017).
\bibitem{Ramsey} F. P. Ramsey, On a problem of formal logic, \emph{Proceedings of the London Mathematical Society} (2) \textbf{30} (1930), no. 1, 264--286.
\bibitem{Shitov} Y. Shitov, A universality theorem for nonnegative matrix factorizations, arXiv:1606.09068v2 (2018).
\bibitem{BreuillardGelander57} E. Breuillard and T. Gelander, On dense free subgroups of Lie groups, \emph{Journal of Algebra} \textbf{261} (2003), no. 2, 448--467.
\bibitem{BreuillardGelanderCorrection57} E. Breuillard and T. Gelander, On dense free subgroups of Lie groups---revisited, arXiv:2605.20568 (2026).
\bibitem{Szarek57} S. J. Szarek, Metric entropy of homogeneous spaces, \emph{Banach Center Publications} \textbf{43} (1998), 395--410.
\bibitem{AF57} A. Ambainis and R. Freivalds, 1-way quantum finite automata: strengths, weaknesses and generalizations, arXiv:quant-ph/9802062 (1998).
\end{thebibliography}

\end{document}
